\documentclass{report}
\usepackage{graphicx} % Required for inserting images
\usepackage{amsmath} % AMS Math Package
\usepackage{amsthm} % Theorem Formatting
\usepackage{amssymb} % Math symbols such as \mathbb
\usepackage[hidelinks]{hyperref}
\usepackage{graphicx}
\usepackage{caption}
\usepackage{filecontents}
\usepackage{subcaption}% Allows for eps images
\usepackage[table,xcdraw]{xcolor}
\usepackage{tikz}
\usepackage{siunitx}
\usepackage{tabularx} % Allows for table with custom column width
\usepackage{lipsum}
\usepackage{parskip}
\usepackage{booktabs}
\usepackage{mathtools}
 
% ── Theorem spacing ──────────────────────────────────────────────
% parskip sets \parskip ≈ 0.5\baselineskip, but amsthm's default
% pre/post skip is only \topsep (≈ 4--6pt), so consecutive
% environments feel jammed.  We override all three styles with
% generous 14pt gaps (roughly one blank-line worth of breathing room).
 
\newtheoremstyle{spacedplain}%      name
  {14pt}%                           space above
  {14pt}%                           space below
  {\itshape}%                       body font
  {}%                               indent
  {\bfseries}%                      head font
  {.}%                              punctuation after head
  { }%                              space after head  (a normal space)
  {}%                               head spec (empty = default)
 
\newtheoremstyle{spaceddefinition}%
  {14pt}{14pt}%
  {\normalfont}{}{%
  \bfseries}{.}{ }{}
 
\newtheoremstyle{spacedremark}%
  {14pt}{14pt}%
  {\normalfont}{}{%
  \itshape}{.}{ }{}
 
% Apply the spaced styles
\theoremstyle{spacedplain}
\newtheorem{theorem}{Theorem}[section]
\newtheorem{proposition}[theorem]{Proposition}
\newtheorem{lemma}[theorem]{Lemma}
\newtheorem{corollary}[theorem]{Corollary}
 
\theoremstyle{spaceddefinition}
\newtheorem{definition}[theorem]{Definition}
\newtheorem{example}[theorem]{Example}
 
\theoremstyle{spacedremark}
\newtheorem{remark}[theorem]{Remark}

\title{Quantitative Finance: Mathematical Foundations and Models}
\author{Liam O'Shaughnessy}
\date{March 2026}

\begin{document}

\maketitle

\chapter*{Preface}
These notes were written to consolidate the mathematical foundations of modern quantitative finance, with an emphasis on arbitrage pricing, stochastic calculus, and practical modeling considerations. They are not intended as a textbook, but as a structured reference bridging theory and implementation. This is more theoretically and academically focused, which I wanted to make sure I had a thorough understanding of - I would like to continue developing practical approaches to effectuate these methods.

\tableofcontents

\pagenumbering{arabic}

\chapter{Introduction}
This has been my attempt to touch on most topics relevent to quantitative finance theory. I start with the mathematics of modeling assets as stochastic processes, which culminates with the Kolmogorov equations/generator framework for Markov processes that gets used for pricing. I then formalize risk-neutral pricing and how the Girsanov's theorem establishes the criteria to change to a martingale measure. In the next chapter, I price flavors of options, and discuss models for fitting implied volatility. Following that, I model rates stochastically, and introduce credit and counterparty risk for complete pricing. I then address portfolio management and risk assessment, as well as market microstructure dynamics. I round things up with time series and an overview of numerical methods for solving the proposed equations.

These notes are drawn from the following books and resources: \textit{Stochastic Calculus for Finance II} by Shreve, \textit{Arbitrage in Continuous Time} by Björk, \textit{Probability Theory} by Sinai and Koralov, \textit{Applied Stochastic Analysis} by Holmes-Cerfon, \textit{Brownian Motion} by Schilling and Partzsch, \textit{Kolmogorov Equations} by Ludvigsson, \textit{All of Statistics} by Wasserman, \textit{Convex Optimization} by Boyd and Vandenberghe, \textit{Monte Carlo Methods in Financial Engineering} by Glasserman, 1\textit{Advances in Financial Machine Learning} by de Prado, 8.S096 on MIT OCW, ECO362 from Princeton, Time Series Analysis from Baruch's MFE program, \textit{Credit Risk Modeling} by Hurd from McMaster, and, of course, LLMs.



\chapter{Mathematical Foundations}

The first half of ``quantitative finance'' is \emph{quantitative}. Every pricing formula, hedging argument, and risk model rests on a mathematical engine built from three layers: a rigorous probability framework that can handle continuous randomness, a theory of stochastic processes that models how randomness evolves through time, and a stochastic calculus that lets us differentiate and integrate processes driven by that randomness. This chapter constructs those three layers in sequence, always keeping sight of \emph{why} a financial modeler needs each piece.

The guiding thread is a single question: \emph{How do we write down, manipulate, and solve equations for asset prices that move randomly?} By the end of this chapter, we will have the tools to formulate stochastic differential equations for asset dynamics, connect them to partial differential equations that can be solved numerically, and change the probability measure under which we evaluate expectations, which will be the central trick of derivative pricing.

\section{Probability Theory}

To model a stock price, we need to formalize what it means for a quantity to be ``random.'' Intuition alone is insufficient: continuous-time finance requires integration over uncountably many outcomes, conditional expectations that serve as ``best forecasts,'' and the ability to switch between different probability measures. The Kolmogorov axioms provide the foundation.

\subsection{Axioms of Probability}

\begin{definition}
Let $\Omega$ be a non-empty set, the \textbf{sample space} of all possible outcomes. A collection $\mathcal{F}$ of subsets of $\Omega$ is called a \textbf{$\sigma$-algebra} if:
\begin{enumerate}
    \item $\emptyset \in \mathcal{F}$,
    \item $A \in \mathcal{F} \implies A^c \in \mathcal{F}$,
    \item $A_1, A_2, \ldots \in \mathcal{F} \implies \bigcup_{n=1}^\infty A_n \in \mathcal{F}$.
\end{enumerate}
The elements of $\mathcal{F}$ are called \textbf{measurable sets} (or \textbf{events}), and $(\Omega, \mathcal{F})$ is a \textbf{measurable space}.
\end{definition}

\begin{remark}
Why do we need a $\sigma$-algebra rather than just the power set of $\Omega$? In finite settings, we can indeed take $\mathcal{F} = 2^\Omega$. But when $\Omega$ is uncountable, as it will be when we model continuous price paths, there exist pathological subsets of $\mathbb{R}$ that cannot be assigned a probability in a consistent way (the Vitali construction). The $\sigma$-algebra restricts attention to ``well-behaved'' sets. In practice, this is almost never a binding constraint; it is the price of mathematical coherence.
\end{remark}

\begin{definition}
For $A$ a collection of subsets of $\Omega$, the \textbf{$\sigma$-algebra generated by $A$}, denoted $\sigma(A)$, is the intersection of all $\sigma$-algebras containing $A$, or equivalently, the smallest $\sigma$-algebra containing $A$. The \textbf{Borel $\sigma$-algebra} on $\mathbb{R}$, written $\mathcal{B}(\mathbb{R})$, is the $\sigma$-algebra generated by the open subsets of $\mathbb{R}$.
\end{definition}

\begin{definition}
Let $(\Omega, \mathcal{F})$ be a measurable space. A function $\mathbb{P}: \mathcal{F} \to [0,1]$ is a \textbf{probability measure} if $\mathbb{P}(\Omega) = 1$ and $\mathbb{P}\bigl(\bigcup_n A_n\bigr) = \sum_n \mathbb{P}(A_n)$ for any countable collection of disjoint sets $A_n \in \mathcal{F}$. The triple $(\Omega, \mathcal{F}, \mathbb{P})$ is a \textbf{probability space}. An event with $\mathbb{P}(A)=0$ has \textbf{measure zero}; an event with $\mathbb{P}(A)=1$ holds \textbf{almost surely (a.s.)}.
\end{definition}

\begin{definition}
A function $X: \Omega \to \mathbb{R}$ is \textbf{$\mathcal{F}$-measurable} if for every Borel set $B \in \mathcal{B}(\mathbb{R})$, the preimage $X^{-1}(B) \in \mathcal{F}$. An $\mathcal{F}$-measurable function on a probability space is called a \textbf{random variable}. The $\sigma$-algebra generated by $X$ is
\begin{align}
    \sigma(X) = \{X^{-1}(B) : B \in \mathcal{B}(\mathbb{R})\},
\end{align}
the smallest $\sigma$-algebra making $X$ measurable.
\end{definition}

\begin{remark}
Measurability is the probabilistic analogue of continuity in topology. Its practical importance: if $X$ is $\mathcal{F}$-measurable, then for any Borel set $B$, the probability $\mathbb{P}(X \in B)$ is well-defined. In finance, measurability with respect to a \emph{filtration} (defined later) encodes which information is available at which time - the mathematical formalization of ``no peeking at the future.''
\end{remark}

\begin{example}[Discrete probability space]
Consider tossing a fair coin twice. The sample space is $\Omega = \{HH, HT, TH, TT\}$. Take $\mathcal{F} = 2^\Omega$ and define the uniform measure $\mathbb{P}(\{\omega\}) = 1/4$ for each $\omega \in \Omega$. Then $(\Omega, \mathcal{F}, \mathbb{P})$ is a probability space. For instance,
\begin{align}
    \mathbb{P}(\{HH, HT\}) = \mathbb{P}(\{HH\}) + \mathbb{P}(\{HT\}) = \frac{1}{2}.
\end{align}
Define the random variable $X(\omega) = \text{number of heads}$. Then $\sigma(X) = \bigl\{\emptyset, \{TT\}, \{HT,TH\}, \{HH\}, \ldots\bigr\}$ partitions $\Omega$ by the value of $X$.
\end{example}

\begin{example}[Financial interpretation]
Consider a one-period market with two states: the economy is ``up'' ($\omega_u$) or ``down'' ($\omega_d$). A stock has payoff $S(\omega_u) = 110$, $S(\omega_d) = 90$. The random variable $S$ is $\mathcal{F}$-measurable for $\mathcal{F} = 2^{\{\omega_u,\omega_d\}}$. The $\sigma$-algebra $\sigma(S) = \mathcal{F}$ because $S$ distinguishes all states. If instead we had $S(\omega_u) = S(\omega_d) = 100$, then $\sigma(S) = \{\emptyset, \Omega\}$: the stock reveals nothing about which state occurred. This distinction becomes critical when we formalize what ``information'' a price process reveals.
\end{example}

\subsection{Distributions}

\begin{definition}
Let $X$ be a random variable on $(\Omega, \mathcal{F}, \mathbb{P})$. The \textbf{distribution measure} $\mu_X$ assigns to each Borel set $B \subseteq \mathbb{R}$ the mass
\begin{align}
    \mu_X(B) = \mathbb{P}(X \in B) = \mathbb{P}(\{\omega : X(\omega) \in B\}).
\end{align}
The \textbf{cumulative distribution function (CDF)} is $F(x) = \mathbb{P}(X \leq x)$.
\end{definition}

\begin{definition}
The distribution $\mu_X$ admits a \textbf{probability density function (pdf)} $f$ if
\begin{align}
    \mathbb{P}(a \leq X \leq b) = \int_a^b f(x)\, dx.
\end{align}
If $X$ takes values in a countable set $\{x_1, x_2, \ldots\}$, the \textbf{probability mass function (pmf)} is $p_i = \mathbb{P}(X = x_i)$.
\end{definition}

The following distributions appear repeatedly in quantitative finance. Rather than simply listing them, we emphasize \emph{where} each one arises.

\begin{example}[Gaussian]
A \textbf{Gaussian} random variable $X \sim \mathcal{N}(\mu, \sigma^2)$ has density
\begin{align}
    f(x) = \frac{1}{\sqrt{2\pi\sigma^2}}\exp\Bigl(-\frac{(x-\mu)^2}{2\sigma^2}\Bigr).
\end{align}
The Gaussian is the default distribution in finance because of the Central Limit Theorem and because Brownian motion, the engine of continuous-time models, has Gaussian increments. Log-returns over short intervals are approximately Gaussian under many models, and the Black-Scholes formula is ultimately an exercise in Gaussian integration.

Key properties: $\mathbb{E}[X] = \mu$, $\text{Var}(X) = \sigma^2$, and if $X_1 \sim \mathcal{N}(\mu_1, \sigma_1^2)$ and $X_2 \sim \mathcal{N}(\mu_2, \sigma_2^2)$ are independent, then $X_1 + X_2 \sim \mathcal{N}(\mu_1 + \mu_2, \sigma_1^2 + \sigma_2^2)$---stability under addition, which is why sums of independent Gaussians remain Gaussian.
\end{example}

\begin{example}[Exponential]
An \textbf{exponential} random variable $\tau \sim \text{Exp}(\lambda)$ has density
\begin{align}
    f(t) = \lambda e^{-\lambda t}, \quad t \geq 0.
\end{align}
This models waiting times: the time until the next default in a credit portfolio, or the time until the next jump in a jump-diffusion model. The defining property is \textbf{memorylessness}:
\begin{align}
    \mathbb{P}(\tau > t + s \mid \tau > s) = \frac{e^{-\lambda(t+s)}}{e^{-\lambda s}} = e^{-\lambda t} = \mathbb{P}(\tau > t).
\end{align}
Given that no event has occurred by time $s$, the residual waiting time has the same distribution as the original, so the process ``forgets'' how long it has been waiting. This is why the exponential distribution is the natural model for default intensity in reduced-form credit models: the hazard rate $\lambda$ summarizes all information about default timing.
\end{example}

\begin{example}[Poisson]
A \textbf{Poisson} random variable $N \sim \text{Pois}(\lambda)$ has pmf
\begin{align}
    \mathbb{P}(N = k) = \frac{\lambda^k}{k!}e^{-\lambda}, \quad k \in \mathbb{N}_0.
\end{align}
Setting $\lambda = r \cdot t$ counts the expected number of events (jumps, defaults, trades) in a window of length $t$ when the arrival rate is $r$. The Poisson process, constructed from these random variables, will drive the jump component of asset price models later in this chapter.
\end{example}

\begin{definition}
Two $\sigma$-algebras $\mathcal{G}, \mathcal{H} \subseteq \mathcal{F}$ are \textbf{independent} if $\mathbb{P}(A \cap B) = \mathbb{P}(A)\mathbb{P}(B)$ for all $A \in \mathcal{G}$, $B \in \mathcal{H}$. Random variables $X$ and $Y$ are independent if $\sigma(X)$ and $\sigma(Y)$ are independent.
\end{definition}

\begin{remark}
Independence is a modeling assumption, not a fact about the world. In finance, independence is routinely assumed for computational convenience (e.g., i.i.d.\ returns in the random walk model) and routinely violated in practice (e.g., volatility clustering, contagion). Recognizing where independence holds and where it breaks down is a recurring theme.
\end{remark}

\subsection{Lebesgue Integration and Expectation}

With random variables defined, we need a notion of ``average value.'' The Riemann integral is inadequate here: the domain of a random variable is $\Omega$, an abstract sample space with no natural partition into intervals. The Lebesgue integral solves this by partitioning the \emph{range} instead.

\begin{definition}
Let $X$ be a non-negative random variable on $(\Omega, \mathcal{F}, \mathbb{P})$ taking values in $[a,b]$. Partition $[a,b]$ into sub-intervals via points $y_0 = a < y_1 < \cdots < y_n = b$. For each sub-interval, define $A_k = \{\omega \in \Omega : y_{k-1} \leq X(\omega) \leq y_k\}$. As the mesh of the partition tends to zero, the sum
\begin{align}
    \sum_k y_k\, \mathbb{P}(A_k) \;\longrightarrow\; \int_\Omega X(\omega)\, d\mathbb{P}(\omega)
\end{align}
converges to the \textbf{Lebesgue integral}.
\end{definition}

\begin{remark}
The intuition: instead of asking ``how does $X$ vary as I sweep across $\Omega$?'' (which is hard, since $\Omega$ has no geometry), we ask ``how much of $\Omega$ does $X$ assign to each slice of the real line?'' and weight accordingly. This is why we need a probability \emph{measure}, which gives the size of each preimage set $A_k$.
\end{remark}

\begin{definition}
For a general (possibly negative) random variable, decompose $X = X^+ - X^-$ where $X^+ = \max(X,0)$ and $X^- = \max(-X,0)$. If both $\int_\Omega X^+ \, d\mathbb{P}$ and $\int_\Omega X^- \, d\mathbb{P}$ are finite, $X$ is called \textbf{integrable}, and
\begin{align}
    \int_\Omega X\, d\mathbb{P} = \int_\Omega X^+\, d\mathbb{P} - \int_\Omega X^-\, d\mathbb{P}.
\end{align}
\end{definition}

\begin{theorem}
$X$ is integrable if and only if $\int_\Omega |X|\, d\mathbb{P} < \infty$.
\end{theorem}

\begin{proof}
Since $|X| = X^+ + X^-$ and $X^+, X^- \leq |X|$, finiteness of $\int|X|$ implies finiteness of $\int X^\pm$ by the comparison property, so $X$ is integrable. Conversely, if both parts have finite integrals, linearity gives $\int|X| = \int X^+ + \int X^- < \infty$.
\end{proof}

\begin{definition}
The \textbf{expectation} of a random variable $X$ on $(\Omega, \mathcal{F}, \mathbb{P})$ is
\begin{align}
    \mathbb{E}[X] = \int_\Omega X(\omega)\, d\mathbb{P}(\omega).
\end{align}
\end{definition}

\begin{definition}
The \textbf{variance} and \textbf{covariance} of integrable random variables $X, Y$ are
\begin{align}
    \text{Var}(X) &= \mathbb{E}[(X - \mathbb{E}[X])^2], \\
    \text{Cov}(X,Y) &= \mathbb{E}[(X - \mathbb{E}[X])(Y - \mathbb{E}[Y])].
\end{align}
\end{definition}

\begin{theorem}[Change of variable / Law of the unconscious statistician]
Let $X$ be a random variable with distribution $\mu_X$ and $g$ a Borel-measurable function with $\int |g(x)|\, d\mu_X(x) < \infty$. Then
\begin{align}
    \mathbb{E}[g(X)] = \int_\Omega g(X(\omega))\, d\mathbb{P}(\omega) = \int_{\mathbb{R}} g(x)\, d\mu_X(x).
\end{align}
If $X$ has density $f$, then $\mathbb{E}[g(X)] = \int_{\mathbb{R}} g(x) f(x)\, dx$.
\end{theorem}

\begin{remark}[Standard machine]
The proof proceeds by verifying the result for indicator functions $g = \mathbb{I}_B$, extending by linearity to simple functions, then using monotone convergence to handle general non-negative measurable functions, and finally the general case by decomposition into positive and negative parts.
\end{remark}

\begin{example}[Computing a derivative's expected payoff]
A European call option on a stock $S$ with strike $K$ pays $\max(S-K, 0)$ at expiry. If the terminal stock price has density $f_S$, the expected payoff is
\begin{align}
    \mathbb{E}[\max(S-K,0)] = \int_K^\infty (s-K) f_S(s)\, ds.
\end{align}
Under the Black-Scholes model, $\log S$ is Gaussian, and this integral evaluates to the Black-Scholes formula. The entire derivation is ultimately an application of the change-of-variable theorem.
\end{example}

\subsection{Conditioning and Change of Measure}

Conditional expectation is the mathematical formalization of ``best forecast given partial information.'' Change of measure, which reweights probabilities, is the mechanism behind risk-neutral pricing. Both are indispensable.

\begin{definition}
Let $X$ be an integrable random variable on $(\Omega, \mathcal{F}, \mathbb{P})$ and $\mathcal{G} \subseteq \mathcal{F}$ a sub-$\sigma$-algebra. The \textbf{conditional expectation} $\mathbb{E}[X \mid \mathcal{G}]$ is any random variable that is:
\begin{enumerate}
    \item $\mathcal{G}$-measurable, and
    \item satisfies $\int_A \mathbb{E}[X \mid \mathcal{G}]\, d\mathbb{P} = \int_A X\, d\mathbb{P}$ for all $A \in \mathcal{G}$ (partial averaging).
\end{enumerate}
\end{definition}

\begin{remark}
The conditional expectation $\mathbb{E}[X \mid \mathcal{G}]$ is the best $L^2$-approximation of $X$ among $\mathcal{G}$-measurable random variables, or the ``projection'' of $X$ onto the subspace of random variables whose values are determined by $\mathcal{G}$. In financial terms: if $\mathcal{G}$ represents information available at time $s$, then $\mathbb{E}[X \mid \mathcal{G}]$ is the best estimate of the future quantity $X$ based on what is known at time $s$.
\end{remark}

\begin{proposition}
Let $\mathcal{H} \subseteq \mathcal{G} \subseteq \mathcal{F}$ be sub-$\sigma$-algebras and $X, Y$ integrable. Then:
\begin{enumerate}
    \item If $X$ is $\mathcal{G}$-measurable: $\mathbb{E}[XY \mid \mathcal{G}] = X \cdot \mathbb{E}[Y \mid \mathcal{G}]$ (taking out what is known).
    \item $\mathbb{E}[\mathbb{E}[X \mid \mathcal{G}] \mid \mathcal{H}] = \mathbb{E}[X \mid \mathcal{H}]$ (tower property).
    \item If $X$ is independent of $\mathcal{G}$: $\mathbb{E}[X \mid \mathcal{G}] = \mathbb{E}[X]$.
\end{enumerate}
\end{proposition}

\begin{remark}
The tower property is the formal statement of iterated expectations: the forecast of a forecast is the forecast based on the coarser information. This will be invoked constantly, as it is the primary tool for computing risk-neutral prices recursively.
\end{remark}

\begin{example}[Conditional expectation in a two-step model]
Consider two fair coin tosses with $\Omega = \{HH, HT, TH, TT\}$ and uniform measure. Let $X$ = total number of heads, and $\mathcal{G} = \sigma(\text{first toss})$, whose atoms are $\{HH, HT\}$ and $\{TH, TT\}$.

On $\{HH, HT\}$ (first toss is heads): $\mathbb{E}[X \mid \mathcal{G}] = (2+1)/2 = 3/2$.

On $\{TH, TT\}$ (first toss is tails): $\mathbb{E}[X \mid \mathcal{G}] = (1+0)/2 = 1/2$.

Verify partial averaging: $\mathbb{E}[\mathbb{E}[X \mid \mathcal{G}]] = \frac{1}{2}\cdot\frac{3}{2} + \frac{1}{2}\cdot\frac{1}{2} = 1 = \mathbb{E}[X]$. \checkmark
\end{example}

\begin{example}[Conditional expectation as a pricing operator]
In a one-period binomial model, a stock is worth $S_0 = 100$ today. Tomorrow, it goes to $S_1 = 120$ (probability $p$) or $S_1 = 80$ (probability $1-p$). The payoff of a call with strike $K=100$ is $X = \max(S_1 - 100, 0)$, which equals 20 in the up state and 0 in the down state.

The ``price'' of this call today is a conditional expectation: $\mathbb{E}[X \mid \mathcal{G}_0]$, where $\mathcal{G}_0 = \{\emptyset, \Omega\}$ represents no information beyond the current state. Under the risk-neutral measure (which we will construct formally later), $p$ is replaced by a risk-neutral probability $q = (e^r - d)/(u-d)$ that makes the discounted stock a martingale, and the price becomes $e^{-r}\mathbb{E}^{\widetilde{\mathbb{P}}}[X]$.
\end{example}

\bigskip

The ability to switch from one probability measure to another is the mathematical heart of derivative pricing. The following theorem provides the mechanism.

\begin{theorem}[Radon-Nikodym / Change of measure]
Let $(\Omega, \mathcal{F}, \mathbb{P})$ be a probability space. Let $Z$ be a non-negative random variable with $\mathbb{E}[Z] = 1$. Then
\begin{align}
    \widetilde{\mathbb{P}}(A) = \int_A Z(\omega)\, d\mathbb{P}(\omega), \quad A \in \mathcal{F},
\end{align}
defines a new probability measure $\widetilde{\mathbb{P}}$. The random variable $Z$ is called the \textbf{Radon-Nikodym derivative}, written $Z = d\widetilde{\mathbb{P}}/d\mathbb{P}$. Expectations transform as
\begin{align}
    \widetilde{\mathbb{E}}[X] = \mathbb{E}[XZ].
\end{align}
\end{theorem}

\begin{remark}
This theorem says that we can reweight outcomes: events that were unlikely under $\mathbb{P}$ can become likely under $\widetilde{\mathbb{P}}$, and vice versa, as long as the reweighting is normalized. In finance, this reweighting converts the ``real-world'' measure (where stocks earn a risk premium) into a ``risk-neutral'' measure (where discounted stocks are martingales). The Radon-Nikodym derivative $Z$ encodes the market price of risk.
\end{remark}

\begin{example}[Concrete change of measure]
Suppose $\Omega = \{\omega_u, \omega_d\}$ with $\mathbb{P}(\omega_u) = 0.6$, $\mathbb{P}(\omega_d) = 0.4$. Define $Z(\omega_u) = 0.5/0.6$ and $Z(\omega_d) = 0.5/0.4$. Check: $\mathbb{E}[Z] = 0.6 \cdot \frac{0.5}{0.6} + 0.4 \cdot \frac{0.5}{0.4} = 1$. Then $\widetilde{\mathbb{P}}(\omega_u) = 0.5$ and $\widetilde{\mathbb{P}}(\omega_d) = 0.5$---the measure change has turned an asymmetric world into a fair one. This is exactly what risk-neutral pricing does in a one-period binomial model.
\end{example}

\subsection{Inference}

The mathematical framework connects to data through limit theorems that justify using sample averages as estimators and that underpin Monte Carlo simulation.

\begin{definition}
A sequence $\{X_n\}$ of random variables is \textbf{independent and identically distributed (i.i.d.)} if all variables are mutually independent and share a common distribution.
\end{definition}

\begin{definition}
A sequence $\{X_n\}$ \textbf{converges in probability} to $X$ if $\lim_{n \to \infty} \mathbb{P}(|X_n - X| > \epsilon) = 0$ for all $\epsilon > 0$. It \textbf{converges in distribution} to $X$ if $\lim_{n \to \infty} F_{X_n}(x) = F_X(x)$ at all continuity points of $F_X$.
\end{definition}

\begin{theorem}[Weak Law of Large Numbers]
Let $\{X_n\}$ be i.i.d.\ with $\mathbb{E}[X_n] = \mu$. Then the sample mean $\bar{X}_n = \frac{1}{n}\sum_{i=1}^n X_i$ converges in probability to $\mu$.
\end{theorem}

\begin{remark}
This is the theoretical foundation of Monte Carlo simulation. To price a derivative by simulation, we draw $M$ independent realizations of the payoff and average them. The law of large numbers guarantees that this average converges to the true expected payoff as $M \to \infty$.
\end{remark}

\begin{theorem}[Central Limit Theorem]
Let $\{X_n\}$ be i.i.d.\ with $\mathbb{E}[X_n] = \mu$ and $\text{Var}(X_n) = \sigma^2 < \infty$. Then
\begin{align}
    Z_n = \frac{\sqrt{n}(\bar{X}_n - \mu)}{\sigma} \xrightarrow{d} \mathcal{N}(0,1).
\end{align}
\end{theorem}

\begin{remark}
The CLT quantifies the \emph{rate} at which Monte Carlo estimates converge: the standard error of the sample mean is $\sigma/\sqrt{n}$. Halving the error requires quadrupling the number of simulations. This $\mathcal{O}(n^{-1/2})$ rate is dimension-independent, which is why Monte Carlo is competitive in high-dimensional problems where grid-based methods suffer from the curse of dimensionality.
\end{remark}

\begin{example}[Monte Carlo error budget]
Suppose we estimate a European option price by averaging $M = 10{,}000$ simulated payoffs, and the sample standard deviation of the discounted payoffs is $\hat{\sigma} = 5.00$. The 95\% confidence interval for the true price is
\begin{align}
    \bar{V} \pm 1.96 \cdot \frac{\hat{\sigma}}{\sqrt{M}} = \bar{V} \pm 1.96 \cdot \frac{5.00}{\sqrt{10{,}000}} = \bar{V} \pm 0.098.
\end{align}
To tighten this to $\pm 0.01$, we would need roughly $M \approx 10^6$ simulations, or a variance reduction technique.
\end{example}


\section{Stochastic Processes}

A stock price is a family of random variables indexed by time, each representing the price at a given instant. This section develops the theory of such families: stochastic processes, their measurability properties, and the critical notions of filtration, stopping time, and martingale.

\subsection{Foundations of Stochastic Processes}

\begin{definition}
Let $(\Omega, \mathcal{F}, \mathbb{P})$ be a probability space. A \textbf{stochastic process} is a family of random variables $\{X(t)\}_{t \in T}$, where $T$ is a parameter set (typically $[0,\infty)$ or $\mathbb{N}$) and each $X(t): \Omega \to \mathbb{R}^d$.
\end{definition}

\begin{remark}
There are two complementary viewpoints. Fix $\omega \in \Omega$: the mapping $t \mapsto X(t, \omega)$ is a \textbf{sample path} (or realization), or one possible trajectory of the stock price. Fix $t$: $X(t)$ is a random variable describing the distribution of the process at that instant. Quantitative finance needs both perspectives: paths for simulation and hedging, distributions for pricing and risk.
\end{remark}

To place a probability measure on the space of all paths, we need to specify the process through its \emph{finite-dimensional distributions}, which are the joint distributions of any finite collection of time points.

\begin{definition}
Let $X(t)$ be a stochastic process. For any finite set of times $\{t_1, \ldots, t_k\} \subset T$, the \textbf{finite-dimensional distribution} is the push-forward measure
\begin{align}
    \mathbb{P}^X_{t_1,\ldots,t_k}(A) = \mathbb{P}\bigl((X(t_1), \ldots, X(t_k)) \in A\bigr), \quad A \in \mathcal{B}(\mathbb{R}^{dk}).
\end{align}
\end{definition}

\begin{theorem}[Kolmogorov Extension]
Let $\{\mathbb{P}_{t_1,\ldots,t_k}\}$ be a family of finite-dimensional probability measures satisfying the \textbf{consistency conditions}:
\begin{enumerate}
    \item \emph{Permutation invariance:} reordering time indices permutes the coordinates of the measure accordingly.
    \item \emph{Marginal consistency:} $\mathbb{P}_{t_1,\ldots,t_k}(A \times \mathbb{R}) = \mathbb{P}_{t_1,\ldots,t_{k-1}}(A)$ for Borel $A$.
\end{enumerate}
Then there exists a probability measure on the space of all paths whose finite-dimensional projections agree with the given family.
\end{theorem}

\begin{remark}
This theorem is the existence guarantee for processes defined by their distributional properties. We will use it to assert that Brownian motion exists: we specify that increments are independent and Gaussian, verify consistency, and conclude that a process with these properties exists on some probability space.
\end{remark}

\subsection{Filtrations, Stopping Times, and Martingales}

\begin{definition}
A \textbf{filtration} $\{\mathcal{F}(t)\}_{t \in T}$ is an increasing family of $\sigma$-algebras: $s \leq t \implies \mathcal{F}(s) \subseteq \mathcal{F}(t)$. A process $X(t)$ is \textbf{adapted} to $\{\mathcal{F}(t)\}$ if $X(t)$ is $\mathcal{F}(t)$-measurable for every $t$.
\end{definition}

\begin{remark}
The filtration encodes the flow of information over time. $\mathcal{F}(t)$ contains all events that can be determined by observing the process up to time $t$. Adaptedness formalizes causality: the value of an adapted process at time $t$ depends only on information available at time $t$, not on the future. Every trading strategy, every hedging decision, every stopping rule must be adapted.
\end{remark}

\begin{definition}
A random variable $\tau: \Omega \to T$ is a \textbf{stopping time} with respect to $\{\mathcal{F}(t)\}$ if $\{\tau \leq t\} \in \mathcal{F}(t)$ for all $t \in T$. A \textbf{stopped process} is $X(t \wedge \tau)$, where $\wedge$ denotes the minimum.
\end{definition}

\begin{remark}
The definition requires that the decision ``has $\tau$ occurred yet?'' can be answered at each time $t$ using only current information. Examples of stopping times: the first time a stock hits a barrier, the first time a portfolio's drawdown exceeds a threshold, the exercise time of an American option. Examples of \emph{non}-stopping times: ``the time at which the stock reaches its maximum'' - you cannot know in real time whether the current price is the eventual peak.
\end{remark}

\begin{definition}
An adapted, integrable process $M(t)$ is a \textbf{martingale} if, for all $s \leq t$,
\begin{align}
    \mathbb{E}[M(t) \mid \mathcal{F}(s)] = M(s).
\end{align}
A \textbf{submartingale} has $\mathbb{E}[M(t) \mid \mathcal{F}(s)] \geq M(s)$; a \textbf{supermartingale} has $\leq$.
\end{definition}

\begin{remark}
A martingale is a process whose best forecast of its future value, given all current information, is its current value. In financial terms:
\begin{itemize}
    \item Under the risk-neutral measure, \emph{discounted} asset prices are martingales. This is the mathematical content of no-arbitrage.
    \item Under the real-world measure, a discounted stock price is typically a submartingale (it earns a risk premium), which is precisely why we change measure.
    \item Your cumulative P\&L at a fair casino is a martingale. On average, you neither win nor lose.
\end{itemize}
\end{remark}

\begin{proposition}
A stopped martingale is still a martingale.
\end{proposition}

\begin{theorem}[Doob's Optional Stopping]
If $M(t)$ is a martingale and $\tau$ is an almost surely bounded stopping time with $\mathbb{E}[|M(\tau)|] < \infty$, then
\begin{align}
    \mathbb{E}[M(\tau)] = \mathbb{E}[M(0)].
\end{align}
\end{theorem}

\begin{remark}
In words: no stopping rule can extract a profit on average from a fair game. This underpins the impossibility of ``beating'' a martingale by clever timing, and is invoked in proving that risk-neutral prices are unique and that hedging strategies replicate payoffs.
\end{remark}

\begin{example}[Gambler's ruin as a martingale application]
A gambler starts with \$50 and bets \$1 per round on a fair coin flip. Wealth $W_n$ is a martingale. Doob's theorem implies $\mathbb{E}[W_\tau] = 50$ for any bounded stopping time $\tau$. If the gambler stops upon reaching \$100 or \$0, and the stopping time is almost surely finite (which it is for symmetric random walks), then $100 \cdot \mathbb{P}(\text{reach 100}) + 0 \cdot \mathbb{P}(\text{ruin}) = 50$, giving $\mathbb{P}(\text{reach 100}) = 1/2$. The same logic, applied to discounted derivative payoffs, yields pricing formulas.
\end{example}

\begin{definition}
Let $Z$ be an a.s.\ positive random variable with $\mathbb{E}[Z] = 1$, and let $\widetilde{\mathbb{P}}$ be the associated change of measure. The \textbf{Radon-Nikodym derivative process} is
\begin{align}
    Z(t) = \mathbb{E}[Z \mid \mathcal{F}(t)].
\end{align}
\end{definition}

\begin{proposition}
$Z(t)$ is a martingale under $\mathbb{P}$.
\end{proposition}

\begin{proof}
By the tower property: $\mathbb{E}[Z(t) \mid \mathcal{F}(s)] = \mathbb{E}[\mathbb{E}[Z \mid \mathcal{F}(t)] \mid \mathcal{F}(s)] = \mathbb{E}[Z \mid \mathcal{F}(s)] = Z(s)$.
\end{proof}

\begin{lemma}
Let $Z(t)$ be a Radon-Nikodym derivative process and $Y$ an $\mathcal{F}(t)$-measurable random variable. Then:
\begin{align}
    \widetilde{\mathbb{E}}[Y] &= \mathbb{E}[Y Z(t)], \\
    \widetilde{\mathbb{E}}[Y \mid \mathcal{F}(s)] &= \frac{1}{Z(s)}\mathbb{E}[Y Z(t) \mid \mathcal{F}(s)].
\end{align}
\end{lemma}

\begin{proof}
For the first identity: $\widetilde{\mathbb{E}}[Y] = \mathbb{E}[YZ] = \mathbb{E}[\mathbb{E}[YZ \mid \mathcal{F}(t)]] = \mathbb{E}[Y \mathbb{E}[Z \mid \mathcal{F}(t)]] = \mathbb{E}[YZ(t)]$, where the third equality uses that $Y$ is $\mathcal{F}(t)$-measurable.

For the second: let $A \in \mathcal{F}(s)$. Then $\widetilde{\mathbb{E}}[\mathbb{I}_A \cdot \frac{1}{Z(s)}\mathbb{E}[YZ(t) \mid \mathcal{F}(s)]] = \mathbb{E}[\mathbb{I}_A Y Z(t)] = \widetilde{\mathbb{E}}[\mathbb{I}_A Y]$. Since the candidate is $\mathcal{F}(s)$-measurable and satisfies partial averaging, it is the conditional expectation.
\end{proof}

\begin{remark}
This lemma is the workhorse for computing expectations under a changed measure. In the risk-neutral pricing framework, $Z(t)$ is the density process of the risk-neutral measure, and these identities let us convert between real-world and risk-neutral expectations at intermediate times.
\end{remark}

\subsection{Markov Processes}

Markov processes are the tractable workhorses of quantitative finance: their future evolution depends on the past only through the current state. This ``memoryless'' property is what makes PDEs and generators available for pricing.

\begin{definition}
A discrete-time process $\{X(n)\}_{n \in \mathbb{N}}$ on a countable state space $S$ is a \textbf{Markov chain} if
\begin{align}
    \mathbb{P}(X(n+1) = x_j \mid X(n), X(n-1), \ldots) = \mathbb{P}(X(n+1) = x_j \mid X(n) = x_i).
\end{align}
If the right-hand side does not depend on $n$, the chain is \textbf{time-homogeneous} with \textbf{transition matrix} $P = (p_{ij})$, and $\mathbb{P}(X(n) = x_j \mid X(0) = x_i) = (P^n)_{ij}$.
\end{definition}

\begin{definition}
A continuous-time process $X(t)$ taking values in $\mathbb{R}^d$ is a \textbf{Markov process} if, for all bounded Borel-measurable $f$ and $s \leq t$,
\begin{align}
    \mathbb{E}[f(X(t)) \mid \mathcal{F}(s)] = \mathbb{E}[f(X(t)) \mid X(s)].
\end{align}
The \textbf{transition probabilities} are $p(s, t, x, A) = \mathbb{P}(X(t) \in A \mid X(s) = x)$.
\end{definition}

\begin{remark}
The Markov property is both an asset and a limitation. On one hand, it enables the entire generator/PDE framework: if the future depends only on the current state, then the value of any derivative can be written as a function $V(t, X(t))$ and satisfies a PDE. On the other hand, many real-world phenomena (volatility clustering, momentum, path-dependent payoffs) violate the Markov property in the original state space. The standard fix is to \emph{enlarge the state space}: for an Asian option, the state becomes $(S(t), \int_0^t S(u)\,du)$; for stochastic volatility, the state becomes $(S(t), V(t))$. In each case, the enlarged process is Markov even though the original was not.
\end{remark}


\subsection{Brownian Motion}

Brownian motion is the canonical noise process of continuous-time finance. Its mathematical elegance and tractability make it the default model for randomness in asset prices.

\begin{theorem}[Existence of Brownian motion]
There exists a stochastic process $W(t)$, called \textbf{Brownian motion} (or the \textbf{Wiener process}), satisfying:
\begin{enumerate}
    \item $W(0) = 0$ a.s.,
    \item $W(t)$ has continuous sample paths a.s.,
    \item $W(t)$ has independent increments,
    \item $W(t) - W(s) \sim \mathcal{N}(0, t-s)$ for $s < t$.
\end{enumerate}
\end{theorem}

\begin{remark}[Proof idea]
Existence follows from the Kolmogorov extension theorem: the specified finite-dimensional distributions (multivariate Gaussian with the given covariance structure) satisfy the consistency conditions. Path continuity requires a separate argument (the Kolmogorov continuity criterion).
\end{remark}

\begin{definition}
A \textbf{filtration for Brownian motion} $\{\mathcal{F}(t)\}$ is a filtration such that $W(u) - W(t)$ is independent of $\mathcal{F}(t)$ for all $u > t$. The \textbf{natural filtration} $\mathcal{F}(t) = \sigma(W(s) : s \leq t)$ is the smallest such filtration.
\end{definition}

\begin{theorem}
Brownian motion is a martingale.
\end{theorem}

\begin{proof}
$\mathbb{E}[W(t) \mid \mathcal{F}(s)] = \mathbb{E}[(W(t)-W(s)) + W(s) \mid \mathcal{F}(s)] = 0 + W(s)$, using independence of the increment and the $\mathcal{F}(s)$-measurability of $W(s)$.
\end{proof}

\begin{theorem}
Brownian motion is a Markov process.
\end{theorem}

\begin{theorem}[Quadratic variation]
Brownian motion accumulates quadratic variation at rate one:
\begin{align}
    [W,W](T) = \lim_{\|\Pi\| \to 0} \sum_j [W(t_{j+1}) - W(t_j)]^2 = T.
\end{align}
For $d$-dimensional Brownian motion, $[W_i, W_j](t) = \delta_{ij} t$.
\end{theorem}

\begin{remark}
Quadratic variation is the property that distinguishes stochastic calculus from ordinary calculus. A smooth function has zero quadratic variation, so the $(dt)^2$ and higher-order terms in Taylor expansions vanish. For Brownian motion, $(dW)^2 = dt$ does \emph{not} vanish, producing the extra second-order term in It\^o's lemma. Financially, quadratic variation is directly related to realized variance: the sum of squared returns converges to the integrated variance of the price process.
\end{remark}

\begin{theorem}[L\'evy's characterization]
Let $M(t)$ be a continuous martingale with $M(0)=0$ and $[M,M](t) = t$. Then $M(t)$ is a Brownian motion. In the multidimensional case, the condition becomes $[M_i, M_j](t) = \delta_{ij} t$.
\end{theorem}

\begin{remark}
L\'evy's theorem is a powerful identification tool: to show that a process is Brownian motion, we need only check two properties: martingale and quadratic variation. We will use it repeatedly when constructing new Brownian motions under changed measures (Girsanov's theorem).
\end{remark}

\begin{example}[Scaling property and the square-root-of-time rule]
For any $a > 0$, the process $\widetilde{W}(t) = \frac{1}{a}W(a^2 t)$ is also a Brownian motion. Verify via L\'evy:

\emph{Martingale:} The natural filtration of $\widetilde{W}$ is $\mathcal{F}(s) = \sigma(W(u), u \leq a^2s)$. Then $\mathbb{E}[\widetilde{W}(t) \mid \mathcal{F}(s)] = \frac{1}{a}\mathbb{E}[W(a^2t) \mid \sigma(W(u), u \leq a^2s)] = \frac{1}{a}W(a^2s) = \widetilde{W}(s)$. \checkmark

\emph{Quadratic variation:} $[\widetilde{W}, \widetilde{W}](t) = \frac{1}{a^2}[W,W](a^2t) = \frac{1}{a^2} \cdot a^2 t = t$. \checkmark

This scaling property is the mathematical origin of the \textbf{square-root-of-time rule}: if daily volatility is $\sigma_{\text{daily}}$, then $T$-day volatility is $\sigma_{\text{daily}}\sqrt{T}$. In particular, the industry convention for annualizing daily volatility is
\begin{align}
    \sigma_{\text{annual}} \approx \sqrt{252}\, \sigma_{\text{daily}},
\end{align}
where 252 is the approximate number of trading days in a year.
\end{example}

\begin{proposition}[Reflection principle]
Let $M(t) = \max_{0 \leq s \leq t} W(s)$ and $\tau_m = \inf\{t : W(t) = m\}$. Then:
\begin{align}
    \mathbb{P}(M(t) \geq a) &= 2\,\mathbb{P}(W(t) \geq a), \\
    \mathbb{P}(\tau_m \leq t,\, W(t) \leq w) &= \mathbb{P}(W(t) \geq 2m - w).
\end{align}
\end{proposition}

\begin{remark}
The reflection principle is the key to pricing barrier options. An up-and-out call, for instance, requires computing probabilities involving both the terminal value and the running maximum of the stock price. The joint distribution of $(W(t), M(t))$ is derived from reflection, and transforms these pricing problems into exercises in Gaussian integration.
\end{remark}


\section{Stochastic Calculus}

Brownian motion, despite being the engine of continuous-time models, is famously nowhere differentiable. The classical chain rule breaks down. This section builds a consistent calculus for processes driven by Brownian motion, culminating in the deep connection between stochastic differential equations (SDEs) and partial differential equations (PDEs) that makes derivative pricing computationally feasible.

\subsection{It\^o Calculus}

\begin{definition}[It\^o integral]
Let $W(t)$ be Brownian motion with filtration $\{\mathcal{F}(t)\}$ on $[0,T]$, and let $\Delta(t)$ be an adapted process with $\mathbb{E}[\int_0^T \Delta(t)^2\, dt] < \infty$. The \textbf{It\^o integral} is defined as the $L^2$-limit of left-endpoint Riemann sums:
\begin{align}
    I(T) = \int_0^T \Delta(t)\, dW(t) \coloneqq \lim_{\|\Pi\| \to 0} \sum_j \Delta(t_j)\bigl[W(t_{j+1}) - W(t_j)\bigr].
\end{align}
\end{definition}

\begin{remark}
Two points deserve emphasis. First, the integrand is evaluated at the \emph{left} endpoint $t_j$, not the right or midpoint. This is not a convention - it is forced by the requirement that the integrand be adapted (known at the time of evaluation). Using a midpoint would produce the Stratonovich integral, which has different properties. Second, the limit is in the $L^2$ sense, not pointwise; convergence is in expectation, not path-by-path.
\end{remark}

\begin{theorem}[Properties of the It\^o integral]
The It\^o integral $I(t) = \int_0^t \Delta(u)\, dW(u)$ satisfies:
\begin{enumerate}
    \item $I(t)$ is a martingale (in particular, $\mathbb{E}[I(t)] = 0$).
    \item $I(t)$ has continuous paths and is $\mathcal{F}(t)$-measurable.
    \item Quadratic variation: $[I,I](t) = \int_0^t \Delta(u)^2\, du$.
    \item It\^o isometry: $\mathbb{E}[I(t)^2] = \mathbb{E}\!\left[\int_0^t \Delta(u)^2\, du\right]$.
\end{enumerate}
\end{theorem}

\begin{remark}
The martingale property of the It\^o integral is why stochastic integrals have zero drift. The It\^o isometry converts a second-moment calculation into a deterministic integral.
\end{remark}

\begin{definition}[It\^o process]
An \textbf{It\^o process} $X(t)$ with deterministic initial condition $X(0)$ is
\begin{align}
    X(t) = X(0) + \int_0^t \Theta(u)\, du + \int_0^t \Delta(u)\, dW(u),
\end{align}
or in differential notation,
\begin{align}
    dX(t) = \Theta(t)\, dt + \Delta(t)\, dW(t),
\end{align}
where $\Theta(t)$ is the \textbf{drift} and $\Delta(t)$ is the \textbf{diffusion coefficient}.
\end{definition}

\begin{lemma}[It\^o's lemma]
Let $f(t,x)$ have continuous partial derivatives $f_t, f_x, f_{xx}$, and let $X(t)$ be an It\^o process. Then:
\begin{align}
    df(t, X(t)) = f_t\, dt + f_x\, dX + \frac{1}{2}f_{xx}\, (dX)^2,
\end{align}
where the multiplication rules are $dt \cdot dt = dW \cdot dt = 0$ and $dW \cdot dW = dt$.

For a product of two It\^o processes:
\begin{align}
    d(X(t)Y(t)) = X\, dY + Y\, dX + dX \cdot dY.
\end{align}

In multiple dimensions, with processes $X_1, \ldots, X_d$:
\begin{align}
    df(t, \vec{X}(t)) = f_t\, dt + \sum_i f_{x_i}\, dX_i + \frac{1}{2}\sum_{i,j} f_{x_ix_j}\, dX_i\, dX_j.
\end{align}
\end{lemma}

\begin{remark}
It\^o's lemma is the chain rule of stochastic calculus, with an extra second-order correction term that has no analogue in ordinary calculus. This correction arises because Brownian motion has non-zero quadratic variation. The term $\frac{1}{2}\sigma^2 S^2 f_{SS}$ in the Black-Scholes PDE, the $-\frac{1}{2}\sigma^2$ drift adjustment in geometric Brownian motion, and the convexity correction in interest rate models all trace back to this single lemma.
\end{remark}

\begin{example}[Geometric Brownian motion]
Suppose the stock price satisfies the SDE
\begin{align}
    dS(t) = \mu S(t)\, dt + \sigma S(t)\, dW(t),
\end{align}
where $\mu$ is the expected return and $\sigma$ is the volatility (both constant for now). To solve, try $f(S) = \log S$. By It\^o's lemma:
\begin{align}
    d(\log S) &= \frac{1}{S}\, dS - \frac{1}{2}\frac{1}{S^2}(dS)^2 \\
    &= \frac{1}{S}\bigl(\mu S\, dt + \sigma S\, dW\bigr) - \frac{1}{2}\frac{1}{S^2}\sigma^2 S^2\, dt \\
    &= \Bigl(\mu - \frac{1}{2}\sigma^2\Bigr) dt + \sigma\, dW(t).
\end{align}
This is a deterministic drift plus a stochastic integral, and both sides can be integrated directly:
\begin{align}
    S(t) = S(0)\exp\Bigl(\bigl(\mu - \tfrac{1}{2}\sigma^2\bigr)t + \sigma W(t)\Bigr).
\end{align}
This is \textbf{geometric Brownian motion (GBM)}. The term $-\frac{1}{2}\sigma^2$ is the It\^o correction, sometimes called the ``convexity adjustment.'' Without it (if we naively applied ordinary calculus), the expected value of $S(t)$ would not equal $S(0)e^{\mu t}$.


\begin{theorem}[Girsanov]
Let $W(t)$ be Brownian motion on $(\Omega, \mathcal{F}, \mathbb{P})$ with filtration $\mathcal{F}(t)$ over $[0,T]$. Let $\Theta(t)$ be an adapted process and define:
\begin{align}
    Z(t) &= \exp\Bigl(-\int_0^t \Theta(u)\, dW(u) - \frac{1}{2}\int_0^t \Theta(u)^2\, du\Bigr), \\
    \widetilde{W}(t) &= W(t) + \int_0^t \Theta(u)\, du.
\end{align}
If $\mathbb{E}\bigl[\int_0^T \Theta(u)^2 Z(u)^2\, du\bigr] < \infty$, then $\mathbb{E}[Z(T)] = 1$, and $\widetilde{W}(t)$ is a Brownian motion under $\widetilde{\mathbb{P}}$ defined by $d\widetilde{\mathbb{P}}/d\mathbb{P} = Z(T)$.
\end{theorem}

\begin{remark}[Proof idea]
One verifies that $\widetilde{W}$ is a martingale under $\widetilde{\mathbb{P}}$ with the correct quadratic variation, then applies L\'evy's characterization.
\end{remark}

\begin{remark}[Financial meaning]
Girsanov's theorem is the mathematical mechanism behind risk-neutral pricing. Under the real-world measure $\mathbb{P}$, the stock has drift $\mu$ and satisfies $dS = \mu S\, dt + \sigma S\, dW$. By choosing $\Theta(t) = (\mu - r)/\sigma$ (the \textbf{market price of risk}), Girsanov produces a new measure $\widetilde{\mathbb{P}}$ under which $dS = rS\, dt + \sigma S\, d\widetilde{W}$, and the stock earns only the risk-free rate. Under this measure, discounted prices are martingales, and pricing reduces to computing expectations. The drift changes; the volatility does not. This is a deep point: hedging depends on volatility, not drift, and Girsanov preserves volatility.
\end{remark}

\begin{example}[Girsanov in action]
Under $\mathbb{P}$, a stock has drift $\mu = 10\%$ and volatility $\sigma = 25\%$, with risk-free rate $r = 3\%$. The market price of risk is $\Theta = (0.10 - 0.03)/0.25 = 0.28$. Under $\widetilde{\mathbb{P}}$, the stock dynamics become
\begin{align}
    dS = 0.03 \cdot S\, dt + 0.25 \cdot S\, d\widetilde{W}.
\end{align}
The drift has been reduced from 10\% to 3\%, exactly the risk-free rate. The ``removed'' drift of 7\% is the equity risk premium, and its removal is encoded in the Radon-Nikodym derivative $Z(T)$.
\end{example}

\begin{theorem}[Martingale representation]
Let $M(t)$ be a martingale with respect to the Brownian filtration $\{\mathcal{F}(t)\}$. Then there exists an adapted process $\Gamma(u)$ such that
\begin{align}
    M(t) = M(0) + \int_0^t \Gamma(u)\, dW(u).
\end{align}
Under $\widetilde{\mathbb{P}}$, if $\widetilde{M}(t)$ is a martingale, there exists $\widetilde{\Gamma}(u)$ with $\widetilde{M}(t) = \widetilde{M}(0) + \int_0^t \widetilde{\Gamma}(u)\, d\widetilde{W}(u)$.
\end{theorem}

\begin{remark}
This theorem guarantees that every martingale adapted to a Brownian filtration can be written as a stochastic integral, so there is no ``hidden'' source of randomness. Financially, it means that every contingent claim can be replicated by trading in the stock and money market, which is the mathematical statement of \emph{market completeness}. The integrand $\Gamma(u)$ is (after rescaling) the delta hedge.
\end{remark}

\subsection{The SDE-PDE Connection}

The deepest structural result in this chapter links stochastic differential equations to partial differential equations. This connection is what makes derivative pricing computationally feasible: instead of simulating thousands of paths and averaging (Monte Carlo), one can solve a single PDE on a grid (finite differences).

\begin{proposition}
Consider the SDE 
\begin{align}
    dX(u) = \beta(u, X(u))\, du + \gamma(u, X(u))\, dW(u).
\end{align} 
Under standard regularity conditions (Lipschitz continuity and linear growth), a unique strong solution exists, and $X(t)$ is a Markov process.
\end{proposition}

\begin{theorem}[Feynman-Kac]
Let $X(t)$ solve the above SDE, and let $h$ be a Borel-measurable function. Define
\begin{align}
    g(t,x) = \mathbb{E}[h(X(T)) \mid X(t) = x].
\end{align}
Then $g$ satisfies the PDE
\begin{align}
    \frac{\partial g}{\partial t} + \beta(t,x)\frac{\partial g}{\partial x} + \frac{1}{2}\gamma(t,x)^2 \frac{\partial^2 g}{\partial x^2} = 0, \quad g(T,x) = h(x).
\end{align}
With discounting at rate $r$, if $f(t,x) = \mathbb{E}[e^{-r(T-t)}h(X(T)) \mid X(t) = x]$, then
\begin{align}
    \frac{\partial f}{\partial t} + \beta(t,x)\frac{\partial f}{\partial x} + \frac{1}{2}\gamma(t,x)^2 \frac{\partial^2 f}{\partial x^2} = rf, \quad f(T,x) = h(x).
\end{align}
\end{theorem}

\begin{remark}[Proof idea]
Show that $e^{-rt}f(t,X(t))$ is a martingale. Apply It\^o's lemma to compute $d(e^{-rt}f)$ and impose that the $dt$ coefficient is zero (otherwise the drift would destroy the martingale property). This produces the PDE.
\end{remark}



The Feynman-Kac theorem is a special case of a much more general relationship between Markov processes and differential operators. This framework of semigroups and generators provides a unified language for pricing across different models.

\begin{definition}
For a time-homogeneous Markov process $X(t)$, define the \textbf{semigroup operators}
\begin{align}
    (P_t f)(x) = \mathbb{E}[f(X(t)) \mid X(0) = x].
\end{align}
These satisfy $P_0 = \text{id}$ and $P_{s+t} = P_s P_t$. For inhomogeneous processes, there is a two-parameter family $P_{t,T}$ with $P_{t,T} = P_{t,s} P_{s,T}$.
\end{definition}

\begin{proof}[Proof of the semigroup property]
$P_0 f(x) = \mathbb{E}[f(X(0)) \mid X(0) = x] = f(x)$. For the composition: $P_{s+t}f(x) = \mathbb{E}[f(X(s+t)) \mid X(0) = x]$. By the tower property and Markov property: $= \mathbb{E}[\mathbb{E}[f(X(s+t)) \mid X(s)] \mid X(0) = x] = \mathbb{E}[(P_t f)(X(s)) \mid X(0) = x] = P_s(P_t f)(x)$.
\end{proof}

\begin{definition}
The \textbf{(infinitesimal) generator} of the semigroup is
\begin{align}
    \mathcal{L}f(x) = \lim_{t \downarrow 0} \frac{P_t f(x) - f(x)}{t} = \lim_{t \downarrow 0} \frac{\mathbb{E}[f(X(t)) \mid X(0) = x] - f(x)}{t},
\end{align}
defined for all $f$ in the domain of $\mathcal{L}$. The \textbf{adjoint} $\mathcal{L}^*$ acts on densities/measures and satisfies $\langle \mathcal{L}f, p \rangle = \langle f, \mathcal{L}^* p \rangle$.
\end{definition}

\begin{proposition}
For the SDE $dX = \beta(t,x)\, dt + \gamma(t,x)\, dW$, the generator acting on spatial variables is
\begin{align}
    \mathcal{L}_t f(x) = \beta(t,x) f'(x) + \frac{1}{2}\gamma(t,x)^2 f''(x),
\end{align}
and the adjoint is
\begin{align}
    \mathcal{L}^*_t p(t,x) = -\frac{\partial}{\partial x}\bigl(\beta(t,x) p\bigr) + \frac{1}{2}\frac{\partial^2}{\partial x^2}\bigl(\gamma(t,x)^2 p\bigr).
\end{align}
\end{proposition}

\begin{remark}[Proof idea]
Apply It\^o's lemma to $f(X(u))$, take conditional expectation, and divide by $t$. The $dW$ integral vanishes in expectation; the $dt$ integral produces the generator. The adjoint follows by integration by parts.
\end{remark}

\begin{definition}
A Markov process is \textbf{affine} if
\begin{align}
    \mathbb{E}[e^{uX(t)} \mid X(0) = x] = e^{A(t,u) + B(t,u)x}
\end{align}
for deterministic functions $A, B$, or equivalently, if $\mathcal{L}e^{ux} = (a(u) + b(u)x)e^{ux}$.
\end{definition}

The generator framework yields two fundamental PDEs governing Markov processes.

\begin{theorem}[Kolmogorov backward equation]
Let $u(t,x) = \mathbb{E}[f(X(T)) \mid X(t) = x]$. Then
\begin{align}
    \frac{\partial u}{\partial t} + \mathcal{L}_t u = 0, \quad u(T,x) = f(x).
\end{align}
With discounting: $u(t,x) = \mathbb{E}[e^{-r(T-t)}f(X(T)) \mid X(t) = x]$ satisfies
\begin{align}
    \frac{\partial u}{\partial t} + \mathcal{L}_t u - ru = 0, \quad u(T,x) = f(x).
\end{align}
The backward equation governs the evolution of \emph{observables} (prices, values).
\end{theorem}

\begin{theorem}[Kolmogorov forward equation / Fokker-Planck]
The transition density $p(t,T,x,y)$ of $X$ satisfies
\begin{align}
    \frac{\partial p}{\partial T} = \mathcal{L}^*_T p,
\end{align}
or explicitly,
\begin{align}
    \frac{\partial p}{\partial T} = -\frac{\partial}{\partial y}\bigl(\beta(T,y)\, p\bigr) + \frac{1}{2}\frac{\partial^2}{\partial y^2}\bigl(\gamma(T,y)^2\, p\bigr),
\end{align}
with initial condition $\lim_{T \downarrow t} p(t,T,x,y) = \delta(y-x)$. The forward equation governs the evolution of \emph{probability distributions}.
\end{theorem}

\begin{proof}
From the backward equation, $\frac{\partial}{\partial t}\int f(x) p(t,x)\, dx = \int \mathcal{L}_t(f(x))\, p(t,x)\, dx$. Applying the definition of the adjoint to the right-hand side: $= \int f(x)\, \mathcal{L}^*_t(p(t,x))\, dx$. Since this holds for all test functions $f$, we obtain $\frac{\partial p}{\partial t} = \mathcal{L}^*_t p$.
\end{proof}

\begin{remark}
To summarize the duality: the backward equation propagates value functions backward from terminal conditions (``given the payoff at expiry, what is the price today?''), while the forward equation propagates probability densities forward from initial conditions (``given the current state, where might the process be in the future?''). Both are determined by the same generator $\mathcal{L}$.
\end{remark}

\section{Jump Diffusion}

Pure diffusion models imply that prices move continuously: there are no sudden jumps, gaps, or discontinuities. This is mathematically convenient but empirically inadequate. Markets experience flash crashes, earnings surprises, central bank announcements, and credit events that cause instantaneous price jumps. Modifying the price process to include jumps captures these features while retaining much of the diffusion framework.

\subsection{L\'evy Processes and Poisson Jumps}

\begin{definition}
A function $f$ is \textbf{c\`adl\`ag} (\emph{continue \`a droite, limite \`a gauche}) if it is right-continuous with left limits at every point. C\`adl\`ag paths allow for jumps: the function value equals the right limit, and the left limit captures the ``pre-jump'' value.
\end{definition}

\begin{definition}
A stochastic process $X(t)$ on a filtered probability space is a \textbf{L\'evy process} if:
\begin{enumerate}
    \item $X(0) = 0$ a.s.,
    \item paths are c\`adl\`ag a.s.,
    \item $X(t)$ is adapted,
    \item increments $X(t) - X(s)$ are independent of $\mathcal{F}(s)$,
    \item increments are stationary: $X(t) - X(s) \overset{d}{=} X(t-s)$,
    \item continuity in probability: $\lim_{h \to 0} \mathbb{P}(|X(t+h) - X(t)| > \epsilon) = 0$ for all $\epsilon > 0$.
\end{enumerate}
\end{definition}

\begin{remark}
L\'evy processes generalize both Brownian motion (continuous, Gaussian increments) and Poisson processes (discontinuous, integer-valued increments). The general theory shows that any L\'evy process decomposes into a Brownian component, a drift, and a jump component.
\end{remark}

\begin{theorem}
There exists a stochastic process $N(t)$, called the \textbf{Poisson process} with intensity $\lambda > 0$, such that $N(0) = 0$ a.s., $N(t)$ has c\`adl\`ag paths, independent increments, and $N(t) - N(s) \sim \text{Pois}(\lambda(t-s))$.
\end{theorem}

\begin{proposition}
The sum of independent Poisson processes with intensities $\lambda_1, \lambda_2$ is a Poisson process with intensity $\lambda_1 + \lambda_2$, and conversely.
\end{proposition}

\begin{definition}
Let $N(t)$ be a Poisson process with intensity $\lambda$, and let $(Y_k)_{k \geq 1}$ be i.i.d.\ random variables with mean $\beta$, independent of $N(t)$. The process
\begin{align}
    Q(t) = \sum_{k=1}^{N(t)} Y_k
\end{align}
is a \textbf{compound Poisson process}. It counts random arrivals and assigns each one a random size.
\end{definition}

\begin{remark}
Think of $N(t)$ as counting ``shocks'' (earnings announcements, credit events, policy changes) and $Y_k$ as the size of each shock. The compound Poisson process models both the timing and magnitude of jumps.
\end{remark}

\begin{definition}
The \textbf{compensated Poisson process} $M(t) = N(t) - \lambda t$ is a martingale: its expected increment is zero. For a compound Poisson process, the compensated martingale is $Q(t) - \lambda \beta t$.
\end{definition}

\begin{remark}
Compensation subtracts the predictable drift, leaving a zero-mean martingale. This parallels the role of Girsanov's theorem for diffusions: the ``surprise'' part of the jump process drives the risk, while the predictable part is absorbed into pricing adjustments.
\end{remark}

\subsection{Jump Diffusions and Extended It\^o Calculus}

\begin{theorem}[L\'evy-It\^o decomposition]
Every L\'evy process admits the decomposition
\begin{align}
    X(t) = bt + \sigma W(t) + \int_{|x| < 1} x\, \widetilde{N}(t, dx) + \int_{|x| \geq 1} x\, N(t, dx),
\end{align}
where $b \in \mathbb{R}$, $\sigma \geq 0$, $W(t)$ is Brownian motion, $N(t, dx)$ is a Poisson random measure, $\widetilde{N}$ is its compensated version, and $\nu(dx)$ is the \textbf{L\'evy measure} satisfying $\int_{\mathbb{R}} \min(1, x^2)\, \nu(dx) < \infty$.
\end{theorem}

\begin{remark}
The L\'evy measure $\nu$ completely characterizes the jump behavior: $\nu(A)$ gives the expected number of jumps per unit time with sizes in the set $A$. Large jumps ($|x| \geq 1$) occur finitely often and are summed directly; small jumps ($|x| < 1$) may occur infinitely often but have finite total quadratic variation, so they are compensated to ensure convergence.
\end{remark}

\begin{definition}
A \textbf{jump diffusion} combines continuous Brownian fluctuations with discrete Poisson jumps:
\begin{align}
    X(t) = X(0) + \int_0^t \mu(s)\, ds + \int_0^t \sigma(s)\, dW(s) + \sum_{k=1}^{N(t)} Y_k,
\end{align}
where $W(t)$, $N(t)$, and the $Y_k$ are mutually independent.
\end{definition}

\begin{theorem}[It\^o's lemma with jumps]
Let $X(t)$ be a jump diffusion and $f \in C^2(\mathbb{R})$. Then
\begin{align}
    df(X(t)) &= f'(X(t^-))\, dX^c(t) + \frac{1}{2}f''(X(t^-))\, \sigma^2(t, X(t^-))\, dt \\
    &\quad + \sum_{0 < s \leq t} \Bigl[f(X(s)) - f(X(s^-)) - f'(X(s^-))\Delta X(s)\Bigr],
\end{align}
where $X^c$ denotes the continuous part and $\Delta X(s) = X(s) - X(s^-)$ is the jump size.
\end{theorem}

\begin{remark}
The jump correction replaces the second-order Taylor term with the exact nonlinear change $f(X(s)) - f(X(s^-))$ at each jump time. For small jumps, this is approximately $\frac{1}{2}f''(\Delta X)^2$; for large jumps, the full nonlinear effect matters. This is why jump-diffusion option pricing formulas involve expectations over the jump size distribution rather than simple Gaussian integrals.
\end{remark}

\begin{example}[Merton's jump-diffusion model]
In Merton's (1976) model, the stock price follows
\begin{align}
    \frac{dS}{S(t^-)} = (\mu - \lambda \beta)\, dt + \sigma\, dW(t) + J\, dN(t),
\end{align}
where $J = e^Y - 1$ for $Y \sim \mathcal{N}(\mu_J, \sigma_J^2)$, and $\beta = \mathbb{E}[J] = e^{\mu_J + \sigma_J^2/2} - 1$. The drift is adjusted by $-\lambda\beta$ to compensate for the expected jump, ensuring the discounted price remains a martingale under the risk-neutral measure.

Between jumps, the stock follows GBM. At each Poisson arrival time, the price is multiplied by $e^Y$. The resulting return distribution has heavier tails than the pure GBM model, capturing the empirical observation that extreme daily returns occur more frequently than a Gaussian distribution would predict.

The European call price under this model decomposes as a Poisson-weighted sum of Black-Scholes prices:
\begin{align}
    C_{\text{Merton}} = \sum_{k=0}^{\infty} \frac{e^{-\lambda' T}(\lambda' T)^k}{k!}\, C_{\text{BS}}(S_0, K, T, r_k, \sigma_k),
\end{align}
where $\lambda' = \lambda(1+\beta)$ and $r_k, \sigma_k$ are adjusted for the number of jumps. This formula has a clear interpretation: condition on the number of jumps $k$, compute the Black-Scholes price given $k$ jumps, and average over $k$.
\end{example}

\begin{remark}
Jump diffusions provide a mathematically tractable middle ground between the oversimplified GBM model and more complex L\'evy-driven processes. They capture sudden price moves, produce implied volatility smiles (even for a single maturity), and allow semi-analytical pricing for European options. The trade-off is that jumps destroy the completeness of the market: there is no longer a unique risk-neutral measure or a perfect hedge. This echoes a recurring theme in quantitative finance.
\end{remark}

\chapter{Risk-Neutral Formulation}

The mathematical foundations of the previous chapter, like stochastic calculus, Girsanov's theorem, and the martingale representation theorem, are powerful but abstract. This chapter puts them to work. The central question is: \emph{What is a derivative worth?}

The answer rests on a single economic principle: \textbf{no arbitrage}. If two portfolios produce the same future payoff in every state of the world, they must have the same price today. This simple statement, combined with the stochastic calculus machinery, yields a complete pricing framework. The key insight is that no-arbitrage is equivalent to the existence of a probability measure (the \emph{risk-neutral measure}) under which discounted asset prices are martingales. Under this measure, every derivative price is simply a discounted expected payoff.

The chapter proceeds in two stages. First, we develop risk-neutral pricing in the single-asset and multi-asset settings, prove the fundamental theorems of asset pricing, and establish the conditions under which hedging is possible. Second, we generalize the framework via change of num\'eraire, showing that the money market account is just one of many possible reference assets, and that choosing the right num\'eraire can dramatically simplify specific pricing problems.


\section{Risk-Neutral Pricing}

\subsection{One Asset}

We begin with the simplest setting: one stock, one Brownian motion, one risk-free rate.

\begin{definition}
A \textbf{portfolio} (or \textbf{wealth process}) $X(t)$ represents the value of a self-financing combination of a stock $S(t)$ and a money market account earning the risk-free rate $R(t)$. Let $\Delta(t)$ be the number of shares of stock held at time $t$, adapted to the filtration. The remaining wealth $X(t) - \Delta(t)S(t)$ is invested in the money market. The self-financing condition gives the portfolio differential:
\begin{align}
    dX(t) = \Delta(t)\,dS(t) + R(t)\bigl(X(t) - \Delta(t)S(t)\bigr)\,dt.
\end{align}
\end{definition}

\begin{remark}
The self-financing condition means that changes in portfolio value come \emph{only} from capital gains on the stock and interest on the money market position; no external cash is added or withdrawn. This is the mathematical embodiment of ``no free lunch,'' that any trading strategy must fund itself.
\end{remark}

To compare values at different times, we need to account for the time value of money.

\begin{definition}
For an adapted interest rate process $R(t)$, the \textbf{discount process} is
\begin{align}
    D(t) = \exp\!\Bigl(-\int_0^t R(u)\,du\Bigr),
\end{align}
with differential $dD(t) = -R(t)D(t)\,dt$. The \textbf{discounted stock price} is $D(t)S(t)$, which removes the effect of the risk-free rate and isolates the ``excess'' dynamics of the stock.
\end{definition}

When the stock follows geometric Brownian motion with drift $\mu(t)$ and volatility $\sigma(t)$, the discounted stock price has an explicit form.

\begin{proposition}
The discounted stock price satisfies
\begin{align}
    D(t)S(t) = S(0)\exp\!\Bigl(\int_0^t\bigl(\mu(s) - R(s) - \tfrac{1}{2}\sigma(s)^2\bigr)\,ds + \int_0^t \sigma(s)\,dW(s)\Bigr),
\end{align}
with differential
\begin{align}
    d\bigl(D(t)S(t)\bigr) = \sigma(t)D(t)S(t)\bigl[\Theta(t)\,dt + dW(t)\bigr],
\end{align}
where we write the \textbf{market price of risk} as
\begin{align}
    \Theta(t) = \frac{\mu(t) - R(t)}{\sigma(t)}.
\end{align}
\end{proposition}

\begin{remark}
The market price of risk $\Theta(t)$ measures the excess return per unit of volatility (the Sharpe ratio of the stock). It quantifies how much compensation the market demands for bearing one unit of diffusive risk. Under the real-world measure, the discounted stock price has a non-zero drift proportional to $\Theta$: it is \emph{not} a martingale. The entire point of what follows is to remove this drift.
\end{remark}

The key step is to apply Girsanov's theorem with the market price of risk as the drift-removal process. This changes the probability measure, not the volatility.

\begin{definition}
Apply Girsanov's theorem with $\Theta(t) = (\mu(t)-R(t))/\sigma(t)$ to define
\begin{align}
    \widetilde{W}(t) = W(t) + \int_0^t \Theta(u)\,du.
\end{align}
Under the new measure $\widetilde{\mathbb{P}}$ (with Radon-Nikodym derivative $Z(T) = \exp(-\int_0^T\Theta\,dW - \frac{1}{2}\int_0^T\Theta^2\,du)$), the process $\widetilde{W}(t)$ is a Brownian motion, and the discounted stock price becomes
\begin{align}
    d\bigl(D(t)S(t)\bigr) = \sigma(t)D(t)S(t)\,d\widetilde{W}(t).
\end{align}
The measure $\widetilde{\mathbb{P}}$ is called the \textbf{risk-neutral measure} (or \textbf{equivalent martingale measure}).
\end{definition}

\begin{remark} Under $\mathbb{P}$, the stock earns $\mu$; under $\widetilde{\mathbb{P}}$, it earns $R$. The excess return has been absorbed into the measure change. However, under both measures, the diffusion coefficient is $\sigma(t)S(t)$. This is why hedging depends on volatility (which is observable) rather than drift (which is not). The discounted stock price $D(t)S(t)$ is now a martingale under $\widetilde{\mathbb{P}}$, as its $dt$ coefficient is zero.
\end{remark}

The martingale property extends from the stock to any self-financing portfolio.

\begin{proposition}
Discounted portfolio values are martingales under the risk-neutral measure:
\begin{align}
    d\bigl(D(t)X(t)\bigr) = \Delta(t)\sigma(t)D(t)S(t)\,d\widetilde{W}(t).
\end{align}
In particular, every self-financing portfolio earns the risk-free rate on average under $\widetilde{\mathbb{P}}$.
\end{proposition}

\begin{remark}
This is the mathematical formalization of ``risk-neutral'': under $\widetilde{\mathbb{P}}$, investors are indifferent to risk, demanding no premium beyond the risk-free rate. Of course, investors in the real world \emph{do} demand risk premia; the risk-neutral measure is a computational device. Prices computed as discounted expected payoffs under $\widetilde{\mathbb{P}}$ are \emph{arbitrage-free} prices, regardless of investors' actual risk preferences.
\end{remark}

With the martingale property in hand, pricing becomes an exercise in computing conditional expectations.

\begin{theorem}[Risk-neutral pricing]
Let $V(T)$ be an $\mathcal{F}(T)$-measurable payoff. Under the risk-neutral measure $\widetilde{\mathbb{P}}$, its value at time $t < T$ is
\begin{align}
    V(t) = \widetilde{\mathbb{E}}\!\left[e^{-\int_t^T R(u)\,du}\,V(T)\;\middle|\;\mathcal{F}(t)\right].
\end{align}
\end{theorem}

\begin{proof}
Under $\widetilde{\mathbb{P}}$, discounted asset prices are martingales. If a replicating portfolio $X(t)$ exists with $X(T) = V(T)$ a.s., then $D(t)X(t) = \widetilde{\mathbb{E}}[D(T)X(T) \mid \mathcal{F}(t)] = \widetilde{\mathbb{E}}[D(T)V(T) \mid \mathcal{F}(t)]$. Since $V(t) = X(t)$ (the portfolio replicates the derivative), dividing by $D(t)$ gives $V(t) = \widetilde{\mathbb{E}}[e^{-\int_t^T R(u)\,du}\,V(T) \mid \mathcal{F}(t)]$.
\end{proof}

\begin{remark}
This is the master formula of derivative pricing. Every closed-form option price (Black-Scholes, bond pricing formulas, cap/floor formulas) is a specific evaluation of this expectation for a particular payoff $V(T)$ and a particular model for the dynamics of $S$ and $R$ under $\widetilde{\mathbb{P}}$. When the expectation cannot be evaluated analytically, it is computed by Monte Carlo simulation (averaging over simulated paths under $\widetilde{\mathbb{P}}$) or by solving the associated PDE (via Feynman-Kac).
\end{remark}

The risk-neutral pricing formula assumed the existence of a replicating portfolio. The next result establishes when such a portfolio exists.

\begin{proposition}[Existence of hedges]
Consider a payoff $V(T)$ that is $\mathcal{F}(T)$-measurable. A replicating portfolio $X(t)$ with $X(T) = V(T)$ a.s.\ exists if one chooses
\begin{align}
    X(0) = V(0), \qquad \Delta(t) = \frac{\widetilde{\Gamma}(t)}{\sigma(t)D(t)S(t)},
\end{align}
where $\widetilde{\Gamma}(t)$ is the integrand from the martingale representation of $D(t)V(t)$ under $\widetilde{\mathbb{P}}$. This requires that $\sigma(t) \neq 0$ (the stock is not deterministic) and that Brownian motion is the only source of uncertainty. A market in which every contingent claim can be replicated is called \textbf{complete}.
\end{proposition}

\begin{proof}
Under $\widetilde{\mathbb{P}}$, the process $D(t)V(t) = \widetilde{\mathbb{E}}[D(T)V(T) \mid \mathcal{F}(t)]$ is a martingale. By the martingale representation theorem, there exists an adapted process $\widetilde{\Gamma}(t)$ such that
\begin{align}
    D(t)V(t) = V(0) + \int_0^t \widetilde{\Gamma}(u)\,d\widetilde{W}(u).
\end{align}
The discounted portfolio satisfies $D(t)X(t) = X(0) + \int_0^t \Delta(u)\sigma(u)D(u)S(u)\,d\widetilde{W}(u)$. Matching initial conditions and integrands at all times gives the stated formulas.
\end{proof}

\begin{remark}
The martingale representation theorem guarantees that every risk-neutral martingale can be written as a stochastic integral against Brownian motion, and the hedge ratio translates that abstract integrand into a concrete stock position.

In practice, continuous rebalancing is impossible: one trades at discrete times, incurring hedging error. The error is proportional to the gamma $\partial^2 V/\partial S^2$ of the derivative and the square of the rebalancing interval, which is why gamma-heavy positions (short-dated options near the strike) are the hardest to hedge.
\end{remark}


\subsection{Multiple Assets}

Real markets contain many stocks, each driven by potentially overlapping sources of randomness. The multi-asset generalization introduces the question of whether a risk-neutral measure exists and, if so, whether it is unique.

\begin{definition}
Under the \textbf{multidimensional market model}, consider $m$ stocks driven by $d$ independent Brownian motions $W_1(t), \ldots, W_d(t)$. Propose that stock $i$ has dynamics
\begin{align}
    dS_i(t) = \mu_i(t)S_i(t)\,dt + S_i(t)\sum_{j=1}^d \sigma_{ij}(t)\,dW_j(t),
\end{align}
where $\sigma_{ij}(t)$ is the loading of stock $i$ on the $j$-th Brownian motion.
\end{definition}

The total volatility and the instantaneous correlations between stocks then emerge naturally from this representation.

\begin{proposition}
Define the aggregate volatility of stock $i$ and the correlated Brownian motions $B_i(t)$ as
\begin{align}
    \sigma_i(t) = \sqrt{\sum_{j=1}^d \sigma_{ij}(t)^2}, \qquad dB_i(t) = \sum_{j=1}^d \frac{\sigma_{ij}(t)}{\sigma_i(t)}\,dW_j(t).
\end{align}
Then $B_i(t)$ is a Brownian motion by L\'evy's theorem, and
\begin{align}
    dS_i(t) &= \mu_i(t)S_i(t)\,dt + \sigma_i(t)S_i(t)\,dB_i(t), \\
    dB_i(t)\,dB_k(t) &= \rho_{ik}(t)\,dt, \qquad \rho_{ik}(t) = \sum_{j=1}^d \frac{\sigma_{ij}(t)\sigma_{kj}(t)}{\sigma_i(t)\sigma_k(t)}.
\end{align}
\end{proposition}

\begin{remark}
We have ''rotated'' from stocks being driven by multiple uncorrelated Brownian motions to being driven by individual, correlated Brownian motions. The correlation $\rho_{ik}$ arises because stocks $i$ and $k$ share exposure to the same underlying factors $W_j$. This is the structure used in portfolio theory: systematic risk (shared factor exposure) generates correlation, while idiosyncratic risk (exposure unique to one stock) does not.
\end{remark}

The question of whether a risk-neutral measure exists therefore reduces to a system of linear equations.

\begin{proposition}
Discounted stock prices under the multidimensional model can be made martingales by a Girsanov change of measure if there exists a $d$-dimensional adapted process $\Theta(t) = (\Theta_1(t), \ldots, \Theta_d(t))$ satisfying the $m$ equations
\begin{align}
    \mu_i(t) - R(t) = \sum_{j=1}^d \sigma_{ij}(t)\,\Theta_j(t), \qquad i = 1, \ldots, m.
\end{align}
If such a $\Theta$ exists, the discounted value of any self-financing portfolio is a martingale under the associated $\widetilde{\mathbb{P}}$.
\end{proposition}

\begin{proof}
Via It\^o's lemma, $d(D(t)S_i(t)) = D(t)S_i(t)\bigl[(\mu_i - R)\,dt + \sum_j \sigma_{ij}\,dW_j\bigr]$. For Girsanov to produce a martingale, this must take the form $D(t)S_i(t)\sum_j \sigma_{ij}(\Theta_j\,dt + dW_j)$. Matching the $dt$ terms gives the stated linear system. If a solution $\Theta$ exists, define $\widetilde{W}_j(t) = W_j(t) + \int_0^t \Theta_j(u)\,du$; under the associated $\widetilde{\mathbb{P}}$, the discounted prices have no drift.
\end{proof}

\begin{remark}
The system $\boldsymbol{\sigma}\,\Theta = \boldsymbol{\mu} - R\mathbf{1}$ is $m$ equations in $d$ unknowns. Three cases arise:
\begin{itemize}
    \item \emph{$m = d$} (as many stocks as Brownian motions): the system generically has a unique solution. There is one risk-neutral measure, and the market is complete.
    
    \item \emph{$m < d$} (fewer stocks than risk factors): the system is underdetermined. Multiple solutions exist, so the risk-neutral measure is not unique. The market is incomplete, so there are risks that cannot be hedged. This is the generic situation in practice (stochastic volatility, jumps, untraded factors).

    \item \emph{$m > d$} (more stocks than risk factors): the system is overdetermined and generically inconsistent. If the system has no solution, there is no risk-neutral measure, and arbitrage opportunities exist. In equilibrium, market forces eliminate such opportunities, which is why redundant assets must be priced consistently.
\end{itemize}
\end{remark}

\begin{example}[Two stocks, one factor]
Suppose $d = 1$ (one Brownian motion) and $m = 2$ stocks with $dS_i = \mu_i S_i\,dt + \sigma_i S_i\,dW$. The market-price-of-risk equations become $\mu_1 - R = \sigma_1 \Theta$ and $\mu_2 - R = \sigma_2 \Theta$. Dividing: $(\mu_1 - R)/\sigma_1 = (\mu_2 - R)/\sigma_2$. This is the single-factor CAPM relation: stocks with higher volatility must have proportionally higher excess returns. If this equality is violated, an arbitrage exists: go long the stock with the higher Sharpe ratio, short the other, and earn riskless profit.
\end{example}

We now formalize the concept of arbitrage and state the two fundamental theorems.

\begin{definition}
An \textbf{arbitrage} is a self-financing portfolio $X(t)$ such that
\begin{align}
    X(0) = 0, \qquad \mathbb{P}(X(T) \geq 0) = 1, \qquad \mathbb{P}(X(T) > 0) > 0.
\end{align}
In words: zero initial cost, no possibility of loss, positive probability of gain.
\end{definition}

\begin{theorem}[First fundamental theorem of asset pricing]
If a market model admits a risk-neutral measure $\widetilde{\mathbb{P}}$, then it does not admit arbitrage opportunities.
\end{theorem}

\begin{proof}
Suppose $\widetilde{\mathbb{P}}$ exists and $X(t)$ is a self-financing portfolio with $X(0) = 0$. Then $D(t)X(t)$ is a $\widetilde{\mathbb{P}}$-martingale, so $\widetilde{\mathbb{E}}[D(T)X(T)] = D(0)X(0) = 0$. Since $D(T) > 0$, if $\mathbb{P}(X(T) \geq 0) = 1$, then $\widetilde{\mathbb{P}}(X(T) \geq 0) = 1$ (the measures are equivalent). But a non-negative random variable with zero expectation must be zero a.s. Hence $\widetilde{\mathbb{P}}(X(T) > 0) = 0$, which implies $\mathbb{P}(X(T) > 0) = 0$. The portfolio cannot be an arbitrage.
\end{proof}

\begin{remark}
The converse, that absence of arbitrage implies existence of a risk-neutral measure, is also true but is substantially harder to prove. It requires functional-analytic arguments (the Hahn-Banach theorem or variants) and is stated precisely in the Delbaen-Schachermayer fundamental theorem for general semimartingale models. Together, the two directions establish: \emph{no arbitrage $\Longleftrightarrow$ existence of an equivalent martingale measure}.
\end{remark}

\begin{theorem}[Second fundamental theorem of asset pricing]
For a market admitting a risk-neutral measure, the market is complete if and only if the risk-neutral measure is unique.
\end{theorem}

\begin{remark}
Completeness means every contingent claim can be replicated. Uniqueness of $\widetilde{\mathbb{P}}$ means every claim has an unambiguous price. The equivalence is intuitive: if a claim can be perfectly hedged, its price is determined by the cost of the hedge, leaving no room for a second martingale measure to assign a different value. Conversely, if the measure is not unique, some payoffs receive a range of arbitrage-free prices, corresponding to risks that cannot be hedged.

Incomplete markets (stochastic volatility, jumps, illiquid assets) are the norm in practice. Pricing in incomplete markets requires additional criteria, such as utility maximization, mean-variance hedging, or market convention, to select among the continuum of equivalent martingale measures.
\end{remark}

\section{Change of Num\'eraire}

\subsection{Motivation and Definitions}

Risk-neutral pricing uses the money market account as the ``unit of account'': all prices are discounted by $D(t) = 1/M(t)$. But there is nothing special about the money market. Any strictly positive traded asset can serve as the reference, and switching to a different reference (a different \emph{num\'eraire}) produces a different (but equivalent) martingale measure under which pricing may be simpler.

\begin{definition}
A \textbf{num\'eraire} is a strictly positive adapted process $N(t) > 0$ a.s.\ for all $t \geq 0$, where $N$ is a traded asset or a self-financing portfolio. The \textbf{relative price} of an asset $S(t)$ with respect to $N$ is
\begin{align}
    S^{(N)}(t) = \frac{S(t)}{N(t)}.
\end{align}
\end{definition}

\begin{definition}
Let $M(t) = \exp\bigl(\int_0^t R(u)\,du\bigr)$ be the money market account. A probability measure $\mathbb{Q}$ is a \textbf{martingale measure associated with num\'eraire $N$} if every traded asset $S(t)$, when expressed in units of $N$, is a $\mathbb{Q}$-martingale:
\begin{align}
    \frac{S(t)}{N(t)} = \mathbb{E}^{\mathbb{Q}}\!\left[\frac{S(T)}{N(T)}\;\middle|\;\mathcal{F}(t)\right].
\end{align}
Taking $N = M$ recovers the standard risk-neutral measure $\widetilde{\mathbb{P}}$.
\end{definition}

\begin{remark}
The idea is simple: in any consistent pricing system, \emph{relative} prices (price of one asset in units of another) must be martingales under some equivalent measure. Which measure depends on which asset you choose as the measuring stick.
\end{remark}


\subsection{The Change of Num\'eraire Theorem}

We now show how to construct the new measure explicitly, using Girsanov's theorem.

\begin{theorem}[Stochastic representation of traded assets]
Consider the multidimensional market model with $m$ assets, and let $\widetilde{\mathbb{P}}$ be the risk-neutral measure with respect to the money market num\'eraire, with $d$-dimensional Brownian motion $\widetilde{W}(t)$. Let $N(t)$ be a positive, non-dividend-paying traded asset. Then there exists a vector volatility process $\nu(t) = (\nu_1(t), \ldots, \nu_d(t))$ such that
\begin{align}
    dN(t) = R(t)N(t)\,dt + N(t)\,\nu(t) \cdot d\widetilde{W}(t),
\end{align}
or equivalently,
\begin{align}
    D(t)N(t) = N(0)\exp\!\Bigl(-\frac{1}{2}\int_0^t |\nu(u)|^2\,du + \int_0^t \nu(u) \cdot d\widetilde{W}(u)\Bigr).
\end{align}
\end{theorem}

\begin{proof}
Under $\widetilde{\mathbb{P}}$, the discounted price $D(t)N(t)$ is a martingale. By the martingale representation theorem, there exists a $d$-dimensional adapted process $\widetilde{\Gamma}(t)$ with $d(D(t)N(t)) = \widetilde{\Gamma}(t) \cdot d\widetilde{W}(t)$. Defining $\nu_i(t) = \widetilde{\Gamma}_i(t)/(D(t)N(t))$ gives $d(D(t)N(t)) = D(t)N(t)\,\nu(t) \cdot d\widetilde{W}(t)$. Applying It\^o's formula to the exponential form verifies the representation.
\end{proof}

\begin{remark}
The vector $\nu(t)$ is the volatility structure of $N$ with respect to the underlying factor Brownian motions. For the money market account $M(t)$, which is locally deterministic, $\nu = 0$. For a zero-coupon bond $B(t,T)$, $\nu(t)$ encodes the bond's exposure to interest rate risk factors.
\end{remark}

The next theorem is the main result: it constructs the new martingale measure and shows that relative prices are martingales under it.

\begin{theorem}[Change of num\'eraire]
Let $S(t)$ and $N(t)$ be traded assets with discounted dynamics
\begin{align}
    d\bigl(D(t)S(t)\bigr) &= D(t)S(t)\,\sigma(t) \cdot d\widetilde{W}(t), \\
    d\bigl(D(t)N(t)\bigr) &= D(t)N(t)\,\nu(t) \cdot d\widetilde{W}(t).
\end{align}
Take $N(t)$ as the num\'eraire. Then there exists a probability measure $\widetilde{\mathbb{P}}^{(N)}$, equivalent to $\widetilde{\mathbb{P}}$, under which the relative price $S^{(N)}(t) = S(t)/N(t)$ is a martingale. The relative price has differential
\begin{align}
    dS^{(N)}(t) = S^{(N)}(t)\bigl[\sigma(t) - \nu(t)\bigr] \cdot d\widetilde{W}^{(N)}(t),
\end{align}
where $\widetilde{W}^{(N)}(t) = \widetilde{W}(t) - \int_0^t \nu(u)\,du$ is a Brownian motion under $\widetilde{\mathbb{P}}^{(N)}$.
\end{theorem}

\begin{proof}
We construct the measure $\widetilde{\mathbb{P}}^{(N)}$ by applying Girsanov's theorem a second time, starting from $\widetilde{\mathbb{P}}$. Set $\Theta_j(t) = -\nu_j(t)$, so the Radon-Nikodym derivative is
\begin{align}
    Z(t) = \exp\!\Bigl(\int_0^t \nu(u) \cdot d\widetilde{W}(u) - \frac{1}{2}\int_0^t |\nu(u)|^2\,du\Bigr) = \frac{D(t)N(t)}{N(0)},
\end{align}
where the last equality follows from the exponential representation of $D(t)N(t)$. The new Brownian motion is $\widetilde{W}^{(N)}(t) = \widetilde{W}(t) - \int_0^t \nu(u)\,du$.

For the relative price, divide the exponential forms of $D(t)S(t)$ and $D(t)N(t)$:
\begin{align}
    S^{(N)}(t) = \frac{S(0)}{N(0)}\exp\!\Bigl(-\frac{1}{2}\int_0^t\bigl(|\sigma|^2 - |\nu|^2\bigr)\,du + \int_0^t(\sigma - \nu) \cdot d\widetilde{W}\Bigr).
\end{align}
Applying It\^o's formula to this exponential gives
\begin{align}
    dS^{(N)}(t) = S^{(N)}(t)\,(\sigma - \nu) \cdot \bigl(-\nu\,dt + d\widetilde{W}\bigr) = S^{(N)}(t)\,(\sigma - \nu) \cdot d\widetilde{W}^{(N)}(t),
\end{align}
which has no $dt$ term, confirming the martingale property.
\end{proof}

\begin{remark}
The Radon-Nikodym derivative $Z(t) = D(t)N(t)/N(0)$ has an elegant interpretation: the density of the new measure is proportional to the discounted value of the new num\'eraire. This makes economic sense, as the new measure ``reweights'' states of the world by how valuable the num\'eraire is in those states. States where $N$ is large receive more weight, reflecting the fact that those states matter more when $N$ is your unit of account.
\end{remark}

The general pricing formula under any num\'eraire follows immediately.

\begin{theorem}[General num\'eraire pricing]
Let $V(T)$ be a contingent claim payable at time $T$. Its price at time $t$, expressed in units of the num\'eraire $N$, is
\begin{align}
    \frac{V(t)}{N(t)} = \mathbb{E}^{(\mathbb{Q}^N)}\!\left[\frac{V(T)}{N(T)}\;\middle|\;\mathcal{F}(t)\right].
\end{align}
Equivalently, $V(t) = N(t)\,\mathbb{E}^{(\mathbb{Q}^N)}[V(T)/N(T) \mid \mathcal{F}(t)]$.
\end{theorem}

\begin{remark}
For $N = M$ (money market), this reduces to the standard risk-neutral formula with discounting by $e^{-\int r\,du}$. For $N = B(\cdot, T)$ (zero-coupon bond), this produces the forward measure $\mathbb{Q}^T$. For $N = S$ (the stock itself), this produces the ``stock measure'' used in some exotic option pricing contexts. The choice of num\'eraire is dictated by convenience: pick whichever $N$ makes the expectation easiest to evaluate.
\end{remark}


\subsection{The Forward Measure}

The most important application of change of num\'eraire in practice is the \textbf{forward measure}, obtained by taking a zero-coupon bond as the num\'eraire.

\begin{example}[Forward contracts under the forward measure]
Let $B(t,T)$ be the time-$t$ price of a zero-coupon bond maturing at $T$. Consider a forward contract on an asset $S$, with delivery at $T$ and strike $K$, whose payoff is $S(T) - K$.

Under the money market num\'eraire, the price at time $t$ is
\begin{align}
    V(t) = \widetilde{\mathbb{E}}\!\left[e^{-\int_t^T r(u)\,du}\bigl(S(T) - K\bigr)\;\middle|\;\mathcal{F}(t)\right].
\end{align}
This is awkward: the discount factor $e^{-\int r\,du}$ and the payoff $S(T) - K$ are both random and potentially correlated (if interest rates are stochastic), so the expectation does not factor.

Now switch to the num\'eraire $N(t) = B(t,T)$. Under the associated forward measure $\mathbb{Q}^T$, the relative price $S(t)/B(t,T)$ is a martingale, so
\begin{align}
    \frac{V(t)}{B(t,T)} = \mathbb{E}^{(\mathbb{Q}^T)}\!\left[\frac{S(T) - K}{B(T,T)}\;\middle|\;\mathcal{F}(t)\right] = \mathbb{E}^{(\mathbb{Q}^T)}\!\left[S(T) - K\;\middle|\;\mathcal{F}(t)\right],
\end{align}
since $B(T,T) = 1$. By the martingale property of $S/B$ under $\mathbb{Q}^T$:
\begin{align}
    \mathbb{E}^{(\mathbb{Q}^T)}[S(T) \mid \mathcal{F}(t)] = \frac{S(t)}{B(t,T)}.
\end{align}
Therefore $V(t) = B(t,T)\bigl(S(t)/B(t,T) - K\bigr) = S(t) - KB(t,T)$, and the forward price that makes the contract worth zero at inception is
\begin{align}
    F(t,T) = \frac{S(t)}{B(t,T)}.
\end{align}
\end{example}

\begin{remark}
Under the forward measure, the forward price $F(t,T) = S(t)/B(t,T)$ is a martingale. This is enormously useful. Interest rate derivatives such as caplets and swaptions have payoffs that depend on forward rates. Under the forward measure, these forward rates are driftless, allowing direct application of the Black formula without worrying about the drift induced by stochastic discounting. This is the foundation of the LIBOR Market Model discussed in the fixed-income chapter.
\end{remark}

\begin{example}[Numerical illustration]
Consider a 1-year forward on a stock with $S(0) = 100$. In a constant-rate world with $r = 5\%$, the bond price is $B(0,1) = e^{-0.05} = 0.9512$, and the forward price is $F(0,1) = 100/0.9512 = 105.13$.

Now suppose rates are stochastic and the current bond price is $B(0,1) = 0.9480$ (corresponding to a yield of $5.34\%$). The forward price becomes $F(0,1) = 100/0.9480 = 105.49$. The higher yield increases the forward price because the cost of carrying the stock for one year is higher.

The key point is that the forward price is simply $S(0)/B(0,1)$ regardless of the interest rate model. The change of num\'eraire has fortunately eliminated the need to specify the dynamics of $r(t)$.
\end{example}


\chapter{Asset Pricing and Derivatives}

The risk-neutral framework of the preceding chapter established a universal pricing formula: the value of any derivative is the discounted expected payoff under the risk-neutral measure. This chapter applies that formula to concrete financial instruments, progressing from the simplest (forwards and futures) through the canonical (European options and Black-Scholes) to the structurally richer (exotic options, American options) and finally to volatility models that address the empirical failure of constant-volatility assumptions.

Each new instrument introduces a new mathematical feature. Forwards and futures test the machinery on linear payoffs. European options require Gaussian integration and produce the Black-Scholes formula. Asian options break the Markov property and require state-space enlargement. Barrier options invoke the reflection principle. American options introduce optimal stopping and free boundary problems. Stochastic volatility models add a second state variable and demand Fourier methods for tractable pricing. By the end of this chapter, the reader will have a concrete sense of both the power and the limitations of the no-arbitrage framework.


\section{Forwards and Futures}

We begin with the simplest derivatives: agreements to buy or sell an asset at a future date. Their pricing requires little more than the martingale property of discounted prices, yet the results are fundamental: the forward price appears throughout option pricing, and the distinction between forwards and futures illuminates the role of interest rate randomness.

\subsection{Cash Flows and Bond Prices}

\begin{proposition}[Pricing with cash flows]
For an asset generating cumulative cash flows $C(u)$, the value of a portfolio holding this asset is
\begin{align}
    X(t) = \frac{1}{D(t)}\,\widetilde{\mathbb{E}}\!\left[\int_t^T D(u)\,dC(u)\;\middle|\;\mathcal{F}(t)\right].
\end{align}
\end{proposition}

\begin{proof}[Proof sketch]
The discounted portfolio has differential $d(D(t)X(t)) = D(t)\,dC(t)$. Integrating from $t$ to $T$ and taking the risk-neutral conditional expectation, the left-hand side is a martingale, so $D(t)X(t) = \widetilde{\mathbb{E}}[D(T)X(T) - \int_t^T D(u)\,dC(u) \mid \mathcal{F}(t)]$ reduces to the stated formula.
\end{proof}

\begin{remark}
This is the continuous-time analogue of the discounted cash flow (DCF) formula from corporate finance: the value of an asset equals the present value of its future cash flows. The difference is that ``present value'' is computed under the risk-neutral measure, not using a risk-adjusted discount rate.
\end{remark}

The most basic cash-flow instrument is a zero-coupon bond, which pays exactly 1 at maturity.

\begin{proposition}
The price at time $t$ of a \textbf{zero-coupon bond} paying 1 at time $T$ is
\begin{align}
    B(t,T) = \frac{1}{D(t)}\,\widetilde{\mathbb{E}}\bigl[D(T)\;\big|\;\mathcal{F}(t)\bigr] = \widetilde{\mathbb{E}}\!\left[\exp\!\Bigl(-\int_t^T R(u)\,du\Bigr)\;\middle|\;\mathcal{F}(t)\right].
\end{align}
\end{proposition}

\begin{definition}
The \textbf{yield to maturity (YTM)} of a zero-coupon bond is the constant rate $Y(t,T)$ satisfying
\begin{align}
    B(t,T) = e^{-Y(t,T)(T-t)}.
\end{align}
The mapping $T \mapsto Y(t,T)$ is the \textbf{yield curve} at time $t$.
\end{definition}

\begin{remark}
When rates are deterministic, $B(t,T) = e^{-\int_t^T R(u)\,du}$ and the yield is the average of the short rate over $[t,T]$. When rates are stochastic, the yield reflects both the expected path of rates and a convexity adjustment (the bond price is a convex function of the rate, so Jensen's inequality creates a wedge between the yield and the expected average rate). This convexity effect is small for short maturities but significant for long-dated bonds.
\end{remark}

\subsection{Forward Contracts}

\begin{definition}
A \textbf{forward contract} on asset $S(t)$ with delivery time $T$ is an agreement in which the holder pays a fixed amount $K$ at time $T$ to receive the asset. The contract requires no upfront payment. The \textbf{forward price} $\text{For}_S(t,T)$ is the value of $K$ that makes the contract worth zero at time $t$.
\end{definition}

\begin{theorem}
Assuming zero-coupon bonds of all maturities are available, the forward price is
\begin{align}
    \mathrm{For}_S(t,T) = \frac{S(t)}{B(t,T)}.
\end{align}
\end{theorem}

\begin{proof}
The payoff at $T$ is $S(T) - K$. Under risk-neutral pricing:
\begin{align}
    V(t) = \frac{1}{D(t)}\,\widetilde{\mathbb{E}}\bigl[D(T)(S(T)-K)\;\big|\;\mathcal{F}(t)\bigr] = S(t) - K\,B(t,T),
\end{align}
using linearity and the martingale property of $D(t)S(t)$. Setting $V(t) = 0$ and solving for $K$ gives $K = S(t)/B(t,T)$.
\end{proof}

\begin{example}[Forward pricing with dividends]
Consider a stock at $S(0) = 100$ with a continuous dividend yield $q = 2\%$, risk-free rate $r = 5\%$, and $T = 1$ year. The dividend-adjusted forward price is
\begin{align}
    \mathrm{For}_S(0,1) = S(0)\,e^{(r-q)T} = 100 \cdot e^{0.03} \approx 103.05.
\end{align}
The forward is below the na\"ive ``growth at $r$'' level of $105.13$ because the dividend yield reduces the cost of carry: a holder of the forward misses the dividends that a holder of the spot collects.

This relationship can be exploited as an arbitrage check. If the observed forward trades at $104$, an arbitrageur borrows \$100 at $r = 5\%$, buys the stock, earns dividends at $q = 2\%$, and delivers at $T$ for \$104, netting $104 - 100 e^{0.05} + 100(e^{0.02}-1) \approx 0.95$ in profit.
\end{example}


\subsection{Futures Contracts}

\begin{definition}
A \textbf{futures contract} on asset $S(t)$ with delivery time $T$ is an arrangement in which:
\begin{enumerate}
    \item No payment is required at initiation.
    \item The contract is \textbf{marked to market}: holding a long position over $[s,t]$ generates a cash flow of $\text{Fut}_S(t,T) - \text{Fut}_S(s,T)$ via a margin account.
    \item At delivery, the holder pays $\text{Fut}_S(T,T)$ to receive the asset, so $\text{Fut}_S(T,T) = S(T)$.
\end{enumerate}
\end{definition}

\begin{theorem}
The futures price is
\begin{align}
    \mathrm{Fut}_S(t,T) = \widetilde{\mathbb{E}}\bigl[S(T)\;\big|\;\mathcal{F}(t)\bigr].
\end{align}
The futures price is a martingale under $\widetilde{\mathbb{P}}$, and the value of holding a futures contract over any period is zero.
\end{theorem}

\begin{remark}
When interest rates are deterministic, futures and forward prices coincide: $\text{Fut}_S(t,T) = \text{For}_S(t,T)$. When rates are stochastic, they differ. The difference arises from the daily mark-to-market: gains on the futures position are reinvested at the prevailing rate, and losses must be financed at that rate. If the asset price is positively correlated with interest rates (as for bonds), the futures price exceeds the forward price because gains are reinvested at high rates and losses financed at low rates. This \textbf{convexity bias} is small for short maturities but can matter for long-dated contracts (e.g., Eurodollar futures vs.\ FRA rates).
\end{remark}


\section{European Options}

European options are the canonical derivative: their pricing theory is the most developed, their formulas the most widely used, and their properties the most thoroughly understood. The Black-Scholes model, despite its well-known limitations, remains the lingua franca of option markets.

\begin{definition}
A \textbf{European call option} on $S(t)$ with strike $K$ and maturity $T$ gives the holder the right, but not the obligation, to buy the asset at time $T$ for price $K$. Its payoff is $\max(S(T)-K,0)$.

A \textbf{European put option} gives the right to sell, with payoff $\max(K-S(T),0)$.
\end{definition}


\subsection{Black-Scholes}

\begin{theorem}[Black-Scholes PDE]
Let $S(t)$ follow geometric Brownian motion with constant drift $\mu$ and volatility $\sigma$, and let $r$ be the constant risk-free rate. The price $c(t,x)$ of a European call satisfies the PDE
\begin{align}
    \frac{\partial c}{\partial t} + rx\,\frac{\partial c}{\partial x} + \frac{1}{2}\sigma^2 x^2\,\frac{\partial^2 c}{\partial x^2} = r\,c,
\end{align}
with terminal condition $c(T,x) = \max(x-K,0)$.
\end{theorem}

\begin{proof}
Construct a portfolio that is long one call and short $\Delta(t)$ shares: $\Pi(t) = c(t,S) - \Delta(t) S(t)$, with the remaining value $\Pi - c + \Delta S$ invested at rate $r$. By It\^o's lemma on $c(t,S)$ with $dS = \mu S\,dt + \sigma S\,dW$, matching the discounted portfolio and discounted call differentials yields two equations from the $dt$ and $dW$ coefficients:
\begin{align}
    \Delta(t)\,\sigma S\,dW &= \sigma S\,\frac{\partial c}{\partial x}\,dW, \\
    \Delta(t)\,(\mu-r)S\,dt &= \Bigl[\frac{\partial c}{\partial t} + \mu S\frac{\partial c}{\partial x} + \frac{1}{2}\sigma^2 S^2\frac{\partial^2 c}{\partial x^2} - rc\Bigr]dt.
\end{align}
The first equation determines the hedge: $\Delta(t) = \partial c/\partial x$. Substituting into the second eliminates $\mu$ and produces the Black-Scholes PDE.
\end{proof}

\begin{remark}
The most striking feature of the Black-Scholes PDE is what is \emph{absent}: the drift $\mu$. The expected return on the stock does not appear in the pricing equation. This is because the hedge eliminates exposure to the stock's direction. An option seller who delta-hedges continuously is indifferent to whether the stock drifts up or down; she cares only about how much it fluctuates. This is the practical content of risk-neutral pricing: the hedge converts the option into a volatility bet, stripping away directional risk.
\end{remark}


\begin{theorem}[Black-Scholes formula]
The solution to the Black-Scholes PDE is
\begin{align}
    c(t,x) &= x\,\Phi(d_+) - K e^{-r\tau}\,\Phi(d_-), \\
    d_{\pm} &= \frac{\log(x/K) + (r \pm \frac{1}{2}\sigma^2)\tau}{\sigma\sqrt{\tau}},
\end{align}
where $\tau = T - t$ and $\Phi$ is the standard normal CDF.
\end{theorem}

\begin{proof}[Proof sketch]
By the Feynman-Kac theorem, the solution to the PDE is
\begin{align}
    c(t,x) = \widetilde{\mathbb{E}}\bigl[e^{-r\tau}\max(S(T)-K,0)\;\big|\;S(t)=x\bigr].
\end{align}
Under $\widetilde{\mathbb{P}}$, $\log S(T) = \log x + (r - \frac{1}{2}\sigma^2)\tau + \sigma\sqrt{\tau}\,Z$ with $Z \sim \mathcal{N}(0,1)$. The option finishes in the money when $Z > -d_-$. Splitting the expectation into the stock-delivery and strike-payment components and evaluating each Gaussian integral produces the formula.
\end{proof}

\begin{remark}
Each term has a financial interpretation. The first term $x\,\Phi(d_+)$ is the present value of receiving the stock, weighted by the risk-neutral probability (under a stock-num\'eraire measure) that $S(T) > K$. The second term $Ke^{-r\tau}\Phi(d_-)$ is the present value of paying the strike, weighted by the risk-neutral probability that the option is exercised. The quantity $\Phi(d_-)$ is sometimes called the ``exercise probability'' and $\Phi(d_+)$ is the ``delta.''
\end{remark}

\begin{example}[Black-Scholes computation]
Consider a call on a stock with $S = 100$, $K = 105$, $\sigma = 25\%$, $r = 3\%$, $T - t = 0.5$ years. Then:
\begin{align}
    d_- &= \frac{\log(100/105) + (0.03 - 0.03125) \times 0.5}{0.25\sqrt{0.5}} = \frac{-0.04879 - 0.000625}{0.17678} = -0.2797, \\
    d_+ &= -0.2797 + 0.17678 = -0.1029.
\end{align}
So $\Phi(d_+) = 0.4590$, $\Phi(d_-) = 0.3899$, and
\begin{align}
    c = 100 \times 0.4590 - 105 \times e^{-0.015} \times 0.3899 = 45.90 - 40.32 = 5.58.
\end{align}
The call is worth \$5.58 despite being out of the money ($S < K$), reflecting the value of optionality over the remaining 6 months.
\end{example}


\subsection{Greeks and Put-Call Parity}

\begin{definition}
The \textbf{Greeks} are partial derivatives of the option price with respect to model inputs. The most important are:
\begin{align}
    \Delta &= \frac{\partial c}{\partial x} = \Phi(d_+), &\quad &\text{(sensitivity to spot price)} \\
    \Gamma &= \frac{\partial^2 c}{\partial x^2} = \frac{\phi(d_+)}{x\sigma\sqrt{\tau}}, &\quad &\text{(curvature / hedging cost)} \\
    \Theta &= \frac{\partial c}{\partial t}, &\quad &\text{(time decay)} \\
    \mathcal{V} &= \frac{\partial c}{\partial \sigma} = x\phi(d_+)\sqrt{\tau}, &\quad &\text{(sensitivity to volatility)}
\end{align}
where $\phi$ is the standard normal density.
\end{definition}

\begin{remark}
The Greeks are tools of practical risk management. A market maker holding a portfolio of options monitors its aggregate delta (directional exposure), gamma (exposure to large moves), vega (exposure to implied volatility changes), and theta (daily time decay). The Black-Scholes PDE itself can be rewritten as a relationship among Greeks:
\begin{align}
    \Theta + rS\Delta + \frac{1}{2}\sigma^2 S^2 \Gamma = rc.
\end{align}
This says that the portfolio's time decay $\Theta$ plus its interest-rate-adjusted drift $rS\Delta$ plus its gamma profit $\frac{1}{2}\sigma^2 S^2 \Gamma$ must equal the financing cost $rc$. If you are long gamma ($\Gamma > 0$), you profit from realized volatility but pay for it through theta decay.
\end{remark}

\begin{theorem}[Put-call parity]
For European call and put options with the same strike and maturity on a non-dividend-paying asset:
\begin{align}
    c(t,S) - p(t,S) = S(t) - K e^{-r(T-t)}.
\end{align}
\end{theorem}

\begin{proof}
At time $T$, the difference $c(T,S) - p(T,S) = \max(S-K,0) - \max(K-S,0) = S(T) - K$. This equals the payoff of a forward with delivery price $K$. By no-arbitrage, this relationship must hold at all times: $c - p = S - KB(t,T) = S - Ke^{-r(T-t)}$.
\end{proof}

\begin{remark}
Put-call parity is model-free: it requires only no-arbitrage and European exercise, not any assumption about the dynamics of $S$. It is therefore one of the most robust relationships in finance. Violations of put-call parity indicate either arbitrage opportunities, transaction costs, or deviations from the European exercise assumption (e.g., early exercise features in American options).

In practice, traders use put-call parity to convert between call and put implied volatilities, to check for data errors, and to synthesize forwards from options.
\end{remark}


\section{Exotic Options}

European options depend only on the terminal value $S(T)$. Many traded derivatives, however, depend on the \emph{path} of $S$ over $[0,T]$: its running average, its maximum, whether it crossed a barrier. These \textbf{exotic options} introduce new mathematical challenges (the breakdown of the simple Markov structure in $(t, S(t))$) and require the state-space enlargement and generator techniques developed in the mathematical foundations chapter.

\subsection{Asian Options}

\begin{definition}
An \textbf{Asian option} (or \textbf{average-price option}) has a payoff that depends on the time-average of the stock price. A fixed-strike Asian call with continuous sampling pays
\begin{align}
    V(T) = \max\!\Bigl(\frac{1}{T}\int_0^T S(u)\,du - K,\;0\Bigr).
\end{align}
\end{definition}

\begin{remark}
Asian options are widely used in commodity and FX markets, where averaging reduces the impact of price manipulation or extreme spot moves near expiry. They are cheaper than European options because averaging reduces the effective volatility of the payoff: $\text{Var}(\bar{S}) < \text{Var}(S(T))$ due to the correlation structure of the average.
\end{remark}

The challenge is that the payoff depends on $A(T) = \int_0^T S(u)\,du$, which is not a function of $S(T)$ alone. The price cannot be written as $c(t, S(t))$, sinceit must depend on the accumulated average as well.

\begin{remark}[State-space enlargement]
The process $S(t)$ is Markov, but the payoff functional $A(T)$ introduces path dependence. The remedy is to enlarge the state space by defining a joint process $Y(t) = (S(t), A(t))$ where $A(t) = \int_0^t S(u)\,du$. The joint process is Markov: given $(S(t), A(t))$, the future evolution of $A$ depends only on the future of $S$, which depends only on $S(t)$. This technique of making the ``memory'' explicit as an additional state variable recurs throughout exotic option pricing.
\end{remark}

\begin{theorem}
Define the Asian option price as
\begin{align}
    u(t,s,a) = \widetilde{\mathbb{E}}\!\left[e^{-r(T-t)}\,h\!\Bigl(\frac{A(T)}{T}\Bigr)\;\middle|\;S(t) = s,\;A(t) = a\right].
\end{align}
Then $u$ satisfies the PDE
\begin{align}
    \frac{\partial u}{\partial t} + rs\,\frac{\partial u}{\partial s} + s\,\frac{\partial u}{\partial a} + \frac{1}{2}\sigma^2 s^2\,\frac{\partial^2 u}{\partial s^2} - ru = 0,
\end{align}
with terminal condition $u(T,s,a) = h(a/T)$.
\end{theorem}

\begin{proof}
Under $\widetilde{\mathbb{P}}$, the dynamics are $dS = rS\,dt + \sigma S\,dW$ and $dA = S\,dt$. The joint process $(S,A)$ is Markov with generator
\begin{align}
    \mathcal{L}f(s,a) = rs\,\frac{\partial f}{\partial s} + s\,\frac{\partial f}{\partial a} + \frac{1}{2}\sigma^2 s^2\,\frac{\partial^2 f}{\partial s^2},
\end{align}
obtained by applying It\^o's formula to $f(S,A)$ and reading off the $dt$ coefficient. The term $s\,\partial f/\partial a$ arises from $dA = S\,dt$: the average accumulates deterministically given the current stock price, introducing a first-order transport term with no associated diffusion.

By the Feynman-Kac theorem, the price function $u$ satisfies $\partial_t u + \mathcal{L}u - ru = 0$.
\end{proof}

\begin{remark}
The Asian PDE is two-dimensional in space $(s,a)$ plus time, making finite-difference methods expensive. In practice, Asian options are typically priced by Monte Carlo simulation (straightforward, since the average is easily computed along each path) or by reducing the dimension via a change of variables that exploits the homogeneity of the payoff.
\end{remark}


\subsection{Barrier Options}

\begin{definition}
A \textbf{barrier option} is a European option whose existence depends on whether the underlying crosses a specified level $B$ during the option's life.
\begin{enumerate}
    \item A \textbf{knock-out} option becomes worthless if the barrier is hit.
    \item A \textbf{knock-in} option becomes active only if the barrier is hit.
\end{enumerate}
Combining with the position of the barrier relative to the initial price gives four types: up-and-out, up-and-in, down-and-out, down-and-in. By no-arbitrage, knock-in $+$ knock-out $=$ vanilla European (same strike and maturity).
\end{definition}

\begin{remark}
Barrier options are cheaper than vanilla options because they ``die'' or fail to activate under certain scenarios. An up-and-out call with a barrier modestly above the strike has a substantially lower premium than a vanilla call, because the scenarios in which the stock rises strongly kill the barrier option. This makes barrier options popular for structured products where cost reduction matters.
\end{remark}

Pricing barrier options requires the joint distribution of the terminal value and the running maximum (or minimum) of the stock price.

\begin{proposition}
Let $\hat{W}(t) = \alpha t + \widetilde{W}(t)$ be a Brownian motion with drift, and $\hat{M}(t) = \max_{0 \leq u \leq t} \hat{W}(u)$ its running maximum. The joint density under $\widetilde{\mathbb{P}}$ is
\begin{align}
    \widetilde{f}_{(\hat{M}(t),\,\hat{W}(t))}(m,w) = \frac{2(2m-w)}{t\sqrt{2\pi t}}\,\exp\!\Bigl(\alpha w - \frac{1}{2}\alpha^2 t - \frac{(2m-w)^2}{2t}\Bigr),
\end{align}
for $m \geq \max(0,w)$.
\end{proposition}

\begin{proof}[Proof idea]
Under zero drift ($\alpha = 0$), the joint density of $(M(t), W(t))$ follows from the reflection principle. For non-zero drift, apply Girsanov's theorem: the density under the drifted measure acquires a Radon-Nikodym factor $\exp(\alpha w - \frac{1}{2}\alpha^2 t)$.
\end{proof}

\begin{theorem}
An up-and-out call with barrier $B > S(0)$ and strike $K < B$ has price $v(t,S)$ satisfying the Black-Scholes PDE on the restricted domain $[0,T] \times [0,B]$:
\begin{align}
    \frac{\partial v}{\partial t} + rx\,\frac{\partial v}{\partial x} + \frac{1}{2}\sigma^2 x^2\,\frac{\partial^2 v}{\partial x^2} = rv,
\end{align}
with boundary conditions
\begin{align}
    v(t,0) = 0, \qquad v(t,B) = 0, \qquad v(T,x) = \max(x-K,0).
\end{align}
The price can also be computed as a risk-neutral expectation using the joint density above.
\end{theorem}

\begin{remark}
The condition $v(t,B) = 0$ is the mathematical encoding of ``worthless at the barrier.'' This transforms the Black-Scholes PDE from an initial-value problem on the half-line $[0,\infty)$ to a boundary-value problem on $[0,B]$. The solution can be obtained by the method of images (reflecting the vanilla solution across the barrier) or by direct integration against the joint density.
\end{remark}


\subsection{Lookback Options}

\begin{definition}
A \textbf{lookback option} has a payoff that depends on the extremum of the stock price over the option's life. A \textbf{floating-strike lookback call} pays
\begin{align}
    V(T) = S(T) - \min_{0 \leq u \leq T} S(u),
\end{align}
guaranteeing the holder the difference between the terminal price and the lowest price achieved. A \textbf{floating-strike lookback put} pays $\max_{0 \leq u \leq T}S(u) - S(T)$.
\end{definition}

The running maximum $Y(t) = \max_{0 \leq u \leq t} S(u)$ introduces path dependence. As with Asian options, the remedy is state-space enlargement: the price depends on $(t, S(t), Y(t))$.

\begin{theorem}
The price $v(t,x,y)$ of a floating-strike lookback put satisfies the Black-Scholes PDE in the region $\{x \leq y\}$:
\begin{align}
    \frac{\partial v}{\partial t} + rx\,\frac{\partial v}{\partial x} + \frac{1}{2}\sigma^2 x^2\,\frac{\partial^2 v}{\partial x^2} = rv,
\end{align}
with boundary and terminal conditions
\begin{align}
    v(t,0,y) = e^{-r(T-t)}y, \qquad \frac{\partial v}{\partial y}(t,y,y) = 0, \qquad v(T,x,y) = y - x.
\end{align}
\end{theorem}

\begin{remark}
The Neumann condition $\partial v/\partial y = 0$ at $x = y$ (the ``diagonal'') reflects the fact that the running maximum only changes when the stock hits a new high. At that instant, $S = Y$ and an infinitesimal increase in $Y$ has no first-order effect on the option value. Lookback options are expensive because they eliminate the risk of poor market timing; they are sometimes called ``perfect hindsight'' options.
\end{remark}


\section{American Options}

American options can be exercised at any time up to maturity, making their pricing fundamentally different from European options. The holder faces an \textbf{optimal stopping problem}: exercise too early and sacrifice the remaining time value; exercise too late and risk the option moving out of the money.

\begin{definition}
An \textbf{American option} with payoff $h(x)$ on asset $S(t)$ can be exercised at any stopping time $\tau \in [t,T]$. Its value is the solution to the optimal stopping problem:
\begin{align}
    v(t,x) = \sup_{\tau \in \mathcal{T}_{t,T}} \widetilde{\mathbb{E}}\!\left[e^{-r(\tau-t)}\,h(S(\tau))\;\middle|\;S(t) = x\right],
\end{align}
where $\mathcal{T}_{t,T}$ is the set of $\mathcal{F}$-stopping times in $[t,T]$.
\end{definition}

\begin{proposition}[Snell envelope characterization]
The American option value satisfies:
\begin{enumerate}
    \item $v(t,x) \geq h(x)$ for all $(t,x)$: the option is worth at least its intrinsic value, since immediate exercise is always available.
    \item The discounted value process $e^{-rt}v(t,S(t))$ is the smallest supermartingale dominating $e^{-rt}h(S(t))$.
\end{enumerate}
This characterization, known as the \textbf{Snell envelope}, provides the probabilistic foundation for the PDE formulation.
\end{proposition}

\begin{remark}
The supermartingale property is the key distinction from European pricing, where discounted values are martingales. An American option's discounted value is a supermartingale because the holder can capture value by exercising. The ``optional'' exercise feature creates a downward pressure on the expected future discounted value.
\end{remark}

\begin{theorem}[Variational inequality]
Let $\mathcal{L}$ be the generator of the risk-neutral stock dynamics,
\begin{align}
    \mathcal{L}f(x) = rx\,f'(x) + \frac{1}{2}\sigma^2 x^2\,f''(x).
\end{align}
The value function $v(t,x)$ satisfies the \textbf{free boundary problem}:
\begin{align}
    \max\!\Bigl(\frac{\partial v}{\partial t} + \mathcal{L}v - rv,\;\; h(x) - v\Bigr) = 0, \qquad v(T,x) = h(x).
\end{align}
The state space is partitioned into the \textbf{continuation region} $\mathcal{C} = \{(t,x) : v(t,x) > h(x)\}$, where $v$ satisfies the Black-Scholes PDE, and the \textbf{exercise region} $\mathcal{E} = \{(t,x) : v(t,x) = h(x)\}$, where immediate exercise is optimal. The boundary between them is an unknown, time-dependent \textbf{free boundary}.
\end{theorem}

\begin{remark}
The variational inequality encodes the following logic: at each point $(t,x)$, either the option is worth more than its exercise value (so the holder continues, and $v$ satisfies the PDE) or the option equals its exercise value (so the holder exercises, and the PDE inequality becomes strict). The ``max'' operator selects the binding constraint. This is a nonlinear problem that generally prohibits closed-form solutions: the exercise boundary must be determined simultaneously with the value function.
\end{remark}

\begin{theorem}[American put]
For an American put with payoff $h(x) = \max(K-x, 0)$, early exercise may be optimal. The value function satisfies the \textbf{linear complementarity} conditions:
\begin{align}
    &rv - v_t - rxv_x - \frac{1}{2}\sigma^2 x^2 v_{xx} \geq 0, \\
    &v(t,x) \geq K - x, \\
    &\bigl(rv - v_t - rxv_x - \tfrac{1}{2}\sigma^2 x^2 v_{xx}\bigr)\bigl(v - (K-x)\bigr) = 0.
\end{align}
There exists a free boundary $x = b(t)$ such that exercise is optimal for $x \leq b(t)$ and continuation for $x > b(t)$. At the boundary, the \textbf{value-matching} and \textbf{smooth-pasting} conditions hold:
\begin{align}
    v(t,b(t)) = K - b(t), \qquad \frac{\partial v}{\partial x}(t,b(t)) = -1.
\end{align}
\end{theorem}

\begin{remark}
Value-matching says the option value is continuous across the exercise boundary (no jump in value at the moment of exercise). Smooth-pasting says the delta is also continuous (no kink). Together, they pin down the free boundary. The exercise boundary $b(t)$ increases toward $K$ as $t \to T$ and satisfies $b(T) = K$ (at expiry, exercise when $S < K$). For $t$ well before expiry, $b(t) < K$: the holder tolerates some in-the-money value because the remaining time value exceeds the interest on the exercise proceeds.
\end{remark}

\begin{example}[Why early exercise of a put can be optimal]
Consider a put with $K = 100$, $r = 5\%$, $\sigma = 30\%$, $T = 1$. Suppose the stock drops to $S = 5$. The intrinsic value is $95$. If the holder exercises, she receives \$95 immediately and earns interest: $95 \times e^{0.05} \approx 99.87$ by expiry. If she waits, the maximum payoff at expiry is \$100 (if $S(T) = 0$), but the expected payoff is less than \$95 because the stock is unlikely to fall much further and might recover.

The interest earned on early exercise proceeds exceeds the option's remaining time value, making exercise optimal. This logic never applies to calls on non-dividend-paying stocks, because the ``insurance value'' of holding the call (limited downside) always exceeds the interest foregone.
\end{example}

\begin{theorem}[American call on a non-dividend-paying asset]
If the underlying pays no dividends, the American call price equals the European call price:
\begin{align}
    v(t,x) = c(t,x).
\end{align}
Early exercise is never optimal, and the continuation region is the entire state space.
\end{theorem}

\begin{remark}
Exercising a call early forfeits two sources of value: the time value of money on the strike payment (delaying payment of $K$ saves interest) and the insurance against the stock falling below $K$ (if you exercise and the stock drops, you lose; if you hold the call, your loss is capped at the premium). With dividends, the first effect is offset by the dividend income foregone, and early exercise can become optimal just before an ex-dividend date.
\end{remark}


\section{Local and Stochastic Volatility}

The Black-Scholes model assumes constant volatility, but market option prices imply a volatility that varies with strike and maturity (the \textbf{implied volatility surface}). This section develops three models that address this discrepancy: the local volatility model (which fits the surface exactly but with unrealistic dynamics), the Heston stochastic volatility model (which produces realistic dynamics but requires Fourier methods), and the SABR model (which dominates interest rate markets via an asymptotic formula).

\subsection{Dupire and Local Volatility}

\begin{definition}
A \textbf{local volatility model} specifies the risk-neutral stock dynamics as
\begin{align}
    dS(t) = r(t)S(t)\,dt + \sigma_{\text{loc}}(t, S(t))\,S(t)\,dW(t),
\end{align}
where $\sigma_{\text{loc}}(t,s)$ is a deterministic function of time and the current stock price, called the \textbf{local volatility}. The model is Markovian, time-inhomogeneous, and arbitrage-free under $\widetilde{\mathbb{P}}$.
\end{definition}

\begin{remark}
The question motivating local volatility is: \emph{Does there exist a one-factor diffusion whose model prices reproduce every observed European option price?} Dupire's formula answers this affirmatively, and provides an explicit expression for the volatility function.
\end{remark}

\begin{proposition}[Forward PDE for call prices]
Under the local volatility model, European call prices $c(t, T; K)$ satisfy the forward PDE
\begin{align}
    \frac{\partial c}{\partial T} = -r(T)K\,\frac{\partial c}{\partial K} + \frac{1}{2}\sigma_{\text{loc}}^2(T,K)\,K^2\,\frac{\partial^2 c}{\partial K^2}.
\end{align}
\end{proposition}

\begin{proof}[Proof idea]
Apply the forward Kolmogorov equation to the transition density of $S(t)$, then integrate against the call payoff in strike space.
\end{proof}

\begin{theorem}[Dupire's formula]
If European call prices $c(0,T; K)$ are observed for all strikes and maturities and are sufficiently smooth, the local volatility function is uniquely determined by
\begin{align}
    \sigma_{\text{loc}}^2(T,K) = \frac{\frac{\partial c}{\partial T} + r(T)K\,\frac{\partial c}{\partial K}}{\frac{1}{2}K^2\,\frac{\partial^2 c}{\partial K^2}}.
\end{align}
\end{theorem}

\begin{proof}
Rearrange the forward PDE.
\end{proof}

\begin{remark}
Dupire's formula is elegant but requires numerical differentiation of market prices. In practice, one first fits a smooth parametric model to the implied volatility surface (such as SVI), then applies Dupire's formula to the smoothed surface.

Local volatility models have an important structural limitation: they fit the \emph{current} implied volatility surface exactly, but their \emph{forward} implied volatilities (the model's prediction for the future smile) are unrealistic. Specifically, the forward smile flattens too quickly, producing poor prices for forward-starting options, cliquets, and other volatility-sensitive exotics.
\end{remark}

For parameterizing the implied volatility surface before applying Dupire, the SVI model is widely used.

\begin{definition}
The \textbf{SVI (stochastic volatility inspired)} parameterization models total implied variance $w(\tau,k) = \sigma_{\text{BS}}^2(\tau,k)\,\tau$ as a function of log-moneyness $k = \log(K/F)$ at each expiry $\tau$:
\begin{align}
    w(k) = a + b\bigl(\rho(k-m) + \sqrt{(k-m)^2 + \sigma^2}\bigr),
\end{align}
where $a, b, \rho, m, \sigma$ are parameters fit to market data. The surface is free of calendar arbitrage if $\partial w/\partial \tau \geq 0$ everywhere.
\end{definition}

\begin{remark}
The SVI parameterization has five parameters per slice and fits equity implied volatility surfaces well. The name ``stochastic volatility inspired'' reflects the fact that for large $|k|$, the total variance is asymptotically linear in $k$, consistent with the wing behavior predicted by stochastic volatility models. The constraint $\partial w/\partial \tau \geq 0$ (total variance must increase with maturity) enforces absence of calendar spread arbitrage.
\end{remark}


\subsection{The Heston Model}

\begin{definition}
The \textbf{Heston model} specifies the joint dynamics of the asset price $S(t)$ and its instantaneous variance $V(t)$ under the risk-neutral measure:
\begin{align}
    dS(t) &= rS(t)\,dt + \sqrt{V(t)}\,S(t)\,dW_1(t), \\
    dV(t) &= \kappa(\theta - V(t))\,dt + \xi\sqrt{V(t)}\,dW_2(t),
\end{align}
where $\kappa > 0$ is the rate of mean reversion, $\theta > 0$ is the long-run variance, $\xi > 0$ is the volatility of volatility, $W_1$ and $W_2$ are Brownian motions with $dW_1\,dW_2 = \rho\,dt$, and typically $\rho < 0$ (leverage effect).
\end{definition}

\begin{proposition}
The joint process $(S,V)$ is a two-dimensional Markov diffusion with generator
\begin{align}
    \mathcal{L}f = rs\,\frac{\partial f}{\partial s} + \kappa(\theta-v)\,\frac{\partial f}{\partial v} + \frac{1}{2}vs^2\,\frac{\partial^2 f}{\partial s^2} + \rho\xi vs\,\frac{\partial^2 f}{\partial s\,\partial v} + \frac{1}{2}\xi^2 v\,\frac{\partial^2 f}{\partial v^2}.
\end{align}
\end{proposition}

\begin{proof}[Proof idea]
Apply the multi-dimensional It\^o formula to a smooth test function $f(s,v)$; the generator is the coefficient of $dt$.
\end{proof}

By the Feynman-Kac theorem, the price of any European derivative under Heston satisfies a two-dimensional PDE.

\begin{theorem}
The price $u(t,s,v)$ of a European derivative with payoff $h(S(T))$ satisfies
\begin{align}
    \frac{\partial u}{\partial t} + \mathcal{L}u - ru = 0, \qquad u(T,s,v) = h(s).
\end{align}
\end{theorem}

The variance process must remain positive for the model to be well-defined. This is guaranteed by the Feller condition.

\begin{theorem}[Feller condition]
If the parameters satisfy $2\kappa\theta \geq \xi^2$, the variance process $V(t)$ remains strictly positive a.s. The boundary $V = 0$ is unattainable.
\end{theorem}

\begin{remark}
The Feller condition says that the ``pull'' of mean reversion ($\kappa\theta$) must dominate the ``push'' of volatility diffusion ($\xi^2/2$) near zero. In calibrated Heston models, the Feller condition is often violated, in which case $V$ can touch zero. The model remains well-defined (the CIR process has a unique strong solution regardless), but the boundary behavior requires careful numerical treatment.
\end{remark}

The Heston PDE does not have a closed-form solution in terms of the density of $S(T)$. However, the model is \emph{affine}: the characteristic function of $\log S(T)$ has an exponential-affine form, which enables efficient Fourier-based pricing.

\begin{proposition}[Fourier pricing framework]
Let $X(t) = \log S(t)$ and define the characteristic function
\begin{align}
    \varphi_{X(T)}(u) = \widetilde{\mathbb{E}}\bigl[e^{iuX(T)}\bigr].
\end{align}
If a payoff $\phi(x)$ admits a Fourier representation $\phi(x) = \frac{1}{2\pi}\int_{-\infty}^{\infty} e^{-ipx}\hat{\phi}(p)\,dp$, then the derivative price is
\begin{align}
    V(0) = \frac{e^{-rT}}{2\pi}\int_{-\infty}^{\infty} \varphi_{X(T)}(-p)\,\hat{\phi}(p)\,dp.
\end{align}
\end{proposition}

\begin{remark}
The key insight is that pricing reduces to evaluating the characteristic function, which has a known closed form for affine models, rather than the density, which does not. The Fourier integral can be computed efficiently by quadrature or the FFT.
\end{remark}

\begin{theorem}[Heston characteristic function]
The characteristic function of $X(T) = \log S(T)$ under the Heston model is
\begin{align}
    \varphi_{X(T)}(u) = \exp\bigl(C(T,u) + D(T,u)\,V(0) + iu\,X(0)\bigr),
\end{align}
where $C$ and $D$ solve the Riccati ODEs
\begin{align}
    \frac{dD}{dt} &= -\frac{1}{2}(u^2 + iu) + (\kappa - i\rho\xi u)\,D + \frac{1}{2}\xi^2 D^2, \\
    \frac{dC}{dt} &= \kappa\theta\,D,
\end{align}
with $C(0,u) = D(0,u) = 0$.
\end{theorem}

\begin{proof}[Proof idea]
Guess that the log-characteristic function is affine in $V(0)$ (motivated by the affine structure of the model), substitute into the backward Kolmogorov PDE $\partial_t \varphi + \mathcal{L}\varphi = 0$, and collect terms in powers of $v$ to obtain the Riccati ODEs.
\end{proof}

\begin{remark}
The Riccati ODE for $D$ has an explicit solution in terms of hyperbolic functions. The closed-form availability of $C$ and $D$ is what makes the Heston model practically usable: evaluating the characteristic function at a grid of frequencies is fast, and the FFT produces option prices across an entire strike range in a single pass.
\end{remark}

The European call payoff $\max(S(T)-K, 0)$ is not absolutely integrable, so its Fourier transform does not exist directly. The Carr-Madan approach resolves this.

\begin{example}[Carr-Madan pricing]
Define log-strike $k = \log K$ and introduce a damping factor $\alpha > 0$. The dampened call price $e^{\alpha k}C(K)$ is integrable, and its Fourier inversion gives
\begin{align}
    C(K) = \frac{e^{-\alpha k}}{\pi}\int_0^{\infty} \text{Re}\!\left[e^{-ivk}\,\frac{\varphi_{X(T)}(v - i(\alpha+1))}{(\alpha + iv)(\alpha + 1 + iv)}\right]dv.
\end{align}
The characteristic function $\varphi$ is evaluated at the complex argument $v - i(\alpha+1)$, which requires the analytic continuation of the Heston formula. The integral is computed by numerical quadrature, or by applying the FFT to the integrand sampled on a uniform grid in $v$, producing prices at a uniform grid in $k$ simultaneously.

A typical implementation evaluates $\varphi$ at $N = 2^{12} = 4096$ frequency points, applies the FFT (cost: $\mathcal{O}(N\log N)$), and produces call prices at 4096 strikes in under a millisecond. This efficiency makes the Heston model practical for real-time calibration to the entire implied volatility surface.
\end{example}


\subsection{The SABR Model}

While Heston is the workhorse for equity derivatives, the \textbf{SABR model} dominates interest rate markets. It models forward rates directly under the appropriate forward measure, and its appeal rests on an asymptotic implied volatility formula that is fast to evaluate and easy to calibrate.

\begin{definition}
Let $F(t)$ denote a forward price or forward interest rate. Under the relevant pricing measure, the \textbf{SABR model} is
\begin{align}
    dF(t) &= \alpha(t)\,F(t)^\beta\,dW_1(t), \\
    d\alpha(t) &= \nu\,\alpha(t)\,dW_2(t), \\
    d\langle W_1, W_2\rangle(t) &= \rho\,dt,
\end{align}
where $\alpha(t) > 0$ is the instantaneous volatility, $\beta \in [0,1]$ controls the elasticity, $\nu \geq 0$ is the vol-of-vol, and $\rho \in [-1,1]$ is the correlation.
\end{definition}

\begin{proposition}
The forward $F(t)$ is a martingale under the chosen measure (it has no drift term), preserving no-arbitrage.
\end{proposition}

\begin{remark}
The SABR model nests several classical models as special cases: $\beta = 1$ gives lognormal dynamics (Black's model), $\beta = 0$ gives normal dynamics (Bachelier's model), $\nu = 0$ gives the constant elasticity of variance (CEV) model, and $\rho = 0$ removes the leverage effect.

The implied volatility smile in SABR arises from three sources: stochastic volatility ($\nu > 0$) generates curvature, correlation ($\rho \neq 0$) generates skew, and the elasticity parameter ($\beta \neq 1$) controls the backbone (how the ATM level of implied volatility changes as the forward moves). Negative correlation ($\rho < 0$) produces the left skew commonly observed in rates markets: when rates fall, volatility rises.
\end{remark}

\begin{theorem}[Hagan et al.\ approximation]
No closed-form option prices exist for SABR. However, for $F(0), K > 0$ and small maturity $T$, the Black implied volatility admits the asymptotic expansion
\begin{align}
    \sigma_{\mathrm{imp}}(F_0,K) = \frac{z\,\alpha_0}{\chi(z)\,(F_0 K)^{(1-\beta)/2}}\left[1 + \left(\frac{\alpha_0^2(1-\beta)^2}{24(F_0 K)^{1-\beta}} + \frac{\rho\beta\nu\alpha_0}{4(F_0 K)^{(1-\beta)/2}} + \frac{2-3\rho^2}{24}\nu^2\right)T\right],
\end{align}
where
\begin{align}
    z = \frac{\nu}{\alpha_0}(F_0 K)^{(1-\beta)/2}\log\!\frac{F_0}{K}, \qquad \chi(z) = \log\!\left(\frac{\sqrt{1-2\rho z+z^2}+z-\rho}{1-\rho}\right).
\end{align}
\end{theorem}

\begin{corollary}
At-the-money ($F_0 = K$), the formula simplifies to
\begin{align}
    \sigma_{\mathrm{ATM}} = \frac{\alpha_0}{F_0^{1-\beta}}\left[1 + \left(\frac{(1-\beta)^2\alpha_0^2}{24\,F_0^{2-2\beta}} + \frac{\rho\beta\nu\alpha_0}{4\,F_0^{1-\beta}} + \frac{2-3\rho^2}{24}\nu^2\right)T\right].
\end{align}
\end{corollary}

\begin{remark}
This formula is used almost universally for interest rate option calibration. Its appeal is computational: evaluating $\sigma_{\text{imp}}$ at a given $(F_0, K, T)$ is instantaneous (no integration, no ODE solving), and calibrating the four parameters $(\alpha_0, \beta, \rho, \nu)$ to a set of market quotes is a straightforward nonlinear least-squares problem.

In practice, $\beta$ is often fixed (e.g., $\beta = 0.5$ for rates) and the remaining three parameters are calibrated to the ATM volatility, the skew (slope of the smile at ATM), and the curvature (smile convexity). The approximation is accurate for maturities up to about 10 years in typical rate environments, though it can break down for very low rates or extreme strikes.
\end{remark}

\begin{example}[SABR calibration for a swaption]
Consider a 5Y$\times$5Y swaption (option to enter a 5-year swap in 5 years) with forward swap rate $F_0 = 3.5\%$. Market quotes:

\begin{center}
\begin{tabular}{lccccc}
\toprule
Strike (bps from ATM) & $-200$ & $-100$ & ATM & $+100$ & $+200$ \\
\midrule
Implied vol (bps/day, normal) & 72.1 & 68.5 & 66.0 & 64.8 & 65.1 \\
\bottomrule
\end{tabular}
\end{center}

Fixing $\beta = 0.5$, calibration to these five points yields $\alpha_0 = 0.041$, $\rho = -0.32$, $\nu = 0.48$. The negative $\rho$ captures the downward-sloping left side of the smile (volatility increases as rates fall). The positive curvature on the right side (volatility increases again for very high strikes) is driven by $\nu > 0$.
\end{example}



\chapter{Fixed-Income and Credit}

The preceding chapters priced derivatives on equity-like assets: a stock driven by geometric Brownian motion, with a constant risk-free rate serving as a background parameter. In fixed-income markets, the interest rate is the \emph{asset}. Bond prices, swap rates, and cap/floor prices are all functions of the stochastic evolution of interest rates, and the challenge is to model an entire \emph{curve} of rates (the term structure) rather than a single price.

This chapter develops the theory in three layers of increasing generality. First, \textbf{short rate models} specify the dynamics of the instantaneous rate and derive bond prices as solutions to PDEs. Second, \textbf{forward rate models} (HJM) model the entire forward curve directly, with no-arbitrage imposing a drift restriction. Third, \textbf{market models} (LMM) specify the dynamics of the discretely compounded rates that appear in actual contracts, bridging theory and trading practice.

The chapter then turns to \textbf{credit risk}: the possibility that a counterparty fails to pay. Structural models (Merton, Black-Cox) derive default from the firm's asset dynamics; reduced-form models treat default as an exogenous jump governed by a hazard rate. The tools combine in the pricing of credit default swaps and collateralized debt obligations, where copula models introduce joint default dependence. Finally, \textbf{counterparty risk and XVA} adjustments extend the classical risk-neutral pricing framework to account for the real-world costs of default, funding, and capital.


\section{Term-Structure Models}

\subsection{Short Rate Models}

The simplest approach to modeling the term structure is to specify the dynamics of a single state variable (the instantaneous short rate) and derive everything else from it.

\begin{definition}
The \textbf{short rate} $R(t)$ is the instantaneous yield at time $t$:
\begin{align}
    R(t) = \lim_{T \to t} -\frac{\partial}{\partial T}\log B(t,T),
\end{align}
where $B(t,T)$ is the price of a zero-coupon bond. The short rate determines the money market account via $M(t) = \exp(\int_0^t R(u)\,du)$.
\end{definition}

\begin{remark}
The short rate is not directly observable, since no instrument pays interest for an infinitesimal period. Nevertheless, it serves as a convenient modeling device: once we specify $dR$, the entire yield curve is determined (in principle) by $B(t,T) = \widetilde{\mathbb{E}}[\exp(-\int_t^T R(u)\,du) \mid \mathcal{F}(t)]$. Different specifications of $dR$ produce different models, each with distinct implications for the shape and dynamics of the yield curve.
\end{remark}

\begin{definition}
The \textbf{Vasicek model} assumes that the short rate follows an \textbf{Ornstein-Uhlenbeck process} under the risk-neutral measure:
\begin{align}
    dR(t) = (\alpha - \beta R(t))\,dt + \sigma\,d\widetilde{W}(t),
\end{align}
where $\alpha/\beta$ is the long-run mean, $\beta > 0$ is the speed of mean reversion, and $\sigma > 0$ is the volatility.
\end{definition}

\begin{theorem}[Solution of the Vasicek process]
The short rate at time $t$, given $R(0)$, is
\begin{align}
    R(t) = R(0)e^{-\beta t} + \frac{\alpha}{\beta}(1 - e^{-\beta t}) + \sigma e^{-\beta t}\int_0^t e^{\beta s}\,d\widetilde{W}(s),
\end{align}
with conditional mean $\mathbb{E}[R(t)] = R(0)e^{-\beta t} + \frac{\alpha}{\beta}(1-e^{-\beta t})$ and variance $\text{Var}(R(t)) = \frac{\sigma^2}{2\beta}(1-e^{-2\beta t})$.
\end{theorem}

\begin{proof}
The rate appears on both sides of the SDE, so direct integration fails. The trick is to consider the process $Y(t) = e^{\beta t}R(t)$. By It\^o's product rule:
\begin{align}
    dY = e^{\beta t}\,dR + \beta e^{\beta t}R\,dt = \alpha e^{\beta t}\,dt + \sigma e^{\beta t}\,d\widetilde{W}.
\end{align}
This is a process with no $Y$-dependence on the right side, so it can be integrated directly:
\begin{align}
    Y(t) = Y(0) + \alpha\int_0^t e^{\beta s}\,ds + \sigma\int_0^t e^{\beta s}\,d\widetilde{W}(s).
\end{align}
Dividing by $e^{\beta t}$ gives the stated formula. The mean follows because the It\^o integral is a martingale (hence has zero expectation). The variance follows from the It\^o isometry: $\text{Var}(R(t)) = \sigma^2 e^{-2\beta t}\int_0^t e^{2\beta s}\,ds$.
\end{proof}

\begin{remark}
As $t \to \infty$, the mean converges to $\alpha/\beta$ and the variance converges to $\sigma^2/(2\beta)$. The stationary distribution is $\mathcal{N}(\alpha/\beta, \sigma^2/(2\beta))$. The half-life of a shock is $\log 2/\beta$: if $\beta = 0.5$, a 100bp shock to the short rate decays to 50bp in about 1.4 years.
\end{remark}

The power of the Vasicek model is that bond prices are available in closed form.

\begin{theorem}[Vasicek bond pricing]
Under the Vasicek model, the zero-coupon bond price $B(t,R;T)$ satisfies the PDE
\begin{align}
    \frac{\partial B}{\partial t} + (\alpha - \beta R)\frac{\partial B}{\partial R} + \frac{1}{2}\sigma^2\frac{\partial^2 B}{\partial R^2} - RB = 0, \qquad B(T,R;T) = 1,
\end{align}
and admits the \textbf{affine} solution
\begin{align}
    B(t,R;T) = \exp\bigl(g(t,T) + h(t,T)\,R\bigr),
\end{align}
where $g$ and $h$ solve Riccati ODEs.
\end{theorem}

\begin{proof}
The bond price is $B(t,T) = \widetilde{\mathbb{E}}[\exp(-\int_t^T R(u)\,du) \mid \mathcal{F}(t)]$. Since $R(t)$ is Markov, this is a function of $(t,R(t))$. The Kolmogorov backward equation with generator $\mathcal{L}f = (\alpha-\beta R)f_R + \frac{1}{2}\sigma^2 f_{RR}$ and discounting at rate $R$ gives the stated PDE.

To show the affine structure, note that $\mathcal{L}e^{uR} = [(\alpha u + \frac{1}{2}\sigma^2 u^2) + (-\beta u)R]\,e^{uR}$---the generator maps exponentials to exponentials times an affine function of $R$. This motivates the ansatz $B = e^{g + hR}$, which reduces the PDE to a system of Riccati ODEs for $g$ and $h$.
\end{proof}

\begin{example}[Vasicek yield curve]
Take $\alpha = 0.03$, $\beta = 0.5$, $\sigma = 0.015$, and current rate $R(0) = 4\%$. The long-run mean rate is $\alpha/\beta = 6\%$. Since $R(0) < \alpha/\beta$, the model expects rates to rise, producing an upward-sloping yield curve. The Riccati solutions give, for example:
\begin{align}
    B(0, 0.04; 1) \approx 0.9610, \quad Y(0,1) \approx 3.97\%, \\
    B(0, 0.04; 10) \approx 0.5769, \quad Y(0,10) \approx 5.50\%.
\end{align}
The 10-year yield (5.50\%) is between the current rate (4\%) and the long-run mean (6\%), reflecting the gradual mean reversion built into the model.
\end{example}

\begin{remark}
The Vasicek model allows negative interest rates, since $R(t)$ is Gaussian. Before 2008, this was considered a deficiency. After the era of negative policy rates in Europe and Japan, it became a feature. The choice between Gaussian models (Vasicek, Hull-White) and positive-rate models (CIR) depends on the market environment and the application.
\end{remark}

The Vasicek model is one member of a broader family. The affine structure (bond prices log-linear in the state variable) holds whenever the drift and diffusion of $R$ are affine in $R$.

\begin{proposition}[Affine term structure models]
If the risk-neutral short rate dynamics have the form
\begin{align}
    dR(t) = (\kappa(t)R(t) + \eta(t))\,dt + \sqrt{\gamma(t)R(t) + \delta(t)}\,d\widetilde{W}(t),
\end{align}
then bond prices are $B(t,R;T) = e^{g(t,T) + h(t,T)R}$, where $g$ and $h$ satisfy the \textbf{Riccati ODEs}:
\begin{align}
    \frac{dh}{dt} &= -\kappa(t)\,h - \frac{1}{2}\gamma(t)\,h^2 + 1, \\
    \frac{dg}{dt} &= -\eta(t)\,h - \frac{1}{2}\delta(t)\,h^2,
\end{align}
with terminal conditions $g(T,T) = h(T,T) = 0$.
\end{proposition}

Two important special cases round out the short rate family.

\begin{definition}
The \textbf{Cox-Ingersoll-Ross (CIR) model} modifies the Vasicek diffusion to preserve positivity:
\begin{align}
    dR(t) = \kappa(\theta - R(t))\,dt + \sigma\sqrt{R(t)}\,d\widetilde{W}(t).
\end{align}
Under the Feller condition $2\kappa\theta \geq \sigma^2$, the rate remains strictly positive. CIR is affine (with $\gamma = \sigma^2$, $\delta = 0$) and yields closed-form bond prices.
\end{definition}

\begin{remark}
The state-dependent diffusion $\sigma\sqrt{R}$ means that volatility increases when rates are high and decreases when they are low. This matches the empirical observation that rate volatility is positively correlated with the level of rates, particularly in high-rate environments. However, CIR cannot produce negative rates.
\end{remark}

\begin{definition}
The \textbf{Hull-White model} extends Vasicek by allowing time-dependent parameters:
\begin{align}
    dR(t) = (\theta(t) - a\,R(t))\,dt + \sigma\,d\widetilde{W}(t).
\end{align}
The function $\theta(t)$ is chosen to fit the observed yield curve exactly.
\end{definition}

\begin{proposition}[Hull-White calibration]
Given observed market forward rates $f^{\text{mkt}}(0,t)$, the drift function is
\begin{align}
    \theta(t) = \frac{\partial f^{\text{mkt}}}{\partial t}(0,t) + a\,f^{\text{mkt}}(0,t) + \frac{\sigma^2}{2a}(1 - e^{-at})^2.
\end{align}
The remaining parameters $a$ and $\sigma$ are calibrated to swaption or cap/floor volatilities.
\end{proposition}

\begin{remark}
Hull-White is the industry standard for interest rate derivatives pricing. Its ability to fit the initial yield curve exactly eliminates one source of model error (mispricing of vanilla bonds). The trade-off is that the model has only two free parameters ($a, \sigma$) for fitting the volatility surface, limiting its flexibility for exotic products. For richer dynamics, one turns to multifactor models (e.g., two-factor Hull-White with independent short-rate and slope factors) or to the forward rate and market models discussed next.
\end{remark}


\subsection{Forward Rate Models}

Short rate models derive the forward curve from a single state variable. An alternative approach models the entire forward curve directly.

\begin{definition}
The \textbf{instantaneous forward rate} $f(t,T)$ is the interest rate at time $t$ for instantaneous borrowing at future time $T$:
\begin{align}
    f(t,T) = -\frac{\partial}{\partial T}\log B(t,T).
\end{align}
The short rate is the forward rate at zero maturity: $R(t) = f(t,t)$. Bond prices are recovered by integration:
\begin{align}
    B(t,T) = \exp\!\Bigl(-\int_t^T f(t,v)\,dv\Bigr).
\end{align}
\end{definition}

\begin{definition}
The \textbf{Heath-Jarrow-Morton (HJM) framework} models the dynamics of the forward rate curve under the risk-neutral measure:
\begin{align}
    df(t,T) = \alpha(t,T)\,dt + \sigma(t,T)\,d\widetilde{W}(t),
\end{align}
initialized at the observed curve $f(0,T) = f^{\text{mkt}}(0,T)$.
\end{definition}

\begin{remark}
The HJM framework is fundamentally different from short rate models: rather than specifying one SDE and deriving bond prices, it specifies the dynamics of an infinite-dimensional object (the entire forward curve) and imposes no-arbitrage as a constraint. The result is a restriction on the drift $\alpha(t,T)$ determined by the volatility $\sigma(t,T)$.
\end{remark}

\begin{theorem}[HJM drift restriction]
Absence of arbitrage requires that the drift of the forward rate satisfies
\begin{align}
    \alpha(t,T) = \sigma(t,T)\,\sigma^*(t,T),
\end{align}
under the risk-neutral measure, where $\sigma^*(t,T) = \int_t^T \sigma(t,u)\,du$ is the integrated volatility.
\end{theorem}

\begin{proof}
The discounted bond price $D(t)B(t,T)$ must be a martingale. Computing $dB(t,T)$ via It\^o requires the differential of $-\int_t^T f(t,v)\,dv$, which equals $f(t,t)\,dt - \int_t^T df(t,v)\,dv$. Substituting the HJM dynamics and collecting terms gives
\begin{align}
    d(D(t)B(t,T)) = D(t)B(t,T)\bigl[(-\alpha^*(t,T) + \tfrac{1}{2}(\sigma^*(t,T))^2)\,dt - \sigma^*(t,T)\,d\widetilde{W}(t)\bigr],
\end{align}
where $\alpha^* = \int_t^T \alpha(t,u)\,du$. The martingale condition forces the $dt$ coefficient to zero: $\alpha^* = \frac{1}{2}(\sigma^*)^2$. Differentiating in $T$ gives $\alpha(t,T) = \sigma(t,T)\,\sigma^*(t,T)$.
\end{proof}

\begin{remark}
The HJM drift restriction is one of the most elegant results in mathematical finance: \emph{the drift of the forward curve is completely determined by its volatility structure}. The modeler specifies only the volatility function $\sigma(t,T)$; no-arbitrage does the rest. Different choices of $\sigma(t,T)$ produce different models. For instance, $\sigma(t,T) = \sigma e^{-a(T-t)}$ recovers the Hull-White model.
\end{remark}

Under the drift restriction, the evolution of the forward curve and bond prices takes a clean form.

\begin{theorem}
Under the HJM no-arbitrage dynamics, forward rates evolve as
\begin{align}
    df(t,T) = \sigma(t,T)\,\sigma^*(t,T)\,dt + \sigma(t,T)\,d\widetilde{W}(t),
\end{align}
and bond prices follow
\begin{align}
    B(t,T) = \frac{B(0,T)}{D(t)}\exp\!\Bigl[-\int_0^t \sigma^*(u,T)\,d\widetilde{W}(u) - \frac{1}{2}\int_0^t (\sigma^*(u,T))^2\,du\Bigr].
\end{align}
\end{theorem}

\begin{remark}
The HJM framework is theoretically complete but can be impractical. For general volatility specifications, the forward curve is non-Markovian. Moreover, certain volatility structures (e.g., $\sigma$ proportional to $f$) can cause forward rates to explode in finite time. These limitations motivated the development of market models, which specify the dynamics of directly observable rates.
\end{remark}


\subsection{Market Models}

\begin{remark}[Motivation]
Interest rate derivatives (caps, floors, swaptions) are quoted and settled in terms of discretely compounded forward rates (historically LIBOR, now SOFR-based equivalents) over fixed accrual periods, not in terms of the instantaneous short rate or the continuous forward curve. The LIBOR Market Model (LMM) specifies the dynamics of these observable rates directly, which ensures that the model's state variables match the market's quoting conventions.
\end{remark}

\begin{definition}
Fix a \textbf{tenor structure} $0 = T_0 < T_1 < \cdots < T_N$ with accrual periods $\delta_i = T_{i+1} - T_i$. The \textbf{forward LIBOR rate} for the period $[T_i, T_{i+1}]$, observed at time $t \leq T_i$, is
\begin{align}
    L_i(t) = \frac{1}{\delta_i}\!\left(\frac{B(t,T_i)}{B(t,T_{i+1})} - 1\right).
\end{align}
\end{definition}

\begin{remark}
The forward LIBOR rate $L_i(t)$ is the rate that, if compounded over the accrual period $\delta_i$, equates the bond prices $B(t,T_i)$ and $B(t,T_{i+1})$. It is directly observable from bond prices and is the reference rate for caps and floors.
\end{remark}

The change-of-num\'eraire framework identifies the natural measure for each forward rate.

\begin{proposition}
Under the $T_{i+1}$-forward measure $\mathbb{Q}^{(T_{i+1})}$ (with num\'eraire $B(t,T_{i+1})$), the forward rate $L_i(t)$ is a martingale.
\end{proposition}

\begin{proof}[Proof sketch]
The payoff $\delta_i L_i(T_i)$ at time $T_{i+1}$ can be replicated by a portfolio of bonds paying $B(T_i,T_i)/B(T_i,T_{i+1}) - 1$ at $T_{i+1}$. Expressed in units of the num\'eraire $B(t,T_{i+1})$, this is a martingale by no-arbitrage.
\end{proof}

\begin{definition}[LIBOR Market Model]
Under $\mathbb{Q}^{(T_{i+1})}$, each forward rate follows lognormal dynamics:
\begin{align}
    dL_i(t) = \sigma_i(t)\,L_i(t)\,d\widetilde{W}_i^{(T_{i+1})}(t), \qquad t \leq T_i,
\end{align}
with deterministic volatility $\sigma_i(t)$ and instantaneous correlations $d\widetilde{W}_i\,d\widetilde{W}_j = \rho_{ij}\,dt$.
\end{definition}

\begin{remark}
Each forward rate is lognormal under its \emph{own} forward measure. But the $N$ forward rates live under $N$ different measures. To simulate all rates jointly (necessary for pricing path-dependent or multi-rate products), one must express them under a common measure, introducing measure-change drifts.
\end{remark}

\begin{proposition}[Drift under the terminal measure]
Under the terminal measure $\mathbb{Q}^{(T_N)}$, the dynamics of $L_i(t)$ acquire a drift:
\begin{align}
    dL_i(t) = \mu_i(t)\,dt + \sigma_i(t)\,L_i(t)\,d\widetilde{W}_i^{(T_N)}(t),
\end{align}
where
\begin{align}
    \mu_i(t) = -\sigma_i(t)\,L_i(t)\sum_{j=i+1}^{N-1}\frac{\delta_j L_j(t)\,\sigma_j(t)\,\rho_{ij}}{1 + \delta_j L_j(t)}.
\end{align}
\end{proposition}

\begin{remark}
The drift of $L_i$ depends on \emph{all subsequent} rates $L_{i+1}, \ldots, L_{N-1}$. This coupling is the principal analytical difficulty of the LMM: while each rate is individually lognormal under its own measure, the joint system is not jointly lognormal under any single measure. As a result, multi-rate products (swaptions, Bermudans) generally require Monte Carlo simulation rather than closed-form solutions.
\end{remark}

The LMM prices single-rate products in closed form.

\begin{definition}
A \textbf{caplet} on the period $[T_i, T_{i+1}]$ with strike $K$ pays $\delta_i\max(L_i(T_i) - K, 0)$ at time $T_{i+1}$.
\end{definition}

\begin{proposition}
Under $\mathbb{Q}^{(T_{i+1})}$, the caplet prices as
\begin{align}
    V_{\text{cap}}(t) = B(t,T_{i+1})\,\widetilde{\mathbb{E}}^{(T_{i+1})}\bigl[\delta_i\max(L_i(T_i)-K,0)\;\big|\;\mathcal{F}(t)\bigr].
\end{align}
Since $L_i$ is lognormal under $\mathbb{Q}^{(T_{i+1})}$, the price is given by the \textbf{Black formula} (the same formula as Black-Scholes but with the forward rate replacing the stock price and the bond price replacing the discount factor).
\end{proposition}

\begin{remark}
The LMM prices caplets consistently with the Black formula that traders already use, while providing a self-consistent joint dynamics for all rates. Swaption pricing is more complex because the swap rate is a nonlinear function of multiple LIBOR rates.
\end{remark}

\begin{definition}
The \textbf{forward swap rate} for a swap starting at $T_\alpha$ and ending at $T_\beta$ is
\begin{align}
    S_{\alpha,\beta}(t) = \frac{B(t,T_\alpha) - B(t,T_\beta)}{\sum_{k=\alpha}^{\beta-1}\delta_k\,B(t,T_{k+1})}.
\end{align}
No measure exists under which $S_{\alpha,\beta}$ is exactly lognormal (it is a ratio of sums of lognormals).
\end{definition}

\begin{remark}[Calibration procedure]
LMM calibration proceeds in three steps: (1) calibrate individual volatilities $\sigma_i(t)$ to caplet implied volatilities, (2) fit the correlation matrix $(\rho_{ij})$ to swaption prices (often using a low-rank parameterization for tractability), and (3) validate the joint dynamics by Monte Carlo pricing of exotic payoffs (Bermudans, callable range accruals, CMS products).
\end{remark}

\begin{remark}[Post-LIBOR transition]
Although LIBOR has been discontinued and replaced by backward-looking rates such as SOFR, the modeling framework survives in adapted forms. SOFR-based forwards are compounded in arrears rather than set in advance, requiring modifications to the LMM dynamics and payoff timing.
\end{remark}


\section{Credit Risk Modeling}

The fixed-income framework thus far has treated bond payments as certain. In reality, a bond issuer may default, and the bondholder may receive only a fraction of the promised cash flows. This section introduces models for the random default time and applies them to the pricing of credit-sensitive instruments.

\subsection{Structural Models}

Structural models derive default from the fundamental economics of the firm: default occurs when the firm's assets are insufficient to cover its liabilities.

\begin{theorem}[Merton model]
Let the firm's asset value follow geometric Brownian motion:
\begin{align}
    dA(t) = \mu A(t)\,dt + \sigma A(t)\,dW(t).
\end{align}
The firm has a single debt obligation with face value $D$ due at time $T$. Default occurs if $A(T) < D$. The \textbf{distance to default} is
\begin{align}
    \text{DD} = \frac{\log(A(0)/D) + (\mu - \frac{1}{2}\sigma^2)T}{\sigma\sqrt{T}},
\end{align}
and the probability of default is $\mathbb{P}(A(T) < D) = \Phi(-\text{DD})$.
\end{theorem}

\begin{remark}
The firm's equity is a call option on its assets with strike $D$ (by limited liability), and the firm's debt is a risk-free bond minus a put option on the assets. This connects with equity volatility, but it has two well-known empirical failures: default can only occur at maturity $T$ (not before), and short-term credit spreads are near zero (because $\Phi(-\text{DD})$ is tiny for short $T$ when assets are well above liabilities).
\end{remark}

\begin{definition}
In the \textbf{Black-Cox model}, default occurs at the first time the asset value hits a deterministic barrier:
\begin{align}
    \tau = \inf\{t \geq 0 : A(t) \leq K(t)\}.
\end{align}
This is a first-passage time problem for geometric Brownian motion, and produces non-trivial short-term default probabilities.
\end{definition}

\begin{remark}
Barrier models improve short-term spread behavior but remain sensitive to the choice of barrier and to asset volatility. In practice, structural models are used primarily for relative value analysis (comparing default probabilities across firms via their equity dynamics) rather than for absolute pricing.
\end{remark}


\subsection{Reduced-Form Models}

Reduced-form (or intensity-based) models treat default as an exogenous event, governed by a stochastic intensity process. This approach separates the modeling of \emph{when} default occurs from the modeling of \emph{why}, gaining tractability at the cost of economic transparency.

\begin{definition}
A \textbf{default time} $\tau$ is a non-negative random variable. The \textbf{survival probability} is $Q(t,T) = \mathbb{P}(\tau > T \mid \mathcal{F}(t))$, and the \textbf{survival indicator} is $\mathbb{I}_{\{\tau > t\}}$.
\end{definition}

\begin{definition}
The \textbf{recovery rate} $R \in [0,1]$ is the fraction of notional recovered upon default. In practice, $R$ is typically assumed constant (e.g., $R = 40\%$ for senior unsecured corporate bonds).
\end{definition}

The central modeling choice is the default intensity.

\begin{definition}
A non-negative adapted process $\lambda(t)$ is the \textbf{default intensity} (or \textbf{hazard rate}) if
\begin{align}
    \mathbb{P}(\tau \in [t, t+dt) \mid \mathcal{F}(t),\, \tau > t) = \lambda(t)\,dt.
\end{align}
\end{definition}

\begin{remark}
The intensity $\lambda(t)$ plays the same role for default as the short rate $R(t)$ plays for interest rates. Just as bond prices are expectations of $\exp(-\int R\,du)$, survival probabilities are expectations of $\exp(-\int \lambda\,du)$.
\end{remark}

\begin{proposition}[Survival probability]
If $\lambda(t)$ is the default intensity, the survival probability is
\begin{align}
    Q(0,T) = \mathbb{E}\!\left[\exp\!\Bigl(-\int_0^T \lambda(s)\,ds\Bigr)\right].
\end{align}
When $\lambda$ is deterministic, $Q(0,T) = \exp(-\int_0^T \lambda(s)\,ds)$.
\end{proposition}

\begin{remark}
Under the risk-neutral measure, $\lambda$ embeds both the physical default probability and a risk premium. Market-implied intensities (extracted from CDS spreads) typically exceed historical default rates by a factor of 2--5, reflecting the compensation investors demand for bearing default risk.
\end{remark}

\begin{definition}
The \textbf{credit spread} $s(T)$ is the yield difference between a defaultable and a default-free zero-coupon bond of maturity $T$.
\end{definition}

\begin{proposition}
Under deterministic rates, deterministic intensity, and constant recovery:
\begin{align}
    s(T) \approx (1-R)\,\bar{\lambda}(T), \qquad \bar{\lambda}(T) = \frac{1}{T}\int_0^T\lambda(s)\,ds.
\end{align}
\end{proposition}

\begin{remark}
This is the simplest credit spread formula: the spread is approximately the loss-given-default $(1-R)$ times the average hazard rate. For $R = 40\%$ and $\lambda = 1\%$, the spread is approximately 60bp.
\end{remark}


\subsection{Credit Default Swaps}

The CDS is the fundamental instrument for trading credit risk. Its pricing connects directly to the hazard rate framework.

\begin{definition}
A \textbf{credit default swap (CDS)} is a contract in which the protection buyer pays a periodic premium (the CDS spread $S$) to the protection seller, who in return pays $(1-R)$ upon default of the reference entity.
\end{definition}

\begin{proposition}[CDS pricing]
The CDS spread $S$ is determined by equating the present values of the two legs. The \textbf{protection leg} (what the seller pays upon default) has value
\begin{align}
    \text{PV}_{\text{prot}} = (1-R)\int_0^T D(0,t)\,\lambda(t)\,Q(0,t)\,dt,
\end{align}
and the \textbf{premium leg} (what the buyer pays while the reference survives) has value
\begin{align}
    \text{PV}_{\text{prem}} = S\int_0^T D(0,t)\,Q(0,t)\,dt.
\end{align}
Setting $\text{PV}_{\text{prot}} = \text{PV}_{\text{prem}}$ and solving for $S$ gives the par CDS spread:
\begin{align}
    S = \frac{(1-R)\int_0^T D(0,t)\,\lambda(t)\,Q(0,t)\,dt}{\int_0^T D(0,t)\,Q(0,t)\,dt}.
\end{align}
\end{proposition}

\begin{remark}
For constant $\lambda$ and flat discount curve, the integrals simplify and $S \approx (1-R)\lambda$, reproducing the credit spread approximation. In practice, the premium leg is discretized into quarterly payments, and accrued premium (the fraction of the coupon owed if default occurs between payment dates) must be included.
\end{remark}

\begin{example}[CDS bootstrapping]
Consider a 5-year CDS with $R = 40\%$, flat discount rate $r = 2\%$, and quoted spread $S = 100$bp. Assuming constant $\lambda$:
\begin{align}
    S\int_0^5 e^{-(r+\lambda)t}\,dt = (1-R)\lambda\int_0^5 e^{-(r+\lambda)t}\,dt.
\end{align}
The integrals cancel, giving $S = 0.6\lambda$, hence $\lambda = 100/(0.6 \times 10{,}000) \approx 1.67\%$ per year. The 5-year survival probability is $Q(0,5) = e^{-0.0167 \times 5} \approx 92.0\%$.

In practice, the term structure of CDS spreads (1Y, 3Y, 5Y, 7Y, 10Y) is used to bootstrap a piecewise-constant hazard rate curve, just as the term structure of swap rates is used to bootstrap a discount curve.
\end{example}


\subsection{Portfolio Credit}

Single-name credit models calibrate marginal default distributions from CDS data. For portfolio credit products (CDOs, basket default swaps, index tranches) we need the \emph{joint} distribution of default times across multiple names.

\begin{remark}
This is the fundamental challenge: single-name CDS data determines each $Q_i(t) = \mathbb{P}(\tau_i > t)$ individually, but says nothing about the joint probability $\mathbb{P}(\tau_1 > t_1, \ldots, \tau_n > t_n)$. The copula framework separates the marginals (calibrated from CDS) from the dependence structure (an additional modeling choice).
\end{remark}

\begin{definition}
A \textbf{copula} $C: [0,1]^d \to [0,1]$ is the joint CDF of a random vector $(U_1, \ldots, U_d)$ with uniform marginals:
\begin{align}
    C(u_1, \ldots, u_d) = \mathbb{P}(U_1 \leq u_1, \ldots, U_d \leq u_d).
\end{align}
Given marginal CDFs $F_i$ and a copula $C$, joint default times are constructed as $\tau_i = F_i^{-1}(U_i)$, preserving each marginal while introducing dependence through $C$.
\end{definition}

\begin{definition}
The \textbf{Gaussian copula} with correlation matrix $\Sigma$ is
\begin{align}
    C_\Sigma^{\text{Gauss}}(u_1, \ldots, u_d) = \Phi_\Sigma(\Phi^{-1}(u_1), \ldots, \Phi^{-1}(u_d)),
\end{align}
where $\Phi^{-1}$ is the standard normal inverse CDF and $\Phi_\Sigma$ is the multivariate normal CDF.

In the \textbf{one-factor Gaussian copula model}, the latent variables are
\begin{align}
    Z_i = \sqrt{\rho}\,X + \sqrt{1-\rho}\,\varepsilon_i, \qquad X, \varepsilon_i \sim \mathcal{N}(0,1) \text{ i.i.d.},
\end{align}
with a single correlation parameter $\rho \in [0,1]$ governing the proportion of default risk attributable to the common systemic factor $X$.
\end{definition}

\begin{remark}
The one-factor structure has an elegant conditional independence property: given the systemic factor $X$, the default times $\tau_i$ are independent. This allows the conditional portfolio loss distribution to be computed analytically (as a sum of independent Bernoullis), with the unconditional distribution obtained by integrating over $X$. This is why the model is computationally tractable.
\end{remark}

\begin{definition}
A \textbf{collateralized debt obligation (CDO)} redistributes the losses of a credit portfolio into tranches with different seniority. For a portfolio with cumulative loss $L(t) = \sum_{i=1}^d (1-R_i)\,\mathbb{I}_{\{\tau_i \leq t\}}$ and a tranche with attachment point $A$ and detachment point $D$, the \textbf{tranche loss} is
\begin{align}
    L^{A,D}(t) = \min(D-A,\;\max(L(t)-A, 0)).
\end{align}
\end{definition}

\begin{remark}
The equity tranche (e.g., 0--3\%) absorbs the first losses and is dominated by idiosyncratic default risk. Senior tranches (e.g., 15--30\%) are exposed primarily to systemic risk. The correlation parameter $\rho$ controls the relative pricing of these tranches: higher correlation shifts risk from equity to senior tranches (making equity cheaper and senior more expensive), because systemic events that cause many simultaneous defaults become more likely.
\end{remark}

\begin{example}[CDS to CDO]
Consider 10 identical names, each with 5-year CDS spread $S = 100$bp, recovery $R = 40\%$, and flat rate $r = 2\%$. From the CDS bootstrapping above, $\lambda = 1.67\%$ and $Q(0,5) = 92.0\%$.

Using the one-factor Gaussian copula with $\rho = 0.2$, simulate $N = 10{,}000$ scenarios:
\begin{enumerate}
    \item Draw $X \sim \mathcal{N}(0,1)$ (the systemic factor).
    \item For each name $i$, draw $\varepsilon_i \sim \mathcal{N}(0,1)$ and compute $Z_i = \sqrt{0.2}\,X + \sqrt{0.8}\,\varepsilon_i$.
    \item Default occurs if $\Phi(Z_i) < F_i(5) = 1 - 0.920 = 0.080$, i.e., if $Z_i < \Phi^{-1}(0.080) = -1.405$.
    \item Compute the portfolio loss $L = \frac{1}{10}\sum_i 0.6 \cdot \mathbb{I}_{\{Z_i < -1.405\}}$.
    \item Compute the 0--3\% tranche loss: $L^{0, 0.03} = \min(0.03, \max(L, 0))$.
\end{enumerate}
Averaging across scenarios gives $\mathbb{E}[L^{0,0.03}] \approx 2.1\%$.

Changing $\rho$ to $0.5$ would increase senior tranche losses (more systemic clustering) while decreasing equity tranche losses (fewer idiosyncratic defaults hit only the equity). This $\rho$-sensitivity is why ``implied correlation'' became the market's de facto quoting convention for CDO tranches.
\end{example}

\begin{remark}[Limitations of the Gaussian copula]
The Gaussian copula exhibits zero \textbf{upper tail dependence}:
\begin{align}
    \lambda_U = \lim_{u \uparrow 1}\mathbb{P}(U_2 > u \mid U_1 > u) = 0.
\end{align}
This means that extreme joint defaults are less likely than many alternative copulas (e.g., $t$-copula, Clayton) would predict. In practice, different CDO tranches imply different correlation parameters (\textbf{correlation smile}), indicating that the model cannot simultaneously fit all tranches with a single parameter. These limitations were exposed dramatically during the 2008 financial crisis, when correlated defaults exceeded the model's predictions for senior tranches.
\end{remark}


\section{Counterparty Risk and XVA}

Classical risk-neutral pricing assumes that both parties to a derivative contract will honor their obligations. In reality, counterparties can default, funding is not free, and regulatory capital must be held. The \textbf{XVA framework} adjusts the ``clean'' (risk-free) price to account for these real-world costs.

\begin{definition}
\textbf{Counterparty risk} is the risk that a trading counterparty defaults before the maturity of a contract, causing a loss equal to the replacement cost of the position at the time of default.
\end{definition}

\subsection{Exposure}

\begin{definition}
Let $V(t)$ be the mark-to-market value of a derivative portfolio from the investor's perspective. The \textbf{positive exposure} $E^+(t) = \max(V(t), 0)$ is the amount at risk if the counterparty defaults (the investor is owed money). The \textbf{negative exposure} $E^-(t) = \max(-V(t), 0)$ is the amount the investor owes.
\end{definition}

\begin{definition}
The \textbf{expected positive exposure (EPE)} and \textbf{expected negative exposure (ENE)} are
\begin{align}
    \text{EPE}(t) = \widetilde{\mathbb{E}}[E^+(t)], \qquad \text{ENE}(t) = \widetilde{\mathbb{E}}[E^-(t)].
\end{align}
\end{definition}

\begin{remark}
The exposure profile of a derivative changes over time and depends on market conditions. For a plain vanilla interest rate swap, the EPE typically rises initially (as rates can move further from the fixed rate) and then declines toward zero as the remaining cash flows shrink. The peak of the EPE profile, called the \textbf{potential future exposure (PFE)}, is used for regulatory capital calculations.
\end{remark}

\begin{definition}
\textbf{Collateral agreements} (governed by Credit Support Annexes) specify the posting of \textbf{variation margin} (covering current exposure) and \textbf{initial margin} (covering potential future exposure during a close-out period). With collateral $C(t)$, the \textbf{collateralized exposure} is $\widetilde{E}(t) = V(t) - C(t)$.
\end{definition}


\subsection{Valuation Adjustments}

The separation between ``risk-free'' discounting and real-world costs of default, funding, and capital leads to a series of additive adjustments to the clean price.

\begin{definition}
The \textbf{Credit Valuation Adjustment (CVA)} is the expected loss due to counterparty default:
\begin{align}
    \text{CVA} = (1-R_c)\int_0^T D(0,t)\,\widetilde{\mathbb{E}}[\widetilde{E}^+(t)]\,\lambda_c(t)\,Q_c(0,t)\,dt,
\end{align}
where $R_c$ is the counterparty's recovery rate, $\lambda_c$ its default intensity, and $Q_c$ its survival probability.
\end{definition}

\begin{remark}
CVA is a real economic cost: if a counterparty defaults when the position has positive value to the investor, the investor loses $(1-R_c)$ times the exposure. CVA desks at major banks are profit centers that actively hedge counterparty risk by trading CDS and managing exposure profiles.
\end{remark}

\begin{definition}
The \textbf{Debit Valuation Adjustment (DVA)} accounts for the investor's own default risk:
\begin{align}
    \text{DVA} = (1-R_o)\int_0^T D(0,t)\,\widetilde{\mathbb{E}}[\widetilde{E}^-(t)]\,\lambda_o(t)\,Q_o(0,t)\,dt.
\end{align}
\end{definition}

\begin{remark}
DVA is controversial. It represents the ``benefit'' the investor derives from the possibility of defaulting on obligations. This is real in an accounting sense (it reduces the investor's liabilities) but perverse from a risk management perspective (one's own creditworthiness deteriorating should not be a source of profit). Under IFRS 13, DVA must be recognized in fair value accounting, but many institutions exclude it from internal risk metrics.
\end{remark}

The remaining adjustments capture funding and regulatory costs.

\begin{definition}
The \textbf{Funding Valuation Adjustment (FVA)} is the cost of funding uncollateralized exposures at rates above the risk-free rate:
\begin{align}
    \text{FVA} = \int_0^T D(0,t)\,\widetilde{\mathbb{E}}[\widetilde{E}(t)]\,s_f(t)\,dt,
\end{align}
where $s_f(t)$ is the funding spread. The \textbf{Capital Valuation Adjustment (KVA)} is the cost of holding regulatory capital:
\begin{align}
    \text{KVA} = \int_0^T D(0,t)\,\widetilde{\mathbb{E}}[K(t)]\,h\,dt,
\end{align}
where $K(t)$ is the required capital and $h$ is the hurdle rate. The \textbf{Margin Valuation Adjustment (MVA)} is the funding cost of posting initial margin:
\begin{align}
    \text{MVA} = \int_0^T D(0,t)\,\widetilde{\mathbb{E}}[\text{IM}(t)]\,s_f(t)\,dt.
\end{align}
\end{definition}

\begin{proposition}
The total value of a derivative portfolio, including all adjustments, is
\begin{align}
    V_{\text{total}} = V_{\text{clean}} - \text{CVA} + \text{DVA} - \text{FVA} - \text{KVA} - \text{MVA}.
\end{align}
\end{proposition}

\begin{remark}
Each adjustment has a sign: CVA, FVA, KVA, and MVA reduce value (they are costs), while DVA increases it (it is a benefit from the investor's perspective). The magnitudes can be substantial: for a large uncollateralized swap portfolio, CVA alone can be several percent of notional.
\end{remark}


\subsection{Wrong-Way Risk and Practical Considerations}

\begin{definition}
\textbf{Wrong-way risk (WWR)} arises when exposure and counterparty credit quality are positively correlated:
\begin{align}
    \text{Corr}(E^+(t), \lambda_c(t)) > 0.
\end{align}
When the counterparty is most likely to default, the exposure is largest.
\end{definition}

\begin{example}[Wrong-way risk in practice]
Consider a bank that has sold CDS protection to a hedge fund referencing the bank's own sector. If the sector deteriorates, the bank's exposure (it owes money on the CDS) increases at the same time that the hedge fund's creditworthiness worsens (it is a leveraged player in the same sector). The factorized CVA formula $\text{CVA} \approx (1-R)\,\text{EPE} \times \text{PD}$ dramatically underestimates the loss, because it treats exposure and default probability as independent.
\end{example}

\begin{remark}
Wrong-way risk invalidates the standard factorized CVA formula and requires joint modeling of market risk factors and counterparty credit quality. This requires simulating both market scenarios and credit scenarios on the same paths.
\end{remark}

\begin{definition}
A \textbf{netting set} is a collection of trades with a single counterparty governed by a legally enforceable netting agreement. Upon default, exposures within the netting set are aggregated before close-out: positive and negative values offset, reducing the net exposure. Netting can dramatically reduce CVA.
\end{definition}

\begin{example}[XVA for an interest rate swap]
Consider a 5-year payer swap with notional \$100M, fixed rate 2\%, discounting at $r = 1\%$. Assume deterministic exposure profiles:
\begin{align}
    \text{EPE}(t) = 5\,e^{-0.05t} \text{ (\$M)}, \qquad \text{ENE}(t) = 2\,e^{-0.05t} \text{ (\$M)}.
\end{align}
Counterparty: $R_c = 40\%$, $\lambda_c = 1\%$/yr. Investor: $R_o = 40\%$, $\lambda_o = 0.5\%$/yr.

The CVA is
\begin{align}
    \text{CVA} &= 0.6 \times 0.01 \times 5\int_0^5 e^{-(0.01+0.05+0.01)t}\,dt = 0.03\int_0^5 e^{-0.07t}\,dt \approx 0.03 \times 4.21 \approx \$126\text{K}.
\end{align}
The DVA is
\begin{align}
    \text{DVA} &= 0.6 \times 0.005 \times 2\int_0^5 e^{-(0.01+0.05+0.005)t}\,dt \approx 0.006 \times 4.27 \approx \$26\text{K}.
\end{align}
If the clean value is \$800K:
\begin{align}
    V_{\text{total}} = 800 - 126 + 26 = \$700\text{K}.
\end{align}
The CVA alone reduces the portfolio value by 16\%, highlighting why counterparty risk management is a first-order concern for derivatives-heavy institutions.
\end{example}

\begin{remark}
In practice, exposure profiles are stochastic (computed via Monte Carlo simulation of interest rates, FX, and credit jointly), collateral dynamics are path-dependent (margin calls have thresholds and delays), and netting across trade types introduces complex optionality.
\end{remark}


\chapter{Portfolio Management}

The preceding chapters priced individual derivatives by constructing hedges and computing risk-neutral expectations. Portfolio management asks a different question: \emph{given many available assets, how should an investor allocate capital?} The answer depends on the investor's objectives (return, risk tolerance, benchmark tracking) and constraints (leverage limits, regulatory requirements, liquidity needs).

This chapter develops two complementary perspectives. First, \textbf{portfolio theory} solves the allocation problem: mean-variance optimization identifies the best trade-off between expected return and risk, the CAPM links individual asset pricing to market equilibrium, and factor models decompose returns into systematic and idiosyncratic components. Second, \textbf{portfolio risk} provides the quantitative tools for measuring and controlling downside exposure: Value at Risk, Expected Shortfall, drawdowns, diversification, and hedging.


\section{Portfolio Theory}

\subsection{Mean-Variance Analysis}

The Markowitz framework reduces portfolio selection to a quadratic optimization problem. Despite its simplicity, it remains the foundation for more elaborate models and is the starting point for practical portfolio construction.

\begin{definition}
Consider a single-period economy with $m$ risky assets and one risk-free asset. Let the random return vector on the risky assets be $\vec{R} = (R_1, \ldots, R_m)$ with expected returns $\vec{\mu} = \mathbb{E}[\vec{R}]$ and covariance matrix $\Sigma = \text{Cov}(\vec{R})$, assumed symmetric and positive definite. The risk-free return is $r_0$.

A portfolio is specified by weight vector $\vec{w} = (w_1, \ldots, w_m)$, where $w_i$ is the fraction of wealth invested in risky asset $i$. The risk-free weight is $w_0 = 1 - \vec{w} \cdot \vec{1}_m$ (negative if borrowing). The \textbf{portfolio return} is
\begin{align}
    R_p = \vec{w} \cdot \vec{R} + (1 - \vec{w} \cdot \vec{1}_m)\,r_0.
\end{align}
Under the \textbf{Markowitz model}, the portfolio's mean and variance are
\begin{align}
    \mathbb{E}[R_p] = \vec{w} \cdot \vec{\mu} + (1 - \vec{w} \cdot \vec{1}_m)\,r_0, \qquad \text{Var}(R_p) = \vec{w}^\top \Sigma\, \vec{w}.
\end{align}
\end{definition}

\begin{remark}
The variance $\vec{w}^\top \Sigma\, \vec{w}$ depends only on the risky weights. The expected return is linear in $\vec{w}$, while the variance is quadratic. This asymmetry is what makes the optimization tractable: we are maximizing a linear function subject to a quadratic penalty.
\end{remark}

To make the trade-off between return and risk precise, we need a preference specification.

\begin{definition}
A rational investor has an increasing, concave utility function $U$. The \textbf{portfolio choice problem} is
\begin{align}
    \max_{\vec{w}}\;\mathbb{E}\bigl[U(1 + R_p)\bigr].
\end{align}
For \textbf{power utility} with constant relative risk aversion $\gamma > 0$,
\begin{align}
    U(x) = \frac{x^{1-\gamma}}{1-\gamma},
\end{align}
a second-order Taylor expansion shows that maximizing expected utility is approximately equivalent to the \textbf{mean-variance problem}:
\begin{align}
    \max_{\vec{w}}\;\Bigl(\mathbb{E}[R_p] - \frac{\gamma}{2}\,\text{Var}(R_p)\Bigr).
\end{align}
\end{definition}

\begin{remark}
The parameter $\gamma$ controls the investor's willingness to bear risk. A pension fund with $\gamma \approx 5$--$10$ takes moderate positions; a leveraged hedge fund operating with $\gamma \approx 1$--$2$ takes aggressive ones. The mean-variance approximation is exact when returns are Gaussian or the utility is quadratic; otherwise it is a convenient approximation that becomes less accurate for large positions or heavy-tailed distributions.
\end{remark}

\begin{theorem}[Optimal portfolio]
The solution to the mean-variance problem is
\begin{align}
    \vec{w}^* = \frac{1}{\gamma}\,\Sigma^{-1}(\vec{\mu} - r_0\,\vec{1}_m),
\end{align}
with the remaining wealth $1 - \vec{w}^* \cdot \vec{1}_m$ invested in the risk-free asset.
\end{theorem}

\begin{proof}[Proof sketch]
The objective is $\vec{w} \cdot (\vec{\mu} - r_0\,\vec{1}_m) + r_0 - \frac{\gamma}{2}\vec{w}^\top\Sigma\,\vec{w}$, which is concave in $\vec{w}$. Taking the gradient and setting it to zero: $\vec{\mu} - r_0\,\vec{1}_m - \gamma\Sigma\,\vec{w} = 0$, which gives the stated formula.
\end{proof}

\begin{remark}
The optimal portfolio is proportional to $\Sigma^{-1}(\vec{\mu} - r_0\,\vec{1}_m)$: the inverse covariance matrix ``decorrelates'' the assets, and the excess return vector determines the direction. More risk-averse investors ($\gamma$ large) hold smaller positions in the same proportions.
\end{remark}

\begin{example}[Two-asset portfolio]
Consider two risky assets with $\mu_1 = 8\%$, $\mu_2 = 12\%$, $\sigma_1 = 15\%$, $\sigma_2 = 25\%$, $\rho_{12} = 0.3$, and $r_0 = 3\%$. The covariance matrix is
\begin{align}
    \Sigma = \begin{pmatrix} 0.0225 & 0.01125 \\ 0.01125 & 0.0625 \end{pmatrix}, \qquad
    \Sigma^{-1} = \begin{pmatrix} 51.28 & -9.23 \\ -9.23 & 18.46 \end{pmatrix}.
\end{align}
The excess return vector is $\vec{\mu} - r_0\vec{1} = (0.05, 0.09)^\top$. For $\gamma = 3$:
\begin{align}
    \vec{w}^* = \frac{1}{3}\begin{pmatrix} 51.28 \times 0.05 - 9.23 \times 0.09 \\ -9.23 \times 0.05 + 18.46 \times 0.09 \end{pmatrix} = \frac{1}{3}\begin{pmatrix} 1.733 \\ 1.200 \end{pmatrix} = \begin{pmatrix} 0.578 \\ 0.400 \end{pmatrix}.
\end{align}
The investor puts 57.8\% in asset 1, 40.0\% in asset 2, and $1 - 0.978 = 2.2\%$ in the risk-free asset. The portfolio is nearly fully invested. With $\gamma = 10$, the weights would be $(0.173, 0.120)$, with 70.7\% in the risk-free asset.
\end{example}

The set of all optimal portfolios, as $\gamma$ varies, traces a straight line in mean-standard-deviation space.

\begin{definition}
The \textbf{tangency portfolio} is the unique efficient portfolio fully invested in risky assets ($\vec{w} \cdot \vec{1}_m = 1$). It corresponds to risk aversion
\begin{align}
    \hat{\gamma} = \vec{1}_m^\top \Sigma^{-1}(\vec{\mu} - r_0\,\vec{1}_m),
\end{align}
with weights $\hat{\vec{w}} = \hat{\gamma}^{-1}\Sigma^{-1}(\vec{\mu} - r_0\,\vec{1}_m)$.
\end{definition}

\begin{theorem}[Two-fund separation]
Every optimal portfolio can be written as a combination of the risk-free asset and the tangency portfolio:
\begin{align}
    \vec{w} = \frac{\hat{\gamma}}{\gamma}\,\hat{\vec{w}}.
\end{align}
The fraction $\hat{\gamma}/\gamma$ invested in the tangency portfolio determines the risk level. In mean-standard-deviation space, all optimal portfolios lie on the \textbf{capital allocation line (CAL)}:
\begin{align}
    \mathbb{E}[R_p] = r_0 + \frac{\mathbb{E}[\hat{R}_p] - r_0}{\sqrt{\text{Var}(\hat{R}_p)}}\,\sqrt{\text{Var}(R_p)}.
\end{align}
The slope is the \textbf{Sharpe ratio} of the tangency portfolio.
\end{theorem}

\begin{remark}
Two-fund separation is a remarkable result: regardless of their risk aversion, all investors hold the \emph{same} tangency portfolio. They differ only in how much of it they hold (the rest is in the risk-free asset). This is the mathematical justification for index investing: if all investors agree on $\vec{\mu}$ and $\Sigma$, the tangency portfolio \emph{is} the market portfolio.

In practice, two-fund separation breaks down because investors disagree on expected returns, face different constraints (taxes, regulatory, ESG), and because $\vec{\mu}$ and $\Sigma$ are estimated with error. The sensitivity of $\vec{w}^* = \gamma^{-1}\Sigma^{-1}(\vec{\mu} - r_0\vec{1})$ to estimation error in $\vec{\mu}$ is severe, wherein small changes in expected returns produce large changes in optimal weights. This is the ``Markowitz curse'' and motivates the regularization and robust optimization techniques discussed in the numerical methods chapter.
\end{remark}


\subsection{Factor Models and Performance Evaluation}

Mean-variance analysis takes expected returns and covariances as given. Factor models explain \emph{where} those returns come from, decomposing asset risk into systematic (market-wide) and idiosyncratic (asset-specific) components.

\begin{theorem}[Capital Asset Pricing Model]
Assume all investors are mean-variance optimizers with homogeneous beliefs, a risk-free asset exists, and markets clear. Let $R_m$ be the return on the \textbf{market portfolio} (the value-weighted portfolio of all risky assets). Then, in equilibrium, the expected excess return on any asset $i$ satisfies
\begin{align}
    \mathbb{E}[R_i - r_0] = \beta_i\,\mathbb{E}[R_m - r_0], \qquad \beta_i = \frac{\text{Cov}(R_i, R_m)}{\text{Var}(R_m)}.
\end{align}
Only systematic risk (measured by $\beta_i$) is compensated; idiosyncratic risk earns no premium because it can be diversified away.
\end{theorem}

\begin{remark}
The CAPM is a single-factor model: the market portfolio is the sole risk factor. It explains the broad cross-section of returns (high-beta stocks earn more on average than low-beta stocks) but leaves significant unexplained variation. This motivates multi-factor extensions.
\end{remark}

\begin{definition}
A \textbf{multi-factor model} specifies
\begin{align}
    \mathbb{E}[R_i - r_0] = \sum_{j=1}^J \beta_{i,j}\,\mathbb{E}[F_j],
\end{align}
where $F_1, \ldots, F_J$ are systematic risk factors with $\beta_{i,j} = \text{Cov}(R_i, F_j)/\text{Var}(F_j)$.
\end{definition}

\begin{remark}
The most widely used multi-factor model in empirical finance is the \textbf{Fama-French three-factor model}, which adds size (SMB: small minus big) and value (HML: high book-to-market minus low) to the market factor. Extensions include momentum (Carhart four-factor), profitability and investment (Fama-French five-factor), and macro factors (term spread, credit spread, inflation). In industry, proprietary factor models (e.g., Barra, Axioma) use dozens of factors and are the basis for risk decomposition, performance attribution, and portfolio construction.
\end{remark}

Factor models provide the framework for evaluating portfolio performance.

\begin{definition}
Given a factor model, the time-series regression for fund $i$ is
\begin{align}
    R_{i,t} - r_{0,t} = \alpha_i + \sum_{j=1}^J \beta_{i,j}\,F_{j,t} + \varepsilon_{i,t}.
\end{align}
The intercept $\alpha_i$ (\textbf{alpha}) measures abnormal performance not explained by factor exposures---the manager's ``skill.'' The coefficients $\beta_{i,j}$ (\textbf{betas}) measure the portfolio's systematic exposures.
\end{definition}

\begin{definition}
The \textbf{information ratio} is
\begin{align}
    \text{IR}_i = \frac{\alpha_i}{\sigma(\varepsilon_{i,t})},
\end{align}
measuring risk-adjusted abnormal return per unit of idiosyncratic (non-factor) risk. An IR above 0.5 is considered good; above 1.0 is exceptional and rare over sustained periods.
\end{definition}

\begin{example}[Performance attribution]
A long/short equity fund reports annual excess return of 9\%. Running the Fama-French regression yields $\beta_{\text{MKT}} = 0.2$ (low market exposure), $\beta_{\text{SMB}} = 0.5$ (tilted toward small caps), $\beta_{\text{HML}} = -0.3$ (tilted toward growth), and $\alpha = 4.5\%$. If the factor premia over the period were $\text{MKT} = 8\%$, $\text{SMB} = 3\%$, $\text{HML} = 2\%$, the decomposition is:
\begin{align}
    9\% = \underbrace{4.5\%}_{\alpha} + \underbrace{0.2 \times 8\%}_{1.6\%} + \underbrace{0.5 \times 3\%}_{1.5\%} + \underbrace{(-0.3) \times 2\%}_{-0.6\%} + \underbrace{\varepsilon}_{2.0\%}.
\end{align}
Half of the return came from alpha (stock selection), with the rest from factor tilts and residual. An investor paying 2-and-20 fees for this fund is paying for the 4.5\% alpha, whereas the factor returns could be obtained cheaply via ETFs.
\end{example}


\section{Portfolio Risk}

Portfolio theory tells us how to allocate capital; risk management tells us how to \emph{control} the downside.

\subsection{Risk Measures}

\begin{definition}
Let $L$ be a random variable representing the \textbf{loss} of a portfolio over a fixed horizon, where positive values correspond to losses. The distribution of $L$ summarizes the portfolio's risk profile.
\end{definition}

\begin{definition}
For a confidence level $\alpha \in (0,1)$, the \textbf{Value at Risk (VaR)} at level $\alpha$ is the $\alpha$-quantile of the loss distribution:
\begin{align}
    \text{VaR}_\alpha(L) = \inf\bigl\{\ell \in \mathbb{R} : \mathbb{P}(L \leq \ell) \geq \alpha\bigr\}.
\end{align}
With probability $\alpha$, the loss does not exceed $\text{VaR}_\alpha$ over the given horizon.
\end{definition}

\begin{example}[VaR computation]
A portfolio has normally distributed daily P\&L with mean \$0 and standard deviation \$1M. The 99\% VaR is $\text{VaR}_{0.99} = 2.326 \times \$1\text{M} = \$2.326\text{M}$: with 99\% probability, the daily loss does not exceed \$2.33M. Equivalently, losses exceeding \$2.33M are expected about 2--3 times per year (1\% of 252 trading days $\approx 2.5$).
\end{example}

\begin{remark}
VaR is simple, intuitive, and widely mandated by regulators (Basel III/IV). However, it has a fundamental theoretical flaw: VaR is not \textbf{subadditive}. That is, it is possible to have $\text{VaR}(L_1 + L_2) > \text{VaR}(L_1) + \text{VaR}(L_2)$, wherein combining two portfolios can appear to \emph{increase} risk, even when diversification should reduce it. This pathology occurs for heavy-tailed distributions and makes VaR unsuitable as a basis for capital allocation across desks.
\end{remark}

Expected Shortfall fixes this problem.

\begin{definition}
The \textbf{Expected Shortfall (ES)}, or \textbf{Conditional Value at Risk (CVaR)}, at level $\alpha$ is
\begin{align}
    \text{ES}_\alpha(L) = \mathbb{E}\bigl[L \mid L \geq \text{VaR}_\alpha(L)\bigr].
\end{align}
It measures the expected loss in the worst $(1-\alpha)$ fraction of outcomes.
\end{definition}

\begin{remark}
For the Gaussian portfolio above, $\text{ES}_{0.99} = \frac{\phi(2.326)}{0.01} \times \$1\text{M} \approx 2.665 \times \$1\text{M} = \$2.665\text{M}$, where $\phi$ is the standard normal density. ES is always larger than VaR because it averages over the entire tail beyond the quantile, not just the boundary. 
\end{remark}

\begin{definition}
A risk measure $\rho$ is \textbf{coherent} if it satisfies:
\begin{enumerate}
    \item \emph{Monotonicity:} $L_1 \leq L_2 \implies \rho(L_1) \leq \rho(L_2)$.
    \item \emph{Translation invariance:} $\rho(L + c) = \rho(L) + c$.
    \item \emph{Positive homogeneity:} $\rho(\lambda L) = \lambda\,\rho(L)$ for $\lambda > 0$.
    \item \emph{Subadditivity:} $\rho(L_1 + L_2) \leq \rho(L_1) + \rho(L_2)$.
\end{enumerate}
\end{definition}

\begin{theorem}
Expected Shortfall is a coherent risk measure. Value at Risk is not (it violates subadditivity in general).
\end{theorem}

Beyond point-in-time measures, path-dependent risk is captured by drawdowns.

\begin{definition}
Let $W_t$ denote portfolio value. The \textbf{maximum drawdown} over $[0,T]$ is
\begin{align}
    \text{MDD} = \max_{0 \leq t \leq T}\!\Bigl(\sup_{0 \leq s \leq t} W_s - W_t\Bigr).
\end{align}
\end{definition}

\begin{remark}
Drawdowns measure the largest peak-to-trough decline and capture path-dependent risk that variance-based measures miss entirely. An investor experiencing a 30\% drawdown may be forced to liquidate (margin calls, redemptions) regardless of the portfolio's long-run Sharpe ratio. For this reason, drawdown constraints often dominate allocation decisions in practice: many hedge funds target a maximum drawdown of 10--15\%, which implicitly constrains leverage and concentration far more tightly than a VaR limit.
\end{remark}


\subsection{Diversification and Risk Decomposition}

\begin{proposition}[Diversification]
For a portfolio with weight vector $\vec{w}$ and covariance matrix $\Sigma$,
\begin{align}
    \text{Var}(R_p) = \vec{w}^\top \Sigma\, \vec{w} = \sum_i w_i^2 \sigma_i^2 + \sum_{i \neq j} w_i w_j \rho_{ij}\sigma_i\sigma_j.
\end{align}
When correlations are imperfect ($|\rho_{ij}| < 1$), the portfolio variance is strictly less than $(\sum_i w_i\sigma_i)^2$.
\end{proposition}

\begin{proof}
The expansion follows from the definition of the quadratic form. The inequality $\vec{w}^\top\Sigma\,\vec{w} \leq (\sum_i w_i\sigma_i)^2$ holds because the cross terms $w_iw_j\rho_{ij}\sigma_i\sigma_j \leq w_iw_j\sigma_i\sigma_j$ when $\rho_{ij} \leq 1$, with strict inequality when $|\rho_{ij}| < 1$ for at least one pair.
\end{proof}

\begin{example}[Diversification benefit]
An equal-weight portfolio of $n$ identical assets, each with variance $\sigma^2$ and pairwise correlation $\rho$, has variance
\begin{align}
    \text{Var}(R_p) = \frac{\sigma^2}{n} + \frac{n-1}{n}\rho\sigma^2.
\end{align}
As $n \to \infty$, the variance converges to $\rho\sigma^2$ (the undiversifiable systematic component). For $\sigma = 30\%$ and $\rho = 0.3$: a single asset has variance $9\%$, while a 100-asset portfolio has variance $\approx 2.79\%$ (a 69\% reduction). But the floor $\rho\sigma^2 = 2.70\%$ cannot be breached no matter how many assets are added. Diversification eliminates \textbf{idiosyncratic risk} but not \textbf{systematic risk}.
\end{example}

Understanding \emph{where} portfolio risk comes from is essential for management.

\begin{definition}
The \textbf{marginal risk contribution} of asset $i$ is
\begin{align}
    \text{MRC}_i = \frac{\partial}{\partial w_i}\text{Var}(R_p) = 2(\Sigma\,\vec{w})_i.
\end{align}
The \textbf{total risk contribution} of asset $i$ is
\begin{align}
    \text{RC}_i = w_i\,(\Sigma\,\vec{w})_i.
\end{align}
These satisfy $\sum_i \text{RC}_i = \vec{w}^\top\Sigma\,\vec{w} = \text{Var}(R_p)$: risk contributions decompose total portfolio variance additively.
\end{definition}

\begin{remark}
A \textbf{risk parity} portfolio equalizes risk contributions: $\text{RC}_i = \text{RC}_j$ for all $i,j$. This ensures that no single asset dominates the portfolio's risk profile. In a traditional 60/40 stock/bond portfolio, equities contribute roughly 90\% of the total risk despite being only 60\% of the weight, because equity volatility is much higher than bond volatility. Risk parity would increase the bond allocation (often with leverage) until the contributions balance (as in Bridgewater's All Weather fund).
\end{remark}


\subsection{Hedging and Risk Limits}

\begin{definition}
A \textbf{hedge} is a position constructed to reduce sensitivity to a specified risk factor. If a portfolio has factor exposure $\beta_F$ to factor $F$, a hedge with exposure $\beta_h$ satisfying
\begin{align}
    \beta_F + \beta_h = 0
\end{align}
eliminates first-order exposure to $F$.
\end{definition}

\begin{example}[Delta hedging an equity portfolio]
A portfolio has $\beta_{\text{MKT}} = 1.3$ (30\% more market-sensitive than the benchmark). To hedge, sell S\&P 500 futures with notional equal to $1.3 \times$ the portfolio value. The hedged portfolio has $\beta_{\text{MKT}} \approx 0$ and retains only idiosyncratic and non-market-factor risk, such that the alpha becomes the primary source of return.
\end{example}

\begin{remark}
In practice, hedging reduces expected variance but introduces \textbf{basis risk} (the hedge instrument does not perfectly track the exposure), \textbf{transaction costs} (bid-ask spreads on the hedge), \textbf{roll costs} (if using futures that expire), and \textbf{model risk} (the estimated $\beta$ may be wrong). Perfect hedging is an idealization; practical hedging is a trade-off between residual risk and hedging cost.
\end{remark}

\begin{definition}
A \textbf{risk limit} constrains a portfolio's risk exposure:
\begin{align}
    \text{VaR}_\alpha(L) \leq \ell_{\max}, \qquad \text{ES}_\alpha(L) \leq e_{\max},
\end{align}
or constraints on factor exposures (e.g., $|\beta_{\text{MKT}}| \leq 0.5$), leverage ($\sum_i|w_i| \leq L_{\max}$), concentration ($|w_i| \leq c_{\max}$), or drawdown.
\end{definition}

\begin{proposition}
If the risk measure $\rho$ is convex (which ES is), then the set $\{\vec{w} : \rho(\vec{w}^\top\vec{R}) \leq c\}$ is convex. Imposing convex risk constraints on the mean-variance problem preserves the convexity of the optimization.
\end{proposition}

\begin{remark}
This is why ES-based constraints are preferable to VaR-based constraints from an optimization standpoint: VaR constraints can produce non-convex feasible sets, leading to multiple local optima and computational difficulties.
\end{remark}



\chapter{Markets and Trading}

The models of the preceding chapters operate in a frictionless world: assets trade continuously, at no cost, in unlimited quantities, at a single well-defined price. Real markets are nothing like this. Trades execute through limit order books with bid-ask spreads. Large orders move prices. Information is asymmetric. Transaction costs erode returns.

\textbf{Market microstructure} studies how these frictions affect price formation, trading behavior, and market quality. This chapter introduces the institutional structure of modern markets (exchanges, order books, market makers), then develops the theory of optimal execution.


\section{Market Structure}

\subsection{Venues and Participants}

\begin{definition}
A \textbf{securities exchange} is a regulated venue where financial instruments are listed and traded. Examples include the NYSE and NASDAQ (equities, options, ETFs), the CME (futures), and ICE (commodities, credit). Securities that are not exchange-traded (forwards and  swaps) trade \textbf{over-the-counter (OTC)}, directly between counterparties or through dealer networks.
\end{definition}

\begin{remark}
Exchange-traded products benefit from central clearing (eliminating bilateral counterparty risk), standardized contracts (enhancing liquidity), and price transparency (all participants see the same order book). OTC products offer customization but require bilateral credit management (hence CVA/XVA, as developed in the preceding chapter). Post-2008 regulation (Dodd-Frank, EMIR) has pushed many OTC derivatives toward central clearing and exchange-like execution.
\end{remark}

\begin{definition}
A \textbf{broker} is an intermediary that executes trades on behalf of clients, earning a commission or markup. A \textbf{market maker} is a firm that continuously quotes bid and ask prices, standing ready to buy or sell. Market makers provide liquidity and profit from the bid-ask spread, which compensates them for \textbf{adverse selection} (the risk of trading against better-informed counterparties) and \textbf{inventory risk} (the risk of holding unwanted positions).
\end{definition}

\begin{remark}
Modern electronic market making is a high-frequency, technology-intensive business. Firms like Citadel Securities, Jane Street, and Optiver use sophisticated models to quote prices, manage inventory, and hedge risk across thousands of instruments simultaneously. The bid-ask spread has narrowed dramatically since the advent of electronic trading (from 12.5 cents per share on the NYSE in the 1990s to often less than 1 cent today).
\end{remark}


\subsection{The Limit Order Book}

The central mechanism of modern electronic markets is the limit order book, which aggregates all outstanding buy and sell orders for a given security.

\begin{definition}
A \textbf{central limit order book (CLOB)} records outstanding orders at each price level. At any time $t$, the book consists of:
\begin{enumerate}
    \item \textbf{Bid} (buy) orders at prices $p_1^b > p_2^b > \cdots$ with volumes $q_1^b, q_2^b, \ldots$
    \item \textbf{Ask} (sell) orders at prices $p_1^a < p_2^a < \cdots$ with volumes $q_1^a, q_2^a, \ldots$
\end{enumerate}
The \textbf{best bid} $p_t^b$ and \textbf{best ask} $p_t^a$ define the \textbf{bid-ask spread} $s_t = p_t^a - p_t^b$ and the \textbf{midprice} $m_t = (p_t^a + p_t^b)/2$.
\end{definition}

\begin{remark}
The midprice is commonly used as a proxy for the ``efficient'' or ``true'' price, but this is an approximation. When the order book is asymmetric (more volume on the bid side than the ask side), the efficient price is shifted toward the heavier side (the \textbf{microprice}).
\end{remark}

\begin{definition}
The three fundamental order types are:
\begin{enumerate}
    \item \textbf{Market orders}: execute immediately at the best available price. They \emph{consume} liquidity.
    \item \textbf{Limit orders}: rest in the book until matched or cancelled. They \emph{provide} liquidity.
    \item \textbf{Cancel orders}: remove previously submitted limit orders.
\end{enumerate}
\end{definition}

\begin{remark}
The interaction between market orders (demanding immediacy) and limit orders (providing it) drives price dynamics at short time scales. A market buy order ``lifts'' the best ask; if the ask-side volume is depleted, the price ticks up. Conversely, a sequence of limit buy orders below the best bid deepens the book without moving the price.
\end{remark}

\begin{definition}
\textbf{Adverse selection} in the context of market making arises because some counterparties (``informed traders'') possess private information about the asset's value. A market maker who quotes a tight spread risks being ``picked off'' by informed traders, earning the spread on uninformed flow but losing more on informed trades. The bid-ask spread must be wide enough to compensate for this expected loss.
\end{definition}

\begin{remark}
The Glosten-Milgrom (1985) and Kyle (1985) models formalize this intuition. In Glosten-Milgrom, the spread arises endogenously from Bayesian updating by the market maker, who cannot distinguish informed from uninformed orders. In Kyle, a single informed trader optimizes against a market maker and noise traders, and the equilibrium price impact is linear in order flow.
\end{remark}

\begin{definition}
\textbf{Market liquidity} is the ease with which an asset can be traded without substantially affecting its price. Liquid markets have tight spreads, deep order books, and low price impact. Illiquid markets have wide spreads, thin books, and high impact.
\end{definition}

\begin{definition}
In contrast to a CLOB, a \textbf{Request for Quote (RFQ)} model operates when a trader solicits prices from multiple dealers. The trader sees competing quotes and executes with the best one. RFQ dominates in OTC markets (corporate bonds, structured products, FX options) where order book liquidity is insufficient for continuous matching.
\end{definition}


\section{Optimal Execution}

A portfolio manager who decides to buy 500,000 shares of a mid-cap stock cannot simply submit a market order. The order would sweep through multiple price levels of the order book, moving the price against her by potentially several percent. Optimal execution theory addresses how to break up a large order over time to minimize the combination of market impact costs and price risk.

\subsection{The Almgren-Chriss Framework}

\begin{definition}[Execution problem]
An agent must execute a total quantity $Q$ of shares over a fixed horizon $[0,T]$, choosing a trading rate $q_t$ (shares per unit time). The \textbf{inventory process} is
\begin{align}
    dX_t = -q_t\,dt, \qquad X_0 = Q, \quad X_T = 0.
\end{align}
\end{definition}

\begin{definition}[Almgren-Chriss model]
The unaffected midprice follows a martingale:
\begin{align}
    dS_t = \sigma\,dW_t.
\end{align}
Trading at rate $q_t$ affects the execution price through two channels:
\begin{enumerate}
    \item \textbf{Temporary impact}: a cost $\eta\,q_t$ per share that affects only the current trade, reflecting the immediate liquidity cost of demanding execution.
    \item \textbf{Permanent impact}: a cumulative shift $\gamma\int_0^t q_s\,ds$ that moves all future prices, reflecting the information content of the trade.
\end{enumerate}
The execution price is $\widetilde{S}_t = S_t - \eta\,q_t - \gamma\int_0^t q_s\,ds$.
\end{definition}

\begin{remark}
Empirically, temporary impact is better described by a concave (square-root) function of trading rate, and permanent impact depends on the asset's liquidity and the information content of the trade. Nevertheless, the linear model is analytically tractable and captures the essential trade-off: trading fast incurs high impact costs; trading slowly exposes the portfolio to price risk.
\end{remark}

\begin{proposition}[Cost decomposition]
The total execution cost $C = \int_0^T q_t\,\widetilde{S}_t\,dt - QS_0$ decomposes as
\begin{align}
    C = \underbrace{\eta\int_0^T q_t^2\,dt}_{\text{temporary impact}} + \underbrace{\frac{\gamma}{2}Q^2}_{\text{permanent impact}} - \underbrace{\sigma\int_0^T X_t\,dW_t}_{\text{price risk}}.
\end{align}
\end{proposition}

The optimal execution problem is thus a trade-off between temporary impact (favoring slow execution) and price risk (favoring fast execution).

\begin{proposition}
An agent minimizing the mean-variance criterion $\mathbb{E}[C] + \lambda\,\text{Var}(C)$ faces the reduced problem
\begin{align}
    \min_{(X_t)}\;\int_0^T\bigl(\eta\,\dot{X}_t^2 + \lambda\sigma^2 X_t^2\bigr)\,dt, \qquad X_0 = Q,\;\; X_T = 0,
\end{align}
where $\lambda \geq 0$ is the risk aversion parameter.
\end{proposition}

\begin{theorem}[Optimal execution trajectory]
The optimal inventory path satisfies the Euler-Lagrange equation
\begin{align}
    \ddot{X}_t = \frac{\lambda\sigma^2}{\eta}\,X_t \eqqcolon \kappa^2 X_t,
\end{align}
with unique solution
\begin{align}
    X_t = Q\,\frac{\sinh\bigl(\kappa(T-t)\bigr)}{\sinh(\kappa T)}, \qquad \kappa = \sqrt{\frac{\lambda\sigma^2}{\eta}}.
\end{align}
\end{theorem}

The optimal trading rate follows by differentiation.

\begin{corollary}
The optimal trading rate is
\begin{align}
    q_t = Q\,\frac{\kappa\cosh\bigl(\kappa(T-t)\bigr)}{\sinh(\kappa T)}.
\end{align}
\end{corollary}

\begin{remark}
The solution interpolates between two extremes controlled by the ``urgency'' parameter $\kappa = \sqrt{\lambda\sigma^2/\eta}$:

\emph{Low urgency ($\kappa \to 0$, i.e., $\lambda$ small or $\eta$ large):} The agent trades at a constant rate $Q/T$, which is the \textbf{TWAP (time-weighted average price)} strategy. This minimizes expected impact cost but maximizes exposure to price risk.

\emph{High urgency ($\kappa \to \infty$, i.e., $\lambda$ large or $\eta$ small):} The agent front-loads execution to minimize exposure. In the limit, execution is instantaneous.

In practice, $\kappa$ is estimated from the stock's volatility $\sigma$ (readily available) and the impact parameter $\eta$ (estimated from historical execution data or market impact models). A typical mid-cap stock might have $\eta \approx 10^{-3}$\$/share per share/second and $\sigma \approx 30\%$ annualized.
\end{remark}

\begin{example}[Optimal execution for a block trade]
An asset manager needs to sell $Q = 200{,}000$ shares over $T = 1$ day (6.5 hours = 23{,}400 seconds) of a stock with $\sigma = 2\%$ daily, $\eta = 5 \times 10^{-6}$ \$/share per share/second, and risk aversion $\lambda = 10^{-6}$.

The urgency parameter is $\kappa = \sqrt{10^{-6} \times 0.02^2 / (5 \times 10^{-6})} = \sqrt{8 \times 10^{-5}} \approx 0.00894 \text{ s}^{-1}$, so $\kappa T = 0.00894 \times 23{,}400 \approx 209$.

Since $\kappa T \gg 1$, $\sinh(\kappa T) \approx \frac{1}{2}e^{\kappa T}$, and the trajectory becomes approximately $X_t \approx Q\,e^{-\kappa t}$---exponential front-loading. The agent sells roughly 63\% of the position in the first $1/\kappa \approx 112$ seconds (under 2 minutes), with the remainder sold gradually over the rest of the day.

If instead $\lambda = 0$ (risk-neutral), the optimal strategy is TWAP: sell $200{,}000/23{,}400 \approx 8.5$ shares per second uniformly. The expected cost is the same, but the variance is much higher.
\end{example}



\chapter{Time Series Analysis}

The continuous-time models of the preceding chapters, like geometric Brownian motion, stochastic volatility, and jump diffusions, are idealized descriptions of how prices evolve. But an analyst does not observe a continuous diffusion; they observes a discrete sequence of closing prices, tick-by-tick trades, or daily returns. Time series analysis is the discipline of extracting structure from such sequences.

The goals are different from, but complementary to, the arbitrage pricing framework. There, the question was ``what is a derivative worth under no-arbitrage?'' Here, the questions are empirical: Does a return series have predictable structure? Is volatility persistent, and if so, how should we model that persistence? Are two price series bound together by a long-run equilibrium, and can we profit from temporary deviations? The answers inform everything from risk measurement (VaR and ES require a model of the return distribution) to strategy construction (pairs trading requires cointegration) to model calibration (GARCH volatility feeds into option pricing).

This chapter develops the theory in three stages. First, we establish the foundations: stationarity, the Wold decomposition, and the ARMA/ARIMA family of linear models that describe serial dependence in levels. Second, we address the central empirical fact that financial returns exhibit time-varying volatility, building from ARCH through the GARCH family and its asymmetric and long-memory extensions. Third, we introduce spectral methods and cointegration, which reveal structure that is invisible in the time domain and formalize the notion that nonstationary series can share a common stochastic trend.


\section{Time Series Foundations}

\subsection{Stationarity and Dependence}

A time series model is useful only if its statistical properties are stable enough to learn from data. The concept of stationarity makes this precise.

\begin{definition}
A discrete-time stochastic process $\{X_t\}_{t \in \mathbb{Z}}$ is called a \textbf{time series}. The series is \textbf{strictly stationary} if all finite-dimensional distributions are invariant under time shifts: for all $m \in \mathbb{N}$ and all lags $\tau$,
\begin{align}
    (X_{t_1+\tau}, \ldots, X_{t_m+\tau}) \overset{d}{=} (X_{t_1}, \ldots, X_{t_m}).
\end{align}
\end{definition}

\begin{remark}
Strict stationarity is a strong condition, requiring \emph{every} aspect of the joint distribution to be time-invariant.
\end{remark}

\begin{definition}
The series $\{X_t\}$ is \textbf{weakly stationary} (or \textbf{covariance stationary}) if its first two moments exist and are time-invariant:
\begin{align}
    \mathbb{E}[X_t] = \mu, \quad
    \text{Var}(X_t) = \sigma^2 < \infty, \quad
    \text{Cov}(X_t, X_{t+\tau}) = \Gamma(\tau),
\end{align}
where $\mu$, $\sigma^2$, and the \textbf{autocovariance function} $\Gamma(\tau)$ do not depend on $t$. The \textbf{autocorrelation function (ACF)} is
\begin{align}
    \rho(\tau) = \frac{\Gamma(\tau)}{\Gamma(0)}.
\end{align}
\end{definition}

\begin{remark}
For Gaussian processes, weak and strict stationarity coincide (since a Gaussian distribution is determined by its first two moments). For financial returns, which are non-Gaussian, the distinction matters: a return series can be weakly stationary while its higher moments evolve over time, like in volatility clustering.
\end{remark}

\begin{definition}
A weakly stationary process $\{X_t\}$ is called \textbf{ergodic} if time averages converge to ensemble averages:
\begin{align}
    \lim_{T \to \infty} \frac{1}{T} \sum_{k=1}^{T} X_{t+k} = \mu \quad \text{a.s.}
\end{align}
\end{definition}

\begin{remark}
Ergodicity is the condition that makes statistical estimation possible from a single realization. We observe one path of a stock price, not an ensemble of parallel universes; ergodicity guarantees that the sample mean converges to the population mean. Most standard financial time series models are ergodic when stationary, but the assumption can fail for regime-switching processes observed over short windows.
\end{remark}

Two operators provide compact notation for the models that follow.

\begin{definition}
The \textbf{lag operator} $L$ and \textbf{difference operator} $\Delta$ are defined by
\begin{align}
    L^k X_t = X_{t-k}, \qquad \Delta^k = (1-L)^k.
\end{align}
In particular, $\Delta X_t = X_t - X_{t-1}$ is the first difference (the return, in levels terms).
\end{definition}


\subsection{The Wold Decomposition}

Before introducing specific model classes, we establish a foundational result: every stationary process has a canonical representation as a linear filter of white noise.

\begin{theorem}[Wold representation]
Any zero-mean weakly stationary process $\{X_t\}$ admits the decomposition
\begin{align}
    X_t = V_t + \sum_{i=0}^{\infty} \psi_i \varepsilon_{t-i} = V_t + \psi(L)\varepsilon_t,
\end{align}
where $\{\varepsilon_t\} \sim \text{WN}(0, \sigma^2)$ is white noise (uncorrelated, zero-mean, constant-variance), $\psi_0 = 1$, $\sum_i \psi_i^2 < \infty$, and $V_t$ is a deterministic (perfectly predictable) component.
\end{theorem}

\begin{remark}
The Wold theorem says that after stripping away any deterministic trend, a stationary process is an infinite moving average of unpredictable shocks. The coefficients $\{\psi_i\}$ describe how the process ``remembers'' past shocks - they are the impulse response function. A shock $\varepsilon_{t-k}$ contributes $\psi_k \varepsilon_{t-k}$ to the current value. If the $\psi_i$ decay quickly, the process has short memory; if they decay slowly, the process has long memory. The ARMA models below are parametric approximations to this infinite representation.
\end{remark}

The Wold representation expresses the process as an MA($\infty$). Under certain conditions, the representation can be inverted.

\begin{definition}
If the operator $\psi(L)$ has a convergent inverse $\psi^{-1}(L) = \sum_{i=0}^{\infty} \psi_i^* L^i$, then the moving average representation is \textbf{invertible}, and admits an equivalent autoregressive representation:
\begin{align}
    X_t = \psi(L)\varepsilon_t \quad \Longleftrightarrow \quad \psi^{-1}(L) X_t = \varepsilon_t.
\end{align}
\end{definition}

\begin{remark}
Invertibility means that the white noise innovations $\varepsilon_t$ can be recovered from the observed process, which is necessary for estimation, since we need to compute residuals from data. Operationally, invertibility requires that the roots of $\psi(z)$ lie outside the unit circle, a condition we will encounter repeatedly.
\end{remark}


\subsection{ARMA Models}

The ARMA family provides parsimonious parametric approximations to the Wold representation. We build up from the two simplest building blocks.

\begin{definition}
An \textbf{MA($q$)} (moving average of order $q$) model is
\begin{align}
    X_t - \mu = \theta(L)\varepsilon_t, \qquad \theta(z) = 1 + \theta_1 z + \cdots + \theta_q z^q,
\end{align}
where $\{\varepsilon_t\} \sim \text{WN}(0, \sigma^2)$.
\end{definition}

\begin{proposition}
MA($q$) processes are always weakly stationary.
\end{proposition}

\begin{proof}[Proof sketch]
$\mathbb{E}[X_t] = \mu$ since $\mathbb{E}[\varepsilon_s] = 0$. The autocovariance $\Gamma(\tau) = \text{Cov}(X_t, X_{t+\tau})$ involves products $\mathbb{E}[\varepsilon_s \varepsilon_r]$, which equal $\sigma^2$ when $s=r$ and zero otherwise. The number of ``overlapping'' noise terms depends only on $|\tau|$, not on $t$, so $\Gamma(\tau)$ is time-invariant. In particular, $\Gamma(\tau) = 0$ for $|\tau| > q$: an MA($q$) process has a finite memory horizon.
\end{proof}

\begin{remark}
The sharp cutoff in the ACF at lag $q$ is the diagnostic signature of an MA process and is used in practice to identify the order $q$. An MA(1) with $\theta_1 = 0.5$ has $\rho(1) = \theta_1/(1+\theta_1^2) = 0.4$ and $\rho(\tau) = 0$ for $\tau \geq 2$.
\end{remark}

\begin{definition}
An \textbf{AR($p$)} (autoregressive of order $p$) model is
\begin{align}
    \phi(L)(X_t - \mu) = \varepsilon_t, \qquad \phi(z) = 1 - \phi_1 z - \cdots - \phi_p z^p.
\end{align}
\end{definition}

\begin{remark}
An AR process is causal: the current value depends linearly on its own past values plus a contemporaneous shock. This makes economic sense and makes estimation straightforward via regression.
\end{remark}

The stationarity of an AR process is not automatic; it depends on the roots of the characteristic polynomial.

\begin{proposition}
An AR($p$) process is weakly stationary if and only if all roots of $\phi(z) = 0$ lie strictly outside the unit circle.
\end{proposition}

\begin{proof}
Stationarity is equivalent to the existence of a convergent Wold (MA($\infty$)) representation $X_t = \phi(L)^{-1}\varepsilon_t$. The operator $\phi(L)^{-1}$ can be expanded as a convergent power series in $L$ if and only if $1/\phi(z)$ has no poles within the unit disk. By complex analysis, this requires all roots of $\phi(z)$ to lie strictly outside the unit circle.
\end{proof}

\begin{example}[AR(1) mean-reversion]
The AR(1) process $X_t = c + \phi X_{t-1} + \varepsilon_t$ is the discrete-time analogue of the Ornstein-Uhlenbeck process. The characteristic polynomial $\phi(z) = 1 - \phi z$ has root $z = 1/\phi$, so stationarity requires $|\phi| < 1$.

When stationary, iterating backward gives
\begin{align}
    X_t = \frac{c}{1-\phi} + \sum_{k=0}^{\infty} \phi^k \varepsilon_{t-k},
\end{align}
with unconditional mean $\mu = c/(1-\phi)$ and autocovariance $\Gamma(\tau) = \sigma^2 \phi^{|\tau|}/(1-\phi^2)$, so the ACF decays exponentially: $\rho(\tau) = \phi^{|\tau|}$.

\textbf{Financial interpretation.} If $|\phi|$ is close to 1 but strictly less, shocks are highly persistent but ultimately mean-revert. Daily interest rate changes and credit spreads are often well-described by AR(1) models with $\phi \approx 0.97$--$0.99$. The half-life of a shock is $\log(0.5)/\log(|\phi|)$: for $\phi = 0.98$, this is about 34 periods, meaning a deviation from the mean takes roughly 34 days to halve.
\end{example}

The boundary case $\phi = 1$ is of special importance in finance.

\begin{definition}
An AR($p$) process possesses a \textbf{unit root} if $\phi(z) = 0$ has at least one root with $|z_0| = 1$.
\end{definition}

\begin{example}[Random walk]
Setting $\phi = 1$ in AR(1) yields $X_t = X_{t-1} + \varepsilon_t$, a random walk. Iterating gives $X_t = X_0 + \sum_{j=1}^t \varepsilon_j$, so $\mathbb{E}[X_t] = X_0$ but $\text{Var}(X_t) = t\sigma^2$ - the variance grows without bound. The process is not stationary: shocks accumulate permanently rather than decaying.

\textbf{Financial interpretation.} The efficient market hypothesis in its simplest form states that log-prices follow a random walk: $\log S_t = \log S_{t-1} + \varepsilon_t$. If true, returns $\Delta \log S_t = \varepsilon_t$ are unpredictable, and the price level is non-stationary. The distinction between a stationary AR(1) with $\phi = 0.999$ and a true unit root $\phi = 1$ is economically profound but statistically difficult to detect in finite samples. This is the unit root testing problem.
\end{example}

\begin{theorem}[Dickey-Fuller test]
Consider the AR(1) process rewritten in differences:
\begin{align}
    \Delta X_t = \delta X_{t-1} + \varepsilon_t, \qquad \delta = \phi - 1.
\end{align}
The null hypothesis is $\delta = 0$ (unit root, non-stationary) against the alternative $\delta < 0$ (stationary). The OLS estimator $\hat{\delta}$ can be computed, but under the null, the $t$-statistic does not follow a standard $t$-distribution; it follows the \textbf{Dickey-Fuller distribution}, which has heavier left tails. Critical values must be obtained from simulation or tabulated values.
\end{theorem}

\begin{remark}
The non-standard distribution arises because under the null, $X_{t-1}$ is integrated (non-stationary), violating the regularity conditions for classical OLS asymptotics. The \textbf{Augmented Dickey-Fuller (ADF)} test extends this by including lagged differences $\Delta X_{t-1}, \ldots, \Delta X_{t-k}$ to control for higher-order serial correlation. The \textbf{Phillips-Perron} test uses a non-parametric correction instead. The \textbf{KPSS} test reverses the hypotheses, taking stationarity as the null. Running both ADF and KPSS provides a more robust diagnostic.

Unit root tests are notoriously low-powered: in moderate samples (a few years of daily data), they struggle to distinguish $\phi = 0.99$ from $\phi = 1$. This is a persistent headache in empirical finance.
\end{remark}

With the AR and MA building blocks established, we combine them.

\begin{definition}
A process $\{X_t\}$ follows an \textbf{ARMA($p,q$)} model if
\begin{align}
    \phi(L)(X_t - \mu) = \theta(L)\varepsilon_t,
\end{align}
where $\phi(z) = 1 - \phi_1 z - \cdots - \phi_p z^p$ and $\theta(z) = 1 + \theta_1 z + \cdots + \theta_q z^q$. The Wold representation is
\begin{align}
    X_t = \mu + \psi(L)\varepsilon_t, \qquad \psi(L) = \phi(L)^{-1}\theta(L),
\end{align}
provided the process is stationary (roots of $\phi$ outside the unit circle) and invertible (roots of $\theta$ outside the unit circle).
\end{definition}

\begin{remark}
ARMA models are the workhorse of linear time series analysis. The AR component captures persistence (current values depend on past values), and the MA component captures the shape of the impulse response at short lags. The combination is parsimonious: an ARMA(1,1) with two parameters can approximate AR processes of very high order. In finance, ARMA models are applied to interest rates, exchange rates, and macroeconomic indicators. For equity returns, which are approximately uncorrelated, ARMA models add little; the action is in the \emph{volatility}, which we address in Section~\ref{sec:volatility}.
\end{remark}

\begin{example}[ARMA(1,1) for credit spreads]
Consider the model $X_t = 0.95\, X_{t-1} + \varepsilon_t + 0.3\, \varepsilon_{t-1}$ for the daily change in a CDS spread. The AR coefficient $\phi = 0.95$ captures the persistence of credit conditions; wide spreads tend to stay wide. The MA coefficient $\theta = 0.3$ captures the fact that a liquidity shock on day $t-1$ continues to affect spreads on day $t$, beyond what the AR component predicts. The unconditional variance is
\begin{align}
    \text{Var}(X_t) = \sigma^2 \cdot \frac{1 + 2\phi\theta + \theta^2}{1 - \phi^2} = \sigma^2 \cdot \frac{1 + 0.57 + 0.09}{1 - 0.9025} \approx 17.0\, \sigma^2.
\end{align}

\end{example}

For non-stationary series, differencing removes stochastic trends.

\begin{definition}
A process $\{X_t\}$ is \textbf{integrated of order $d$}, written $I(d)$, if $\Delta^d X_t$ is stationary but $\Delta^{d-1} X_t$ is not. An $I(1)$ process, the most common case in finance, has a unit root; its first difference $\Delta X_t = X_t - X_{t-1}$ is stationary.
\end{definition}

\begin{definition}
If $\Delta^d X_t$ follows an ARMA($p,q$) model, then $X_t$ follows an \textbf{ARIMA($p,d,q$)} model:
\begin{align}
    \phi(L)\Delta^d X_t = \theta(L)\varepsilon_t.
\end{align}
A random walk is ARIMA(0,1,0); a random walk with drift $c$ adds $c$ to the right-hand side.
\end{definition}

\begin{remark}
The ARIMA framework codifies a simple recipe: if a series is non-stationary, difference it until it becomes stationary, then fit an ARMA model to the differences. For log-prices, taking $d=1$ gives log-returns, which are typically stationary. The Box-Jenkins methodology - identify $d$ (via unit root tests), choose $(p,q)$ (via ACF/PACF diagnostics and information criteria), estimate parameters (via MLE), and validate (via residual diagnostics) - remains the standard workflow.
\end{remark}


\subsection{Estimation and Model Selection}

Fitting a model requires choosing \emph{which} model (order selection) and then estimating its parameters. Both steps require care.

\begin{definition}
For an ARMA($p,q$) process with Gaussian innovations, the parameter vector is $\theta = (\phi_1, \ldots, \phi_p, \theta_1, \ldots, \theta_q, \sigma^2)$. The \textbf{maximum likelihood estimator (MLE)} maximizes the log-likelihood
\begin{align}
    \ell(\theta) = \log \mathbb{P}(X_1, \ldots, X_T \mid \theta),
\end{align}
which is typically computed via the Kalman filter (exploiting the state-space representation of ARMA models) or by exact likelihood recursions. MLE is asymptotically efficient under correct specification.
\end{definition}

For pure autoregressive models, a simpler estimator is available.

\begin{definition}
The \textbf{Yule-Walker equations} relate the AR coefficients to the theoretical autocovariances:
\begin{align}
    \Gamma(\tau) = \sum_{j=1}^p \phi_j\, \Gamma(\tau - j), \qquad \tau \geq 1.
\end{align}
In matrix form, $\boldsymbol{\Gamma} \boldsymbol{\phi} = \boldsymbol{\gamma}$, where $\boldsymbol{\Gamma}$ is the Toeplitz autocovariance matrix. Replacing population autocovariances with sample counterparts yields the \textbf{Yule-Walker estimator} $\hat{\boldsymbol{\phi}} = \hat{\boldsymbol{\Gamma}}^{-1}\hat{\boldsymbol{\gamma}}$.
\end{definition}

\begin{remark}
Yule-Walker estimators are closed-form and computationally cheap, but less efficient than MLE in finite samples, especially when MA terms are present. In practice, they serve as initial values for the numerical optimization in MLE.
\end{remark}

Choosing the right model order requires balancing fit against complexity. Overly complex models fit noise; overly simple models miss structure.

\begin{definition}
Let $\hat{\sigma}^2$ denote the estimated innovation variance and $n$ the sample size. The \textbf{Akaike Information Criterion (AIC)} and \textbf{Bayesian Information Criterion (BIC)} are
\begin{align}
    \text{AIC}(p,q) &= \log\hat{\sigma}^2 + \frac{2(p+q)}{n}, \\
    \text{BIC}(p,q) &= \log\hat{\sigma}^2 + \frac{\log(n)\,(p+q)}{n}.
\end{align}
The preferred model minimizes the criterion.
\end{definition}

\begin{remark}
The AIC penalizes complexity at rate $2k/n$ and is motivated by minimizing the Kullback-Leibler divergence between the true and fitted models. It tends to select larger models and is better suited for \emph{forecasting}. The BIC penalizes at rate $(\log n)k/n$, which grows with sample size and is therefore more conservative. BIC is \emph{consistent}: it identifies the true model with probability approaching 1 as $n \to \infty$, assuming the true model is among the candidates. AIC is not consistent but adapts better when the true data-generating process is not in the candidate set.
\end{remark}


\section{Volatility Modeling}
\label{sec:volatility}

The ARMA framework models the conditional mean of a time series. For equity returns, however, the conditional mean is nearly zero, so returns are approximately unpredictable. The action is in the \emph{conditional variance}. This section develops models that capture the systematic dynamics of volatility, the most important statistical feature of financial return series.

\subsection{Stylized Facts}

Before writing down models, we catalog the empirical regularities that any volatility model must capture. These ``stylized facts'' are remarkably robust across asset classes, time periods, and frequencies.

\begin{remark}[Volatility clustering]
Large price movements tend to be followed by large movements of either sign, and small movements by small movements. Formally, while raw returns $\{X_t\}$ exhibit little serial correlation ($\rho_X(\tau) \approx 0$), the squared and absolute returns $\{X_t^2\}$ and $\{|X_t|\}$ display strong, slowly decaying positive autocorrelation. This phenomenon of \textbf{volatility clustering} was first documented by Mandelbrot (1963) and is visible to the naked eye in any return series plot: periods of calm alternate with periods of turbulence.
\end{remark}

\begin{remark}[Heavy tails]
Empirical return distributions have heavier tails than the Gaussian. The kurtosis of daily equity returns is typically 5--10 (versus 3 for the Gaussian), and extreme-event frequencies follow approximate power laws:
\begin{align}
    \mathbb{P}(|X_t| > x) \sim x^{-\alpha}, \quad \alpha \in (3,5),
\end{align}
for large $x$. This means that a daily return of $-5\%$ or worse occurs far more often than a Gaussian model predicts.
\end{remark}

\begin{remark}[Leverage effect]
Negative returns increase future volatility more than positive returns of equal magnitude. This asymmetry, called the \textbf{leverage effect} (Black, 1976), arises partly because a price decline increases a firm's debt-to-equity ratio and hence its riskiness, and partly from behavioral feedback (panic amplifies selling). Empirically, the correlation between $X_t$ and $X_{t+1}^2$ is significantly negative for equity indices.
\end{remark}

\begin{remark}[Long memory in volatility]
Volatility proxies such as $X_t^2$ and $|X_t|$ often exhibit slowly decaying autocorrelations:
\begin{align}
    \text{Corr}(X_t^2, X_{t+k}^2) \sim k^{-\delta}, \quad 0 < \delta < 1,
\end{align}
suggesting \textbf{long-range dependence} rather than the exponential decay implied by short-memory models like GARCH(1,1). This has been quantified via the \textbf{Hurst exponent} $H$, defined through the scaling behavior
\begin{align}
    \mathbb{E}[|X_{t+h} - X_t|^2] \sim h^{2H}.
\end{align}
Brownian motion has $H = 1/2$; volatility estimates for equity indices typically yield $H \approx 0.6$--$0.8$, indicating persistent positive dependence. This observation has motivated the rough volatility literature, where the Hurst exponent of the volatility process itself is $H \approx 0.1$ (\emph{rougher} than Brownian motion) which produces realistic term structures of implied volatility.
\end{remark}


\subsection{ARCH and GARCH}

\begin{definition}
An \textbf{ARCH($q$)} (autoregressive conditional heteroskedasticity) model specifies the return $X_t$ as
\begin{align}
    X_t = \sigma_t z_t, \qquad \sigma_t^2 = \alpha_0 + \sum_{i=1}^q \alpha_i X_{t-i}^2,
\end{align}
where $\{z_t\}$ is i.i.d.\ with $\mathbb{E}[z_t] = 0$, $\mathbb{E}[z_t^2] = 1$, and $\alpha_0 > 0$, $\alpha_i \geq 0$.
\end{definition}

\begin{remark}
ARCH captures volatility clustering through a direct mechanism: the conditional variance $\sigma_t^2$ is large when recent squared returns $X_{t-1}^2, \ldots, X_{t-q}^2$ are large. The unconditional distribution of $X_t$ has heavier tails than $z_t$ (even if $z_t$ is Gaussian), because the variance itself fluctuates. However, ARCH models require many lags $q$ to capture the persistence observed in real data, motivating a more parsimonious extension.
\end{remark}

\begin{definition}
A \textbf{GARCH($p,q$)} (generalized ARCH) model incorporates lagged conditional variances:
\begin{align}
    \sigma_t^2 = \omega + \sum_{i=1}^q \alpha_i X_{t-i}^2 + \sum_{j=1}^p \beta_j \sigma_{t-j}^2,
\end{align}
with $\omega > 0$, $\alpha_i, \beta_j \geq 0$.
\end{definition}

\begin{remark}
The inclusion of lagged $\sigma_{t-j}^2$ terms allows GARCH to capture long-lived volatility persistence with few parameters. A GARCH(1,1) with two parameters ($\alpha_1$ and $\beta_1$) can approximate an ARCH process of arbitrarily high order. This is analogous to how an ARMA(1,1) can approximate a high-order AR: the ``G'' in GARCH plays the same role as the MA term in ARMA.
\end{remark}

\begin{proposition}
For a GARCH(1,1) model, a sufficient condition for covariance stationarity is
\begin{align}
    \alpha_1 + \beta_1 < 1,
\end{align}
in which case the unconditional variance exists and equals
\begin{align}
    \mathbb{E}[\sigma_t^2] = \frac{\omega}{1 - \alpha_1 - \beta_1}.
\end{align}
\end{proposition}

\begin{proof}[Proof sketch]
Taking unconditional expectations of both sides: $\mathbb{E}[\sigma_t^2] = \omega + \alpha_1 \mathbb{E}[X_{t-1}^2] + \beta_1 \mathbb{E}[\sigma_{t-1}^2]$. Since $\mathbb{E}[X_{t-1}^2] = \mathbb{E}[\sigma_{t-1}^2 z_{t-1}^2] = \mathbb{E}[\sigma_{t-1}^2]$ (by independence and $\mathbb{E}[z^2]=1$), stationarity gives $\bar{\sigma}^2 = \omega + (\alpha_1 + \beta_1)\bar{\sigma}^2$, which solves to the stated formula when $\alpha_1 + \beta_1 < 1$.
\end{proof}

\begin{remark}
Empirical GARCH(1,1) fits to equity index returns typically yield $\alpha_1 \approx 0.05$--$0.10$ and $\beta_1 \approx 0.85$--$0.93$, so that $\alpha_1 + \beta_1 \approx 0.95$--$0.99$. This is close to 1, implying highly persistent volatility: after a shock, the conditional variance decays slowly toward the unconditional mean. The half-life of a variance shock is $\log(0.5)/\log(\alpha_1 + \beta_1)$; for $\alpha_1 + \beta_1 = 0.97$, this is about 23 days.
\end{remark}

\begin{example}[GARCH(1,1) for the S\&P 500]
Consider a GARCH(1,1) calibrated to daily S\&P 500 returns:
\begin{align}
    \sigma_t^2 = 2.0 \times 10^{-6} + 0.08\, X_{t-1}^2 + 0.91\, \sigma_{t-1}^2.
\end{align}
The persistence is $\alpha_1 + \beta_1 = 0.99$, and the unconditional variance is
\begin{align}
    \bar{\sigma}^2 = \frac{2.0 \times 10^{-6}}{0.01} = 2.0 \times 10^{-4},
\end{align}
corresponding to an annualized volatility of $\sqrt{252 \times 2.0 \times 10^{-4}} \approx 22.4\%$.

Suppose yesterday's return was $X_{t-1} = -3\%$ and yesterday's conditional variance was $\sigma_{t-1}^2 = 2.5 \times 10^{-4}$ (annualized $\approx 25\%$). Today's conditional variance is
\begin{align}
    \sigma_t^2 = 2.0 \times 10^{-6} + 0.08 \times (0.03)^2 + 0.91 \times 2.5 \times 10^{-4} = 3.02 \times 10^{-4},
\end{align}
corresponding to an annualized volatility of $\sqrt{252 \times 3.02 \times 10^{-4}} \approx 27.6\%$. The large negative return has pushed the one-day-ahead volatility forecast from 25\% to 27.6\%.

In contrast, if yesterday's return had been $X_{t-1} = +3\%$, the forecast would be \emph{identical}; the GARCH(1,1) treats positive and negative shocks symmetrically. This is a known limitation.
\end{example}


\subsection{Asymmetric and Long-Memory Extensions}

Standard GARCH treats positive and negative shocks identically, contradicting the leverage effect. Several extensions address this.

\begin{definition}
A \textbf{GJR-GARCH} model (Glosten, Jagannathan, Runkle, 1993) introduces asymmetry via an indicator function:
\begin{align}
    \sigma_t^2 = \omega + \alpha\, X_{t-1}^2 + \gamma\, X_{t-1}^2\, \mathbb{I}_{\{X_{t-1} < 0\}} + \beta\, \sigma_{t-1}^2.
\end{align}
When $\gamma > 0$, negative returns have an amplified impact on future volatility: the effective $\alpha$ for down moves is $\alpha + \gamma$, while for up moves it remains $\alpha$.
\end{definition}

\begin{definition}
An \textbf{EGARCH} (exponential GARCH, Nelson, 1991) model specifies the \emph{log} of the conditional variance:
\begin{align}
    \log \sigma_t^2 = \omega + \sum_{j=1}^p \beta_j \log \sigma_{t-j}^2 + \sum_{i=1}^q \bigl[\alpha_i z_{t-i} + \gamma_i\bigl(|z_{t-i}| - \mathbb{E}|z_{t-i}|\bigr)\bigr].
\end{align}
The $\alpha_i$ terms introduce asymmetry (negative $z_{t-i}$ increases log-variance), and modeling the logarithm ensures $\sigma_t^2 > 0$ without parameter sign constraints.
\end{definition}

\begin{remark}
The practical difference between GJR-GARCH and EGARCH is minor for most applications. GJR-GARCH is easier to estimate (the conditional variance is linear in the parameters) and easier to interpret (the leverage coefficient $\gamma$ has a direct meaning). EGARCH guarantees positivity of the variance without constraints, which can improve numerical stability. Both are standard in risk management software.
\end{remark}

\begin{example}[Quantifying the leverage effect with GJR-GARCH]
Using the same S\&P 500 data, a fitted GJR-GARCH(1,1) gives:
\begin{align}
    \sigma_t^2 = 2.0 \times 10^{-6} + 0.02\, X_{t-1}^2 + 0.12\, X_{t-1}^2\, \mathbb{I}_{\{X_{t-1} < 0\}} + 0.90\, \sigma_{t-1}^2.
\end{align}
For a positive return of $+3\%$: the effective $\alpha$ is $0.02$, contributing $0.02 \times (0.03)^2 = 1.8 \times 10^{-5}$.

For a negative return of $-3\%$: the effective $\alpha$ is $0.02 + 0.12 = 0.14$, contributing $0.14 \times (0.03)^2 = 1.26 \times 10^{-4}$ (seven times larger).

This asymmetry explains why markets ``take the stairs up and the elevator down'': negative shocks feed into volatility far more aggressively than positive ones.
\end{example}

While GARCH-type models generate persistent volatility, their autocorrelation structure is ultimately exponential in decay. To capture the empirically observed long memory, fractional models are needed.

\begin{definition}
A \textbf{FIGARCH($1,d,1$)} (fractionally integrated GARCH) model replaces the integer differencing in the GARCH recursion with a fractional difference operator:
\begin{align}
    \sigma_t^2 = \omega + \bigl[1 - \beta(L)\bigr]^{-1}\bigl[1 - \Delta^d\bigr]\alpha(L)\, X_t^2,
\end{align}
where $d \in (0,1)$ governs the degree of long memory.
\end{definition}

\begin{remark}
FIGARCH generates hyperbolically decaying autocorrelations in squared returns: $\text{Corr}(X_t^2, X_{t+k}^2) \sim k^{-d}$. Standard GARCH(1,1) is nested as $d = 0$. The cost is increased complexity and weaker theoretical guarantees (strict stationarity may fail for certain parameter configurations). Typical estimates for equity data give $d \approx 0.3$--$0.5$.
\end{remark}


\subsection{Forecasting versus Pricing}

\begin{remark}
The ARCH/GARCH family is designed for \emph{forecasting} conditional variance under the physical measure $\mathbb{P}$. It takes observed returns as input and produces one-step-ahead (or multi-step) variance forecasts, which feed directly into risk measures like VaR and ES.

However, these models do not specify dynamics under a risk-neutral measure $\mathbb{Q}$, nor do they embed the asset in a continuous-time, arbitrage-free framework. Consequently, they are not directly usable for derivative pricing. Stochastic volatility models (Heston, SABR) and local volatility models (Dupire) serve that role: they are calibrated to market option prices and produce $\mathbb{Q}$-dynamics.

Bridging the two worlds (using $\mathbb{P}$-estimated volatility dynamics to inform $\mathbb{Q}$-calibrated pricing models) remains a central challenge. One approach is the \textbf{variance risk premium}: the persistent wedge between GARCH-forecasted variance (under $\mathbb{P}$) and implied variance from options (under $\mathbb{Q}$) is itself a tradable risk premium, which informs strategies like variance swaps and volatility-targeting portfolios.
\end{remark}



\section{Spectral Analysis and Cointegration}

The time-domain tools of the preceding sections, like ACFs, ARMA models, and GARCH, describe dependence through lags. An alternative perspective decomposes variability across \emph{frequencies}. This is the spectral viewpoint, and it reveals structure that can be invisible in the time domain. We then turn to cointegration, the idea that nonstationary series can share a common stochastic trend, which is a low-frequency phenomenon naturally expressed in spectral language.


\subsection{The Spectral Density}

\begin{definition}
Let $\{X_t\}_{t \in \mathbb{Z}}$ be a weakly stationary process with absolutely summable autocovariance: $\sum_{k=-\infty}^{\infty} |\Gamma(k)| < \infty$. The \textbf{spectral density} is the Fourier transform of the autocovariance:
\begin{align}
    f(\lambda) = \frac{1}{2\pi} \sum_{k=-\infty}^{\infty} \Gamma(k)\, e^{-ik\lambda}, \qquad \lambda \in [-\pi, \pi].
\end{align}
\end{definition}

\begin{remark}
The spectral density decomposes the total variance of $X_t$ across frequencies. Low frequencies ($\lambda$ near $0$) correspond to slow, trend-like fluctuations; high frequencies ($\lambda$ near $\pm\pi$) correspond to rapid oscillations. A peak in $f(\lambda)$ at some frequency $\lambda_0$ indicates a periodic or quasi-periodic component with period $2\pi/\lambda_0$. A flat spectrum corresponds to white noise with equal power at all frequencies.
\end{remark}

The autocovariance and spectral density are equivalent descriptions of second-order structure, related by Fourier inversion.

\begin{proposition}
The autocovariance is recovered from the spectral density by
\begin{align}
    \Gamma(k) = \int_{-\pi}^{\pi} e^{ik\lambda} f(\lambda)\, d\lambda.
\end{align}
In particular, the variance is $\Gamma(0) = \int_{-\pi}^{\pi} f(\lambda)\, d\lambda$ (the total ``power'' across all frequencies).
\end{proposition}

For ARMA models, the spectral density has a revealing closed form.

\begin{proposition}
Let $\{X_t\}$ follow an ARMA($p,q$) process $\phi(L)X_t = \theta(L)\varepsilon_t$. The spectral density is
\begin{align}
    f(\lambda) = \frac{\sigma^2}{2\pi} \cdot \frac{|\theta(e^{-i\lambda})|^2}{|\phi(e^{-i\lambda})|^2}.
\end{align}
\end{proposition}

\begin{remark}
This formula makes several phenomena visually transparent:

\emph{Unit roots.} If $\phi(z) = 0$ has a root near $z = 1$, then $|\phi(e^{-i\lambda})|$ is small near $\lambda = 0$, and the spectral density has a large peak at zero frequency. An exact unit root ($z = 1$) makes $\phi(1) = 0$ and the spectral density diverges to $+\infty$ at $\lambda = 0$---infinite power at zero frequency is the spectral signature of non-stationarity.

\emph{Mean reversion.} A strongly mean-reverting process (large negative MA coefficient, or AR root well inside the unit circle) has a spectral peak at high frequencies---it oscillates more rapidly than white noise.

\emph{Seasonality.} A seasonal AR component with period $s$ introduces spectral peaks at $\lambda = 2\pi k/s$ for integer $k$, visible as sharp spikes in the periodogram.
\end{remark}

\begin{example}[Spectral density of AR(1)]
For $X_t = \phi X_{t-1} + \varepsilon_t$ with $|\phi| < 1$:
\begin{align}
    f(\lambda) = \frac{\sigma^2}{2\pi} \cdot \frac{1}{|1 - \phi e^{-i\lambda}|^2} = \frac{\sigma^2}{2\pi} \cdot \frac{1}{1 - 2\phi\cos\lambda + \phi^2}.
\end{align}
When $\phi > 0$ (positive autocorrelation), the spectral density peaks at $\lambda = 0$ and is monotonically decreasing. The process is thus dominated by low-frequency variation, consistent with slowly decaying positive autocorrelations.

When $\phi < 0$ (negative autocorrelation), the peak shifts to $\lambda = \pi$. The process thus oscillates rapidly, alternating between positive and negative deviations.

As $\phi \uparrow 1$, the peak at zero grows without bound: the spectral density of a random walk diverges at zero frequency, reflecting the accumulation of variance over time.
\end{example}


\subsection{Cointegration}

Many financial and economic time series are individually non-stationary, since stock prices, exchange rates, and interest rate levels all contain stochastic trends, yet pairs or groups of these series often move together over the long run. Cointegration formalizes this observation.

\begin{definition}
Let $X_t \in \mathbb{R}^k$ be a vector of $I(1)$ processes. The components are \textbf{cointegrated} if there exists a nonzero vector $\beta \in \mathbb{R}^k$ such that the linear combination
\begin{align}
    \beta^\top X_t \sim I(0)
\end{align}
is stationary. The vector $\beta$ defines a \textbf{cointegrating relationship}, where there exists a long-run equilibrium among the components.
\end{definition}

\begin{remark}
Cointegration is a statement about low-frequency behavior: while each component of $X_t$ has a spectral density that diverges at $\lambda = 0$ (the signature of $I(1)$), the combination $\beta^\top X_t$ has finite spectral mass at zero. The stochastic trends cancel in the linear combination, leaving only stationary fluctuations.
\end{remark}

\begin{example}[Pairs trading motivation]
Consider two stocks in the same sector with log-prices $P_t^{(1)}$ and $P_t^{(2)}$, both $I(1)$. If the companies are exposed to the same fundamental risk factors, their prices may be cointegrated: there exists a $\beta$ such that $P_t^{(1)} - \beta P_t^{(2)}$ is stationary. This spread fluctuates around a constant mean and mean-reverts after deviations.

A \textbf{pairs trading strategy} buys the underperforming stock and sells the overperformer when the spread deviates significantly from its mean, betting on reversion. The strategy is profitable if and only if the cointegrating relationship is stable, which which is an empirical question, not a mathematical certainty. Cointegration can and does break down, particularly during structural shifts (mergers, regulatory changes, sector rotations).
\end{example}

The dynamic relationship between cointegrated variables is captured by the error-correction representation.

\begin{theorem}[Granger representation]
If $X_t$ is a cointegrated $I(1)$ vector process, it admits a \textbf{vector error-correction model (VECM)}:
\begin{align}
    \Delta X_t = \Pi X_{t-1} + \sum_{j=1}^{p-1} \Gamma_j \Delta X_{t-j} + \varepsilon_t,
\end{align}
where $\Pi = \alpha\beta^\top$ has reduced rank. The matrix $\beta$ contains the cointegrating vectors (the long-run equilibria), $\alpha$ contains the \textbf{adjustment coefficients} (the speed at which each variable corrects deviations from equilibrium), and the $\Gamma_j$ capture short-run dynamics.
\end{theorem}

\begin{remark}
The VECM has a clean interpretation. The term $\Pi X_{t-1} = \alpha(\beta^\top X_{t-1})$ is the \textbf{error-correction term}: when $\beta^\top X_{t-1}$ deviates from zero (the equilibrium is violated), the system adjusts at rates determined by $\alpha$. If $\alpha_i < 0$ for variable $i$, then a positive deviation in $\beta^\top X_{t-1}$ causes variable $i$ to decrease in the next period, pulling the system back toward equilibrium.

The rank of $\Pi$ equals the number of cointegrating relationships. In a $k$-variable system, there can be at most $k-1$ independent cointegrating vectors. If $\text{rank}(\Pi) = 0$, the variables are not cointegrated; if $\text{rank}(\Pi) = k$, the variables are all $I(0)$ and no differencing was needed.
\end{remark}

\begin{example}[VECM for spot and futures prices]
Let $S_t$ and $F_t$ be log-prices of a commodity spot and its near-month futures contract. Both are $I(1)$, but the spread $S_t - F_t$ (the basis) is stationary---these are cointegrated with $\beta = (1, -1)$.

A VECM takes the form:
\begin{align}
    \begin{pmatrix} \Delta S_t \\ \Delta F_t \end{pmatrix}
    =
    \begin{pmatrix} \alpha_S \\ \alpha_F \end{pmatrix}
    (S_{t-1} - F_{t-1})
    + \text{(lagged differences)}
    + \varepsilon_t.
\end{align}
If $\alpha_S < 0$ and $\alpha_F > 0$, then when the basis widens (spot above futures), the spot price decreases and the futures price increases, closing the gap. The magnitudes of $\alpha_S$ and $\alpha_F$ determine which price adjusts faster.

Empirically, for liquid commodities, $|\alpha_F| > |\alpha_S|$: the futures price does most of the adjusting, reflecting its role as the ``price discovery'' vehicle.
\end{example}

\begin{remark}[Testing for cointegration]
Two standard approaches are widely used in practice.

The \textbf{Engle-Granger two-step procedure}: (1) regress $X_t^{(1)}$ on $X_t^{(2)}$ (and possibly other components) by OLS to obtain residuals $\hat{u}_t$; (2) test $\hat{u}_t$ for stationarity using an ADF test. If the residuals are stationary, the variables are cointegrated. This is simple and intuitive but handles only one cointegrating vector and suffers from small-sample bias.

The \textbf{Johansen procedure}: estimate the VECM directly and use likelihood ratio tests (the trace and maximum eigenvalue statistics) to determine the rank of $\Pi$. This handles multiple cointegrating relationships simultaneously and provides confidence intervals for $\beta$ and $\alpha$.
\end{remark}

\begin{remark}
Cointegration is a powerful tool, but three cautions are in order. First, cointegration is a long-run, low-frequency property. Two series can be cointegrated over a ten-year window but diverge over any given month; strategies built on cointegration require patience and risk tolerance. Second, cointegration is estimated, not known. Estimated cointegrating vectors change with the sample, and out-of-sample stability is never guaranteed. Third, embedding cointegration-based strategies in no-arbitrage pricing frameworks requires care, because the mean-reverting spread is not itself a traded asset with a well-defined martingale measure.
\end{remark}



\chapter{Numerical Methods}

The continuous-time models of the preceding chapters yield elegant partial differential equations and exact martingale representations. However, aside from a few idealized cases like the Black-Scholes or Vasicek models, these equations rarely admit closed-form solutions. Numerical methods bridge the gap between theoretical models and practical implementation. 

This chapter develops the two primary computational engines of quantitative finance: Monte Carlo simulation (which evaluates expectations by averaging over random paths) and finite difference methods (which solve PDEs by discretizing space and time).

\section{Monte Carlo Simulation}

Monte Carlo methods estimate the price of a derivative by simulating thousands of independent sample paths of the underlying stochastic process, evaluating the payoff on each path, and taking the discounted average.

\begin{definition}[Monte Carlo Estimator]
Let $X_T$ be the state of the market at maturity $T$, and let $\Phi(X_T)$ be the payoff of a derivative. Under the risk-neutral measure $\mathbb{Q}$, the exact time-$0$ value is $V_0 = e^{-rT}\mathbb{E}^{\mathbb{Q}}[\Phi(X_T)]$. By generating $M$ independent paths $X_T^{(1)}, \dots, X_T^{(M)}$, the \textbf{Monte Carlo estimator} is:
\begin{align}
    V_0^{(M)} = e^{-rT}\frac{1}{M}\sum_{m=1}^M \Phi(X_T^{(m)})
\end{align}
\end{definition}

\begin{theorem}[Monte Carlo Convergence]
By the Strong Law of Large Numbers, if the expected squared payoff is finite ($\mathbb{E}^{\mathbb{Q}}[\Phi(X_T)^2] < \infty$), the estimator converges almost surely to the true price: $V_0^{(M)} \xrightarrow{a.s.} V_0$. Furthermore, by the Central Limit Theorem:
\begin{align}
    \sqrt{M}\big(V_0^{(M)} - V_0\big) \xrightarrow{d} \mathcal{N}(0,\sigma^2)
\end{align}
where $\sigma^2 = \mathrm{Var}^{\mathbb{Q}}(\Phi(X_T))$.
\end{theorem}

\begin{remark}
The most profound implication of this theorem is that the root-mean-square error of the Monte Carlo estimator decays at a rate of $\mathcal{O}(M^{-1/2})$, \emph{independent of the dimension of the problem}. Finite difference methods suffer from the ``curse of dimensionality'' (their computational cost grows exponentially with the number of underlying assets). For a basket option on 50 stocks, Monte Carlo is the only viable pricing method.
\end{remark}

Because continuous-time SDEs like $dX_t = \mu(t,X_t)\,dt + \sigma(t,X_t)\,dW_t$ rarely have exact pathwise solutions, we must discretize time. 

\begin{definition}[Discretization Schemes]
Let $\Delta t = T/N$ and $t_n = n\Delta t$. The \textbf{Euler-Maruyama approximation} steps forward via:
\begin{align}
    X_{n+1} = X_n + \mu(t_n,X_n)\Delta t + \sigma(t_n,X_n)\Delta W_n
\end{align}
where $\Delta W_n \sim \mathcal{N}(0,\Delta t)$ are independent Gaussian increments. The Euler-Maruyama scheme has a \emph{strong order} of convergence (path-by-path error) of $1/2$ and a \emph{weak order} (error in expectation) of $1$. 

To achieve higher strong convergence (order $1$), one must use the \textbf{Milstein scheme}, which adds a second-order It\^o correction term: $\frac{1}{2}\sigma(t_n, X_n) \frac{\partial \sigma}{\partial x}(t_n, X_n) (\Delta W_n^2 - \Delta t)$.
\end{definition}

Because Monte Carlo convergence is slow ($\mathcal{O}(M^{-1/2})$), a massive number of paths is often required. 

\begin{definition}[Variance Reduction]
\textbf{Variance reduction techniques} aim to decrease the numerator $\sigma^2$ without increasing the computational burden $M$. 
\begin{itemize}
    \item \textbf{Antithetic Variates:} For every path generated with Brownian increments $\Delta W_n$, generate a shadow path using $-\Delta W_n$. If the payoff is monotonic, the negative correlation between the two paths drastically reduces the variance of their average.
    \item \textbf{Control Variates:} Let $Y$ be a highly correlated asset with a known analytical price $\mathbb{E}[Y]$. The modified estimator $\Phi^\ast = \Phi - \theta(Y - \mathbb{E}[Y])$ has lower variance if $\theta$ is chosen optimally as $\theta^\ast = \mathrm{Cov}(\Phi,Y)/\mathrm{Var}(Y)$. This is exceptionally powerful for pricing exotics by using their vanilla counterparts as controls.
\end{itemize}
\end{definition}

\begin{example}[Pricing American Options via Longstaff-Schwartz]
Historically, Monte Carlo struggled with American options because paths move forward, but optimal stopping requires knowing the future (working backward). The \textbf{Longstaff-Schwartz algorithm} elegantly solves this by approximating the continuation value at each time step using cross-sectional regression. 

At time $t_n$, the algorithm regresses the discounted future payoffs of in-the-money paths onto a set of basis functions (e.g., Laguerre polynomials) of the current state $X_n$. If the intrinsic value exceeds this predicted continuation value, the path exercises early. Though inherently slightly biased, it is the industry standard for multi-factor callable derivatives.
\end{example}

\section{Finite Difference Methods}

While Monte Carlo simulates the process forward in time, \textbf{finite difference methods} solve the backward Kolmogorov PDE derived from the Feynman-Kac theorem.

\begin{definition}[Grids and Operators]
Discretize the state space into a grid $S_i = i\Delta S$ and time into $t_n = n\Delta t$. We approximate the continuous derivatives using Taylor expansions:
\begin{itemize}
    \item $\frac{\partial V}{\partial S} \approx \frac{V_{i+1}^n - V_{i-1}^n}{2\Delta S}$ (Central difference)
    \item $\frac{\partial^2 V}{\partial S^2} \approx \frac{V_{i+1}^n - 2V_i^n + V_{i-1}^n}{(\Delta S)^2}$ (Second-order central difference)
\end{itemize}
Substituting these into the Black-Scholes PDE turns a continuous differential equation into a system of discrete algebraic equations for the grid values $V_i^n \approx V(t_n,S_i)$.
\end{definition}

\begin{definition}[Time-Stepping Schemes]
There are a number of methods to iterate through time steps:
\begin{itemize}
    \item \textbf{Explicit Scheme:} Computes $V_i^n$ (earlier time) directly from known values $V^{n+1}$ (later time). It is fast but \emph{conditionally stable}; if $\Delta t$ is too large relative to $(\Delta S)^2$, the numerical solution violently explodes.
    \item \textbf{Implicit Scheme:} Formulates the spatial derivatives at the unknown time step $n$. This requires inverting a tridiagonal matrix at each step, but it is \emph{unconditionally stable}.
    \item \textbf{Crank-Nicolson:} A 50/50 weighted average of the explicit and implicit schemes. It achieves second-order accuracy in time $\mathcal{O}(\Delta t^2)$ and is the preferred method for most 1D and 2D pricing engines.
\end{itemize}
\end{definition}

\begin{theorem}[Lax Equivalence]
A numerical scheme converges to the true solution of a linear PDE if and only if it is both \textbf{consistent} (the discrete equation approaches the continuous PDE as the grid shrinks) and \textbf{stable} (errors do not compound infinitely during iteration). 
\end{theorem}

\begin{example}[The Free Boundary of an American Put]
Finite differences shine when handling early exercise features. For an American put, the value satisfies a variational inequality: $V(t,S) \ge \max(K-S,0)$. In an implicit finite difference scheme, after solving the linear system for the theoretical continuation values $V_i^n$, we enforce the American constraint pointwise:
\begin{align}
    V_i^n \leftarrow \max(V_i^n,\; K-S_i)
\end{align}
The optimal exercise boundary $S^\ast(t)$ naturally emerges as the precise grid interface where the intrinsic value overtakes the continuation value.
\end{example}



\end{document}
